\chapter{Gromov--Witten 不变量}\label{chap:gromov-witten}
\SourcePage{190}
如\hyperref[sec:three-point]{第~\ref*{sec:three-point} 节}所见，Calabi--Yau 三维簇 $V$ 上的 A 模型关联函数用不变量 $n_\beta$ 表示。朴素地说，$n_\beta$ 是 $V$ 上代表同调类 $\beta$ 的有理曲线的数目，尤其参见式~\eqref{eq:a-yukawa}。本章有两个目标：严格定义 $n_\beta$，并严格定义相关的 \emph{Gromov--Witten 不变量}。事实上，$n_\beta$ 将通过 Gromov--Witten 不变量来定义。\index{Gromov--Witten不变量@Gromov--Witten 不变量|hyperpage}

在以下一般情形中开展讨论。给定射影代数簇 $X$，固定同调类 $\beta\in H_2(X,\mathbb Z)$ 及 $X$ 上的圈 $Z_1,\ldots,Z_n$。基本问题是研究下列曲线集合的结构：
\begin{equation}\label{eq:gw-naive-curves}
C\subset X,\qquad C\text{ 的亏格为 }g,\quad[C]=\beta,\quad
C\cap Z_i\ne\varnothing\text{ 对所有 }i\text{ 成立},
\end{equation}
其中假定各 $Z_i$ 处于一般位置。我们将区分 \emph{Gromov--Witten 不变量}与 \emph{Gromov--Witten 类}：粗略地说，前者是式~\eqref{eq:gw-naive-curves} 中曲线只有有限条时的曲线数目，后者则是数目无限时得到的上同调类。\index{Gromov--Witten类@Gromov--Witten 类|hyperpage}

Gromov--Witten 不变量源于物理学中的\emph{与引力耦合的拓扑 $\sigma$ 模型}。特别地，亏格 $0$ 的 Gromov--Witten 不变量，有时称为\emph{树级}不变量，来自拓扑 $\sigma$ 模型，而后者是一种拓扑量子场论。尤其，Gromov--Witten 不变量的构造本来就要求它在 $X$ 的复结构形变下不变。\hyperref[sec:topological-qft]{附录 B.6} 简要介绍拓扑量子场论。

下面将看到，式~\eqref{eq:gw-naive-curves} 对 Gromov--Witten 不变量或类的描述实在过于朴素。要给出严格定义，需要相当精细的工具。代数几何与辛几何各提供了若干方法。第一个具有相当一般性的定义是在半正辛流形的框架中给出的：亏格 $0$ 的情形见 \SourceCite{RT1}，此前一些特殊情形已见于 \SourceCite{Ruan2}；更高亏格的情形见 \SourceCite{RT2}。此后，又出现了对一般辛流形的若干推广，如 \SourceCite{LTi3,Siebert1,Ruan3,FO}。在代数几何方面，早期首先为 Grassmann 簇构造了这些不变量 \SourceCite{BDW}，然后推广到齐性空间 \SourceCite{LTi1,FP}。一般的代数几何定义见 \SourceCite{LTi2} 以及 \SourceCite{BF,Behrend}。

\hyperref[sec:gw-algebraic-definition]{第~\ref*{sec:gw-algebraic-definition} 节}和\hyperref[sec:gw-symplectic-definition]{第~\ref*{sec:gw-symplectic-definition} 节}将介绍定义 Gromov--Witten 类的主要代数几何方法与辛几何方法。定义种类繁多，可能令人望而生畏；幸运的是，它们都满足\hyperref[subsec:gw-axioms]{第~\ref*{subsec:gw-axioms} 小节}陈述的 Gromov--Witten 类的基本公理。通常相信，各定义在共同的适用范围内彼此一致。\hyperref[subsec:gw-equivalence]{第~\ref*{subsec:gw-equivalence} 小节}将说明，某些情形中已经证明了这一点，但一般情形的等价性尚不清楚。因此，将所有这些对象都称作
\SourcePage{191}
“Gromov--Witten 不变量”，在某种意义上是对记号的宽松使用。幸运的是，只要所做的工作仅涉及 Gromov--Witten 公理，就没有问题；而在讨论各定义猜想中的等价性可能带来的推论时，我们会明确指出所使用的定义。

本章末尾将给出几个计算 Gromov--Witten 不变量的例子。接着，\hyperref[chap:quantum-cohomology]{第~\ref*{chap:quantum-cohomology} 章}将用这些不变量构造\emph{量子上同调环}。

文献中已有一些很好的 Gromov--Witten 理论入门资料。特别推荐 Voisin 的书 \SourceCite{Voisin3}、Morrison 的 Park City 讲义 \SourceCite{Morrison7}，以及 Fulton 和 Pandharipande 的 Santa Cruz 讲义 \SourceCite{FP}。

\section{代数几何中的定义}\label{sec:gw-algebraic-definition}
\index{Gromov--Witten类@Gromov--Witten 类!代数几何中的定义|hyperpage}
在 Gromov--Witten 理论的代数几何方法中，Kontsevich \SourceCite{Kontsevich2} 的关键观察是：式~\eqref{eq:gw-naive-curves} 中的曲线 $C\subset X$ 应当换成带 $n$ 个标记的曲线 $(C,p_1,\ldots,p_n)$ 与全纯映射 $f:C\to X$。这一观察与该概念在非线性 $\sigma$ 模型（\hyperref[sec:nonlinear-sigma-models]{附录 B.2}）以及与引力耦合的二维拓扑量子场论中的起源一致。于是，式~\eqref{eq:gw-naive-curves} 提示应研究映射
\begin{equation}\label{eq:gw-pointed-maps}
f:C\longrightarrow X,\qquad f_*[C]=\beta,\quad
f(p_i)\in Z_i,\quad i=1,\ldots,n.
\end{equation}
这些映射之间的同构概念是显然的，因而需要考虑它们的\emph{模}。为了得到紧的模空间，将允许某些亏格 $g$ 的可约曲线 $C$。带 $n$ 个标记的亏格 $g$ 稳定曲线的 Deligne--Mumford 粗模空间 $\overline M_{g,n}$，将在理论中发挥重要作用。\SourceCite{FP} 给出了关于 $\overline M_{g,n}$ 的很好介绍。

Gromov--Witten 类的粗略思想是：式~\eqref{eq:gw-pointed-maps} 中的映射给出 $\overline M_{g,n}$ 的一个子集，进而给出 $H^*(\overline M_{g,n},\mathbb Q)$ 中的上同调类。采用 $\mathbb Q$ 系数，是因为当 $n+2g\geq3$ 时，$\overline M_{g,n}$ 存在，并且是维数 $3g-3+n$ 的轨形。为了使这一表述更正式，令 $\alpha_i\in H^*(X)$ 为圈 $Z_i$ 的对偶上同调类。则 Gromov--Witten 类
\[
I_{g,n,\beta}(\alpha_1,\ldots,\alpha_n)\in H^*(\overline M_{g,n},\mathbb Q)
\]
直观上应是由式~\eqref{eq:gw-pointed-maps} 中出现的带标记曲线 $(C,p_1,\ldots,p_n)$ 的集合所代表的上同调类。因此，Gromov--Witten 类是一族映射
\begin{equation}\label{eq:gw-class-map}
I_{g,n,\beta}:H^*(X,\mathbb Q)^{\otimes n}
\longrightarrow H^*(\overline M_{g,n},\mathbb Q).
\end{equation}
\hyperref[sec:gw-properties]{第~\ref*{sec:gw-properties} 节}将研究这些映射的性质。

除了上同调类，还能定义数值不变量。具体如下：若 $I_{g,n,\beta}(\alpha_1,\ldots,\alpha_n)$ 如上所述，则得到有理数
\begin{equation}\label{eq:gw-numerical-invariant}
\langle I_{g,n,\beta}\rangle(\alpha_1,\ldots,\alpha_n)
=\int_{\overline M_{g,n}}I_{g,n,\beta}(\alpha_1,\ldots,\alpha_n).
\end{equation}
这就是一个 \emph{Gromov--Witten 不变量}。注意，除非 Gromov--Witten 类
$I_{g,n,\beta}(\alpha_1,\ldots,\alpha_n)$ 在 $H^*(\overline M_{g,n},\mathbb Q)$ 中含有最高次数的分量，式~\eqref{eq:gw-numerical-invariant} 为零。直观地，当式~\eqref{eq:gw-pointed-maps} 只含有限条曲线时，$I_{g,n,\beta}(\alpha_1,\ldots,\alpha_n)$ 应具有最高次数，此时曲线数目就等于
$\langle I_{g,n,\beta}\rangle(\alpha_1,\ldots,\alpha_n)$。不过，\hyperref[sec:computing-gw-i]{第~\ref*{sec:computing-gw-i} 节}将看到，Gromov--Witten 不变量可能是分数，甚至为负，因此它们的枚举意义并非总是直接明了。

\SourcePage{192}
到目前为止，所有讨论仍然极其朴素。要定义 Gromov--Witten 类，必须更仔细地研究映射~\eqref{eq:gw-pointed-maps} 的模。特别地，需要了解\emph{稳定映射}与\emph{虚拟基本类}。下面分别讨论这两个主题。

\subsection{稳定映射及其模}\label{subsec:stable-maps}
\index{稳定曲线|hyperpage}\index{带标记曲线|hyperpage}\index{稳定映射!模空间|hyperpage}
寻找 Gromov--Witten 类的代数定义时，\SourceCite{Kontsevich2} 的另一个关键思想，是构造映射~\eqref{eq:gw-pointed-maps} 的模空间的一个紧化。这通过\emph{稳定映射}实现，定义如下。\index{稳定映射|hyperpage}

\begin{definition}\label{def:stable-map}
带 $n$ 个标记的稳定映射，由连通的带标记曲线 $(C,p_1,\ldots,p_n)$ 与态射 $f:C\to X$ 组成，满足以下性质：
\begin{enumerate}[label=(\roman*)]
\item $C$ 的奇点只能是普通二重点。
\item $p_1,\ldots,p_n$ 是 $C$ 上互不相同、有次序的光滑点。
\item 若 $C_i$ 为 $C$ 的一个分支，$C_i\simeq\mathbb P^1$，且 $f$ 在 $C_i$ 上为常值，则 $C_i$ 至少含三个特殊点，即结点或标记点。
\item 若 $C$ 的算术亏格为 $1$ 且 $n=0$，则 $f$ 不是常值映射。
\end{enumerate}
\end{definition}
在定义的前两个条件成立的前提下，第三与第四个条件等价于：数据
$(f,C,p_1,\ldots,p_n)$ 只有有限多个自同构。无限多个自同构的存在，正是 $n+2g<3$ 时 $\overline M_{g,n}$ 不存在的原因。

如下定义稳定映射族。

\begin{definition}\label{def:stable-map-family}
设 $X$ 为射影代数簇，$S$ 为 $\mathbb C$ 上的概形。$S$ 上带 $n$ 个标记的稳定映射，是平坦的固有态射 $\mathcal C\to S$，连同 $n$ 个截面 $s_1,\ldots,s_n$ 及映射 $f:\mathcal C\to X$；对 $S$ 的每个几何点 $s$，$f$ 在 $\mathcal C$ 位于 $s$ 上的几何纤维上的限制 $f_s:\mathcal C_s\to X$，连同各截面 $s_i$ 的像，定义一个稳定映射。此外：
\begin{enumerate}[label=(\roman*)]
\item 若对 $S$ 的每个几何点 $s$，曲线 $\mathcal C_s$ 的算术亏格都为 $g$，则称 $f:\mathcal C\to X$ 的亏格为 $g$。
\item 给定同调类 $\beta\in H_2(X,\mathbb Z)$，若对 $S$ 的每个几何点 $s$，都有 $(f_s)_*[\mathcal C_s]=\beta$，则称 $f:\mathcal C\to X$ 的类为 $\beta$。
\end{enumerate}
\end{definition}
于是，给定 $\beta\in H_2(X,\mathbb Z)$，可定义反变函子
\[
\overline{\mathcal M}_{g,n}(X,\beta):
(\mathbb C\text{-概形})\longrightarrow(\text{集合}),
\]
将概形 $S$ 映为 $S$ 上所有亏格 $g$、类为 $\beta$、带 $n$ 个标记的稳定映射的同构类集合。

为了解决稳定映射的模问题，需要在某种意义下表示这一函子。最直接的方法，是考虑稳定映射 $f:C\to X$ 的\emph{粗模空间}。粗模空间的思想见 \SourceCite{Mumford1}。有以下基本存在性结果。\index{粗模空间|hyperpage}

\begin{theorem}\label{thm:stable-map-coarse-moduli}
若 $X$ 为射影簇，$\beta$ 为 $H_2(X,\mathbb Z)$ 中的同调类，则函子
$\overline{\mathcal M}_{g,n}(X,\beta)$ 有粗模空间 $\overline M_{g,n}(X,\beta)$，它是 $\mathbb C$ 上的射影概形。
\end{theorem}
这一结果首先由 Alexeev \SourceCite{Alexeev} 证明。\SourceCite{FP} 给出了 $\overline M_{g,n}(X,\beta)$ 的相当显式的构造。

模问题的一种更精细的方法涉及\emph{代数叠}。叠的语言见 \SourceCite{DM} 或 \SourceCite{BEFFGK}。
\SourcePage{193}
对本书的目的，只需指出，叠是从概形范畴到集合范畴、满足某些性质的反变函子；随着讨论展开，将明确说明其中一些性质。概形本身也是叠。对概形 $S$，可赋予由
\begin{equation}\label{eq:representable-scheme-functor}
h_S(X)=\operatorname{Hom}(X,S)
=\{\text{所有概形态射 }f:X\to S\}
\end{equation}
定义的可表示函子 $h_S$，其中 $X$ 为任意概形。\index{代数叠|hyperpage}

用这种语言，\SourceCite{Kontsevich2} 陈述了以下基本存在性结果。

\begin{theorem}\label{thm:stable-map-stack}
若 $X$ 为射影簇，$\beta$ 为 $H_2(X,\mathbb Z)$ 中的同调类，则
$\overline{\mathcal M}_{g,n}(X,\beta)$ 为 $\mathbb C$ 上的固有代数叠。此外，叠 $\overline{\mathcal M}_{0,n}(\mathbb P^r,\beta)$ 为光滑叠。
\end{theorem}
证明见 \SourceCite{BM}。由于这一结果，以下用 $\overline{\mathcal M}_{g,n}(X,\beta)$ 表示上述模叠。还要回忆，每个代数叠都有底层代数空间。对叠 $\overline{\mathcal M}_{g,n}(X,\beta)$，底层空间为粗模空间 $\overline M_{g,n}(X,\beta)$。有自然的叠态射
$\overline{\mathcal M}_{g,n}(X,\beta)\to\overline M_{g,n}(X,\beta)$，将 $S$ 上带 $n$ 个标记的稳定映射，映为其分类映射 $S\to\overline M_{g,n}(X,\beta)$。粗略地说，这一映射的存在表明，叠比粗模空间本身保留了更多信息。

叠是概形的推广，因此会经常用概形的语言讨论叠。我们希望读者能在几何上理解
$\overline{\mathcal M}_{g,n}(X,\beta)$，而不必深入涉及叠的技术问题；代价是引入一些不够精确的表述，希望不至于干扰阅读。想用叠的语言理解我们的确切意思的读者，需要稍多做一些工作。例如，有时把一个亏格 $g$、类为 $\beta$、带 $n$ 个标记的稳定映射
$f:(C,p_1,\ldots,p_n)\to X$ 称为 $\overline{\mathcal M}_{g,n}(X,\beta)$ 的一个\emph{点}。严格地说，函子并没有通常意义下的点，所以这一说法需要解释。按语境，它可能指关联于 $\overline{\mathcal M}_{g,n}(X,\beta)$ 的代数空间上的相应点，也可能指群胚
$\overline{\mathcal M}_{g,n}(X,\beta)(\operatorname{Spec}\mathbb C)$ 中的相应对象。

略说特殊情形 $g=0$、$X=\mathbb P^r$。此时可写 $\beta=d\ell$，其中 $\ell$ 为直线的类。由定理~\ref{thm:stable-map-stack}，
$\overline{\mathcal M}_{0,n}(\mathbb P^r,d\ell)$ 光滑，故其底层代数空间
$\overline M_{0,n}(\mathbb P^r,d\ell)$ 为轨形。可以证明，它是维数
$dr+r+d+n-3$ 的不可约正规射影簇。经常用 $\overline M_{g,n}(\mathbb P^r,d)$ 代替
$\overline M_{g,n}(\mathbb P^r,d\ell)$。

\subsection{从模空间到 Gromov--Witten 类}\label{subsec:moduli-to-gw}
\index{Gysin映射@Gysin 映射|hyperpage}\index{预期维数|hyperpage}
有了模空间 $\overline M_{g,n}(X,\beta)$，便得到以下自然映射：
\begin{equation}\label{eq:stable-map-evaluation-forgetting}
\begin{aligned}
\pi_1:\overline M_{g,n}(X,\beta)&\longrightarrow X^n,\\
\pi_2:\overline M_{g,n}(X,\beta)&\longrightarrow\overline M_{g,n}.
\end{aligned}
\end{equation}
这些映射容易定义。给定稳定映射 $f:(C,p_1,\ldots,p_n)\to X$，因为 $p_1,\ldots,p_n$ 是有次序的点列，得到良定义的 $n$ 元组
\[
(f(p_1),\ldots,f(p_n))\in X^n.
\]
这在集合层面给出了 $\pi_1$，不难验证它是态射。对于 $\pi_2$，先注意：若 $f:C\to X$ 为稳定映射，$C$ 并不必是 Deligne--Mumford 意义下的稳定曲线。
\SourcePage{194}
不过，若 $n+2g\geq3$，依次收缩 $C$ 的不稳定分支，就得到稳定曲线 $\widetilde C$。于是 $\pi_2$ 在集合层面将 $f:(C,p_1,\ldots,p_n)\to X$ 映为 $\widetilde C$ 的同构类。用 \SourceCite{Knudsen} 的技术可证明，它也是态射。

现在说明特殊情形下 Gromov--Witten 类的形式。假定 $X$ 光滑，且 $\overline M_{g,n}(X,\beta)$ 为维数
\begin{equation}\label{eq:stable-map-expected-dimension}
(1-g)(\dim X-3)-\int_\beta\omega_X+n
\end{equation}
的轨形，其中 $\omega_X$ 为 $X$ 的典范类。\hyperref[subsec:normal-cone]{第~\ref*{subsec:normal-cone} 小节}将说明，这是模空间的“预期维数”。在这些假设下，映射~\eqref{eq:stable-map-evaluation-forgetting} 给出自然映射
\[
\begin{aligned}
\pi_1^*:H^*(X,\mathbb Q)^{\otimes n}
&\longrightarrow H^*(\overline M_{g,n}(X,\beta),\mathbb Q),\\
\pi_{2*}:H_*(\overline M_{g,n}(X,\beta),\mathbb Q)
&\longrightarrow H_*(\overline M_{g,n},\mathbb Q),
\end{aligned}
\]
其中仍假定 $n+2g\geq3$。由于涉及的空间都是轨形，Poincaré 对偶（见\hyperref[app:singular-varieties]{附录 A}）与 $\pi_{2*}$ 诱导 Gysin 映射
\[
\pi_{2!}:H^*(\overline M_{g,n}(X,\beta),\mathbb Q)
\longrightarrow H^{2m+*}(\overline M_{g,n},\mathbb Q),
\]
其中 $m=(g-1)\dim X+\int_\beta\omega_X$。当 $n+2g\geq3$ 时，用公式
\begin{equation}\label{eq:gw-gysin-definition}
I_{g,n,\beta}(\alpha_1,\ldots,\alpha_n)
=\pi_{2!}\bigl(\pi_1^*(\alpha_1\otimes\cdots\otimes\alpha_n)\bigr)
\end{equation}
定义 Gromov--Witten 类。注意，这精确地定义了式~\eqref{eq:gw-pointed-maps} 中曲线集合的上同调类；回忆 $\alpha_i$ 是 $Z_i$ 的 Poincaré 对偶。容易验证，
$I_{g,n,\beta}(\alpha_1,\ldots,\alpha_n)$ 的次数为
\[
2m+\sum_{i=1}^n\deg\alpha_i
=2(g-1)\dim X+2\int_\beta\omega_X+\sum_{i=1}^n\deg\alpha_i.
\]
由式~\eqref{eq:gw-numerical-invariant}，还有 Gromov--Witten 不变量
\[
\langle I_{g,n,\beta}\rangle(\alpha_1,\ldots,\alpha_n)
=\int_{\overline M_{g,n}}I_{g,n,\beta}(\alpha_1,\ldots,\alpha_n),
\]
这里 $I_{g,n,\beta}(\alpha_1,\ldots,\alpha_n)$ 的维数须为 $2(3g-3+n)$，即 $\overline M_{g,n}$ 的实维数。利用式~\eqref{eq:gw-gysin-definition}，容易化简为
\begin{equation}\label{eq:gw-moduli-integral}
\langle I_{g,n,\beta}\rangle(\alpha_1,\ldots,\alpha_n)
=\int_{\overline M_{g,n}(X,\beta)}
\pi_1^*(\alpha_1\otimes\cdots\otimes\alpha_n).
\end{equation}
直观地，这恰好是式~\eqref{eq:gw-pointed-maps} 给出的数目。一个重要观察是：当 $n+2g<3$ 时，式~\eqref{eq:gw-moduli-integral} 仍有意义。因此，虽然 Gromov--Witten 类要求 $n+2g\geq3$，Gromov--Witten 不变量却对 $n,g\geq0$ 都有定义。\hyperref[chap:mirror-theorem]{第~\ref*{chap:mirror-theorem} 章}中，当语境已经明确 $n$ 时，经常采用简写
\[
\langle\alpha_1,\ldots,\alpha_n\rangle_{g,\beta}
=\langle I_{g,n,\beta}\rangle(\alpha_1,\ldots,\alpha_n).
\]

到目前为止，所有构造都假定 $\overline M_{g,n}(X,\beta)$ 为具有预期维数~\eqref{eq:stable-map-expected-dimension} 的光滑轨形。一个满足这些假设的良好例子是 $g=0$、$X=\mathbb P^r$。如上所述，
$\overline M_{0,n}(\mathbb P^r,d)$ 是维数 $rd+r+d+n-3$ 的轨形，而这正是本例的预期维数。完整细节见 \SourceCite{BM} 或 \SourceCite{FP}。

遗憾的是，定义~\eqref{eq:gw-gysin-definition} 和~\eqref{eq:gw-moduli-integral} 一般并不适用。例如，当 $g\geq1$ 时，甚至当 $g=0$ 而 $X\ne\mathbb P^r$ 的大多数流形情形中，空间 $\overline M_{g,n}(X,\beta)$
\SourcePage{195}
可能含有维数超过预期维数的分支。为了看清如何修正，需要强调\emph{基本类}的作用。上述情形中，有基本类
\[
\xi=[\overline M_{g,n}(X,\beta)]
\in H_*(\overline M_{g,n}(X,\beta),\mathbb Q),
\]
它在 Poincaré 对偶下对应于 $1\in H^*(\overline M_{g,n}(X,\beta),\mathbb Q)$。可用 $\xi$ 重写式~\eqref{eq:gw-gysin-definition}。首先，若
$e=\dim\overline M_{g,n}(X,\beta)$，则 Poincaré 对偶
$PD:H^*(\overline M_{g,n}(X,\beta),\mathbb Q)\to H_{2e-*}(\overline M_{g,n}(X,\beta),\mathbb Q)$ 为与 $\xi$ 的帽积。式~\eqref{eq:stable-map-evaluation-forgetting} 中的 $\pi_1,\pi_2$ 给出映射
\[
\pi:\overline M_{g,n}(X,\beta)\longrightarrow X^n\times\overline M_{g,n},
\]
满足 $\pi_i=p_i\circ\pi$，其中 $p_1,p_2$ 为自然投影。因此
\[
\begin{aligned}
\pi_{2!}\bigl(\pi_1^*(\alpha_1\otimes\cdots\otimes\alpha_n)\bigr)
&=\pi_{2!}PD^{-1}\bigl(\pi_1^*(\alpha_1\otimes\cdots\otimes\alpha_n)\cap\xi\bigr)\\
&=PD^{-1}\pi_{2*}\bigl(\pi_1^*(\alpha_1\otimes\cdots\otimes\alpha_n)\cap\xi\bigr)\\
&=PD^{-1}p_{2*}\pi_*\bigl(\pi^*p_1^*(\alpha_1\otimes\cdots\otimes\alpha_n)\cap\xi\bigr).
\end{aligned}
\]
最后两个等式中的 $PD$ 指 $\overline M_{g,n}$ 的 Poincaré 对偶。由投影公式及式~\eqref{eq:gw-gysin-definition}，Gromov--Witten 类为
\begin{equation}\label{eq:gw-fundamental-class-definition}
I_{g,n,\beta}(\alpha_1,\ldots,\alpha_n)
=PD^{-1}p_{2*}\bigl(p_1^*(\alpha_1\otimes\cdots\otimes\alpha_n)\cap\pi_*(\xi)\bigr),
\end{equation}
其中 $n+2g\geq3$。类似地，对 $n,g\geq0$，式~\eqref{eq:gw-moduli-integral} 中的 Gromov--Witten 不变量可用 $\xi$ 写为
\[
\langle I_{g,n,\beta}\rangle(\alpha_1,\ldots,\alpha_n)
=\int_\xi\pi_1^*(\alpha_1\otimes\cdots\otimes\alpha_n).
\]

后文中，用各分量表示式~\eqref{eq:stable-map-evaluation-forgetting} 中的 $\pi_1$ 会很方便。对 $1\leq i\leq n$，定义\emph{赋值映射}
\[
e_i:\overline{\mathcal M}_{g,n}(X,\beta)\longrightarrow X,
\]
将稳定映射 $f:(C,p_1,\ldots,p_n)\to X$ 映为 $f(p_i)$；等价地，它是 $\pi_1$ 与 $X^n$ 到第 $i$ 个因子的投影的复合。用此记号，Gromov--Witten 不变量为
\begin{equation}\label{eq:gw-evaluation-integral}
\langle I_{g,n,\beta}\rangle(\alpha_1,\ldots,\alpha_n)
=\int_\xi e_1^*(\alpha_1)\cup\cdots\cup e_n^*(\alpha_n).
\end{equation}
关键在于，只要有适当的“基本类”
$\xi\in H_*(\overline{\mathcal M}_{g,n}(X,\beta),\mathbb Q)$，公式~\eqref{eq:gw-fundamental-class-definition} 和~\eqref{eq:gw-evaluation-integral} 就有意义。具有所需性质的类 $\xi$ 通常称为\emph{虚拟基本类}。不过，很快将看到，其定义相当微妙。\index{虚拟基本类|hyperpage}

\subsection{法锥}\label{subsec:normal-cone}
\index{法锥!与Euler类的关系@与 Euler 类的关系|hyperpage}
先从\emph{法锥}说起，它是虚拟基本类的一个典型范例。设 $Y$ 为维数 $n$ 的光滑簇，$E$ 为其上的秩 $r$ 向量丛。给定截面 $s\in H^0(E)$，令 $Z=Z(s)\subset Y$ 为 $s$ 的零点概形。随着 $s$ 变化，$Z$ 的表现可能很糟糕，甚至维数也会改变。这与空间 $\overline{\mathcal M}_{g,n}(X,\beta)$ 所面临的问题类似。不过，本情形的处理方法是熟知的，见 \SourceCite{Fulton1}；更深入的处理见 \SourceCite{Fulton2}。
\SourcePage{196}
若 $\mathcal I$ 为 $Z$ 在 $Y$ 中的理想层，则 $Y$ 中 $Z$ 的\emph{法锥}定义为
\begin{equation}\label{eq:normal-cone-definition}
C_ZY=\operatorname{Spec}_Z\left(\bigoplus_{k=0}^\infty
\mathcal I^k/\mathcal I^{k+1}\right).
\end{equation}
法锥 $C_ZY$ 是 $Z$ 上纯维数为 $n$ 的仿射锥。与 $s$ 相乘给出的满射
$\mathcal O(E^*)\to\mathcal I$，诱导满射
\[
\bigoplus_k\operatorname{Sym}^k
\bigl(\mathcal O(E^*)/\mathcal I\mathcal O(E^*)\bigr)
\longrightarrow\bigoplus_k\mathcal I^k/\mathcal I^{k+1}.
\]
取 $\operatorname{Spec}$，得到 $C_ZY$ 到 $Z$ 上的丛 $E|_Z$ 的嵌入。法锥定义圈类
$[C_ZY]\in A_n(E|_Z)$。于是，类 $s^*[C_ZY]$ 位于 $A_{n-r}(Z)$。注意，若 $s$ 将 $Z$ 定义为局部完全交，$Z$ 的维数恰好为 $n-r$。因此，$n-r$ 为 $Z$ 的“预期维数”。关于 $s^*[C_ZY]$，还有以下熟知结果，可用 \SourceCite{Fulton2} 的结果与技术证明。\index{法锥|hyperpage}

\begin{lemma}\label{lem:normal-cone-euler-class}
若 $i:Z\hookrightarrow Y$ 为包含映射，则
$i_*(s^*[C_ZY])\in A_{n-r}(Y)$ 为 $E$ 的 Euler 类 $c_r(E)\cap[Y]$。
\end{lemma}
这一引理表明，$s^*[C_ZY]$ 是 Euler 类 $c_r(E)\cap[Y]$ 的精细化。随着截面 $s$ 变化，类 $s^*[C_ZY]$ 的行为也很好。还要指出，有时可把 $s^*[C_ZY]$ 算为关联的\emph{超额法丛}的 Euler 类。有时将用上同调而非 Chow 群开展工作；在这种情形中，秩 $r$ 丛 $E$ 的 Euler 类指最高 Chern 类 $c_r(E)\in H^{2r}(Y,\mathbb Z)$。这与拓扑学中的通行用法稍有不同，后者将 Euler 类写为同调类 $c_r(E)\cap[Y]$。

下面的例子说明法锥与稳定映射的关系。

\SourceExampleNumber{7.1.3}{1}
\begin{example}\label{ex:quintic-lines-normal-cone}
设 $V\subset\mathbb P^4$ 为光滑五次三维簇。由于一般 $V$ 恰好含 2875 条直线，Gromov--Witten 不变量的任何合理定义，都应在 $\ell$ 为 $V$ 中直线的类时满足
$\langle I_{0,0,\ell}\rangle=2875$。注意，$\langle I_{0,0,\ell}\rangle$ 是具有 $n=0$ 个变元的函数，也就是一个数。后面最终将证明这一点；目前先看法锥为何与本例相关。

首先注意，$\overline M_{0,0}(\mathbb P^4,\ell)$ 是 $\mathbb P^4$ 中直线的 Grassmann 簇 $G(2,5)$。因此，包含映射 $V\subset\mathbb P^4$ 诱导自然映射
$\overline M_{0,0}(V,\ell)\to\overline M_{0,0}(\mathbb P^4,\ell)=G(2,5)$，将稳定映射映为其在 $\mathbb P^4$ 中的像直线。还要注意 $\dim G(2,5)=6$。例题~\ref{ex:quintic-virtual-class} 将如下描述这一映射：若 $U$ 为 $G(2,5)$ 上的秩 $2$ 重言子丛，则直线 $\ell$ 处的纤维 $U_\ell$ 是 $\mathbb C^5$ 中射影化为 $\ell$ 的二维子空间。$V$ 的一个方程给出秩 $6$ 丛 $E=\operatorname{Sym}^5U^*$ 的截面 $s$，而 $\overline M_{0,0}(V,\ell)$ 正是 $s$ 的零点概形。

上述构造产生类 $s^*[C_ZY]\in A_0(\overline M_{0,0}(V,\ell))$，其中
$C_ZY$ 是 $\overline M_{0,0}(V,\ell)=Z(s)\subset Y=G(2,5)$ 的法锥。由下文的例题~\ref{ex:quintic-virtual-class}，这个零圈就是 $\overline M_{0,0}(V,\ell)$ 的虚拟基本类。这是合理的，因为 $V$ 的典范丛平凡，故由式~\eqref{eq:stable-map-expected-dimension}，
$\overline M_{0,0}(V,\ell)$ 的预期维数为 $0$，其虚拟基本类就应是零圈。

若 $V$ 为一般五次三维簇，可证明截面 $s$ 的零点轨迹为既约完全交，因而
$\overline M_{0,0}(V,\ell)$ 的实际维数确为 $0$。于是，一般情形中，
$\langle I_{0,0,\ell}\rangle$ 恰好是 $V$ 上直线的数目。
\SourcePage{197}
例题~\ref{ex:quintic-lines-localization} 将算出这一数目为 2875，与\hyperref[chap:quintic]{第~\ref*{chap:quintic} 章}一致。

另一方面，设 $V\subset\mathbb P^4$ 为由
\[
x_0^5+x_1^5+x_2^5+x_3^5+x_4^5=0
\]
定义的 Fermat 五次三维簇。这是含有无限多条直线的光滑 Calabi--Yau 三维簇。为看清这一点，令 $\zeta$ 为任意五次单位根。对任意 $i\ne j$，超平面
$x_i+\zeta x_j=0$ 与 $V$ 相交于一个以平面 Fermat 五次曲线为底的锥，给出 $V$ 中的一个单参数直线族。进一步可证明，按此方式得到的 50 个族包含 $V$ 上的全部直线，每条直线的法丛都是
$N=\mathcal O_{\mathbb P^1}(-3)\oplus\mathcal O_{\mathbb P^1}(1)$。于是，
$\overline M_{0,0}(V,\ell)$ 是一维的，有 50 个分支，每个分支都是平面五次曲线。此外，法丛满足 $h^0(N)=2$，由此可证明 $\overline M_{0,0}(V,\ell)$ 处处非既约。事实上，每个分支在一般点的重数为 $2$；而对应于同时属于不止一个锥的 375 条直线的 375 个特殊点，各点重数为 $5$。例如，由
$x_0+x_1=x_2+x_3=x_4=0$ 定义的直线便是这种情形。这些结果都在 \SourceCite{AKa} 中用形变理论证明。

虽然 $\overline M_{0,0}(V,\ell)$ 很复杂，其虚拟基本类仍是法锥构造给出的零圈
$s^*[C_ZY]\in A_0(\overline M_{0,0}(V,\ell))$，其中 $s$ 由 Fermat 五次方程诱导。
\SourceCite[例题 4.2]{ClK} 描述了法锥 $C=C_ZY$ 的结构，并给出更多细节。
在 $\overline M_{0,0}(V,\ell)$ 的 50 个分支各自上方，法锥 $C$ 都有一个纤维维数 $5$、重数 $2$ 的分支；在对应于同时属于不止一个锥的 375 条直线的各特殊点上方，则有一个纤维维数 $6$、重数 $5$ 的分支。因此，$C$ 的纯维数为 $6$，正如预期。再看虚拟基本类 $s^*[C]$，可证明：$C$ 中位于 $\overline M_{0,0}(V,\ell)$ 的一个分支上方的每个分支，对虚拟基本类贡献 $20$，这里需要计算平面五次曲线上\IndexLocation{indexloc:0280:197}超额法丛的 Euler 类；位于特殊点上方的每个分支则贡献 $5$。由式~\eqref{eq:gw-evaluation-integral}，Gromov--Witten 不变量
$\langle I_{0,0,\ell}\rangle$ 为这一零圈的次数，即
$20\times50+5\times375=2875$，正如预期。

可以验证，上述计算与 \SourceCite{AKa} 完全对应。事实上，我们的论述表明，这一对应来自动态相交理论，见 \SourceCite[第 11 章]{Fulton2}。虚拟基本类在某种意义上延伸了动态相交理论的适用范围。
\end{example}

\subsection{虚拟基本类}\label{subsec:virtual-fundamental-class}
\index{切—障碍复形|hyperpage}\index{Kuranishi映射@Kuranishi 映射|hyperpage}\index{障碍理论|hyperpage}\index{余切复形|hyperpage}\index{虚法锥|hyperpage}\index{预稳定曲线|hyperpage}\index{内蕴法层|hyperpage}
为了对任意 $g,n,X,\beta$ 定义 $\overline{\mathcal M}_{g,n}(X,\beta)$ 的虚拟基本类，需要比法锥构造更一般的工具。因此，不再把 $\overline{\mathcal M}_{g,n}(X,\beta)$ 实现为向量丛截面的零点概形，而是赋予它更多结构。在 \SourceCite{LTi2} 的方法中，这种结构称为\emph{完美切—障碍复形}；在 \SourceCite{BF} 的方法中，则称为\emph{完美障碍理论}。\index{完美障碍理论|hyperpage}

深入这些定义的细节前，先说明障碍为何与虚拟基本类有关。根本原因是：计算式~\eqref{eq:stable-map-expected-dimension}
给出的粗模空间 $\overline M_{g,n}(X,\beta)$ 的预期维数时，障碍自然出现。说明如下。

“预期维数”的思想常出现在代数几何中；也许在紧复流形
\SourcePage{198}
$M$ 的形变情形中最容易理解。无穷小形变由 $H^1(M,\Theta_M)$ 给出，障碍位于
$H^2(M,\Theta_M)$。Kuranishi 理论 \SourceCite{Kuranishi} 将 $M$ 的模空间局部描述为全纯障碍映射
\begin{equation}\label{eq:kuranishi-obstruction-map}
U\longrightarrow H^2(M,\Theta_M)
\end{equation}
的零点轨迹，其中 $U\subset H^1(M,\Theta_M)$ 为开集。因此，$M$ 的模空间的预期维数为
\[
\dim H^1(M,\Theta_M)-\dim H^2(M,\Theta_M).
\]

对 $\overline M_{g,n}(X,\beta)$，推理类似。令 $f:(C,p_1,\ldots,p_n)\to X$ 为稳定映射。叠
$\overline{\mathcal M}_{g,n}(X,\beta)$ 在 $f$ 处的切空间是超 $\operatorname{Ext}$ 群
\[
\operatorname{Ext}_C^1
\left(f^*\Omega_X^1\longrightarrow\Omega_C^1\left(\sum_{i=1}^np_i\right),
\mathcal O_C\right),
\]
而障碍位于
\[
\operatorname{Ext}_C^2
\left(f^*\Omega_X^1\longrightarrow\Omega_C^1\left(\sum_{i=1}^np_i\right),
\mathcal O_C\right).
\]
这里沿用 \SourceCite{LTi2}。\SourceCite{Kontsevich2} 则用相应的切形变复形，将切空间与障碍空间写成
\[
\mathbb H^1(C,\Theta'_C\longrightarrow f^*\Theta_X),
\qquad
\mathbb H^2(C,\Theta'_C\longrightarrow f^*\Theta_X),
\]
其中 $\Theta_X$ 为 $X$ 的切层，$\Theta'_C$ 为在 $p_1,\ldots,p_n$ 处为零的向量场所组成的层；在结点处应采用相应的导出形变复形。

将看到，式~\eqref{eq:kuranishi-obstruction-map} 在这一框架中有对应形式。因此，
$\overline M_{g,n}(X,\beta)$ 的预期维数为
\[
\begin{aligned}
&\dim\operatorname{Ext}_C^1
\left(f^*\Omega_X^1\to\Omega_C^1\left(\sum_{i=1}^np_i\right),\mathcal O_C\right)\\
&\qquad-\dim\operatorname{Ext}_C^2
\left(f^*\Omega_X^1\to\Omega_C^1\left(\sum_{i=1}^np_i\right),\mathcal O_C\right).
\end{aligned}
\]
考察 $\operatorname{Ext}$ 的长正合列，其开头为
\begin{equation}\label{eq:stable-map-ext-sequence}
\begin{aligned}
0&\longrightarrow
\operatorname{Ext}_C^0
\left(f^*\Omega_X^1\to\Omega_C^1\left(\sum_{i=1}^np_i\right),\mathcal O_C\right)\\
&\longrightarrow\operatorname{Ext}_C^0
\left(\Omega_C^1\left(\sum_{i=1}^np_i\right),\mathcal O_C\right)
\longrightarrow\operatorname{Ext}_C^0(f^*\Omega_X^1,\mathcal O_C)\\
&\longrightarrow\operatorname{Ext}_C^1
\left(f^*\Omega_X^1\to\Omega_C^1\left(\sum_{i=1}^np_i\right),\mathcal O_C\right)
\longrightarrow\cdots.
\end{aligned}
\end{equation}
注意，$\operatorname{Ext}_C^0(\Omega_C^1(\sum_{i=1}^np_i),\mathcal O_C)$ 为
$(C,p_1,\ldots,p_n)$ 的无穷小自同构空间；由于 $C$ 本身不必稳定，这些自同构可能存在。
此外，$\operatorname{Ext}_C^0(f^*\Omega_X^1,\mathcal O_C)$ 为固定 $C$ 时映射 $f$ 的一阶形变空间。
由 $f$ 的稳定性，可知连接这两项的映射是单射。因此，
$\operatorname{Ext}_C^0(f^*\Omega_X^1\to\Omega_C^1(\sum_{i=1}^np_i),\mathcal O_C)=0$。
于是，预期维数可改写为
\[
\chi(f^*\Theta_X)+
\dim\operatorname{Ext}_C^1\left(\Omega_C^1\left(\sum_{i=1}^np_i\right),\mathcal O_C\right)
-\dim\operatorname{Ext}_C^0\left(\Omega_C^1\left(\sum_{i=1}^np_i\right),\mathcal O_C\right),
\]
其中使用同构 $\operatorname{Ext}_C^i(f^*\Omega_X^1,\mathcal O_C)\simeq H^i(f^*\Theta_X)$。
Hirzebruch--\allowbreak Riemann--\allowbreak Roch 定理给出
\[
\chi(f^*\Theta_X)=-\int_\beta\omega_X+(1-g)\dim X.
\]
另外两个 $\operatorname{Ext}$ 项中，
\[
\operatorname{Ext}_C^1\left(\Omega_C^1\left(\sum_{i=1}^np_i\right),\mathcal O_C\right)
\]
为 $(C,p_1,\ldots,p_n)$ 的一阶形变空间，而
\[
\operatorname{Ext}_C^0\left(\Omega_C^1\left(\sum_{i=1}^np_i\right),\mathcal O_C\right)
\]
为其无穷小自同构空间。
回忆 $(C,p_1,\ldots,p_n)$ 本身不必稳定。因此，两者维数之差为带标记曲线模叠的维数；在稳定范围内，这就是 $\overline M_{g,n}$ 的轨形切空间的维数 $3g-3+n$。故预期维数为
\[
\begin{aligned}
-\int_\beta\omega_X+(1-g)\dim X+(3g-3+n)
&=(1-g)(\dim X-3)-\int_\beta\omega_X+n,
\end{aligned}
\]
与式~\eqref{eq:stable-map-expected-dimension} 一致。

\SourcePage{199}
现在可以精确定义虚拟基本类。如前所述，\SourceCite{LTi2} 与 \SourceCite{BF} 给出了两种方法。
预期二者在共同适用范围内定义相同的虚拟基本类。证明应当直接，但尚未见于文献。因此，有时会针对关于虚拟基本类的一个论断给出两个证明，各对应一种方法。此外，使用特定方法时，会明确说明采用哪一种。为了强调虚拟基本类的思想，所用记号不区分这些方法。

先介绍 Li 与 Tian 的方法，更多细节见 \SourceCite{LTi2}。令
\[
\mathcal F:(\mathbb C\text{-概形})\longrightarrow(\text{集合})
\]
为反变模函子，将概形 $S$ 映为 $S$ 上某种固定类型的几何对象的所有同构类的集合。
稳定映射的同构类函子 $\overline{\mathcal M}_{g,n}(X,\beta)$ 是其一例。

设 $S$ 为 $\mathbb C$ 上的仿射概形，$\mathcal N$ 为 $\mathcal O_S$ 模。用 $S_{\mathcal N}$ 表示 $S$ 由 $\mathcal N$ 给出的平凡扩张，即
$S_{\mathcal N}=\operatorname{Spec}(\Gamma(\mathcal O_S)\oplus\Gamma(\mathcal N))$，其中
$\Gamma(\mathcal O_S)\oplus\Gamma(\mathcal N)$ 是 $\Gamma(\mathcal O_S)$ 由 $\Gamma(\mathcal N)$ 给出的平凡环扩张。特别地，得到无穷小扩张 $S\hookrightarrow S_{\mathcal N}$。
$\mathcal F$ 的\emph{切函子} $T\mathcal F$ 是以下函子族。对如上所述的元素
$\alpha\in\mathcal F(S)$，有反变函子
\[
T\mathcal F(\alpha):(\mathcal O_S\text{-模})\longrightarrow(\text{集合}),
\]
将 $\mathcal N$ 映为 $\mathcal F(S_{\mathcal N})$ 中在自然“限制”映射
$\mathcal F(S_{\mathcal N})\to\mathcal F(S)$ 下限制为 $\alpha$ 的所有元素的集合。
更精确的描述见 \SourceCite{LTi2}。因此，切函子记录 $\alpha$ 到 $S$ 的各平凡无穷小扩张的所有“扩张”的数据。组成 $T\mathcal F$ 的各函子满足自然的基变换性质。

$\mathcal F$ 的\emph{切—障碍复形}由函子复形
\[
T^1\mathcal F\longrightarrow T^2\mathcal F
\]
组成，其中箭头为零自然变换，$T^1\mathcal F$ 为上述切函子，$T^2\mathcal F$ 给出
$\mathcal F$ 的适当障碍理论。这里不进一步展开障碍理论，只指出：对数据 $(\alpha,S,\mathcal N)$，有障碍类
\[
\operatorname{ob}(\alpha,S\hookrightarrow S',\mathcal N)
\in\Gamma\bigl(T^2\mathcal F(\alpha)\otimes_{\mathcal O_S}\mathcal N\bigr),
\]
其消失，是存在将 $\alpha$ 延拓的元素 $\widetilde\alpha\in\mathcal F(S')$ 的充要条件。不过，障碍理论还包含更多信息；更多细节见 \SourceCite{LTi2}。

若对上述每个 $(\alpha,S)$，至少局部于 $S$，存在局部自由 $\mathcal O_S$ 模层的两项复形
$\mathcal E^1\to\mathcal E^2$，使对任意 $\mathcal N$，
$T^i\mathcal F(\alpha)(\mathcal N)$ 为诱导复形
$\mathcal E^\bullet\otimes_{\mathcal O_S}\mathcal N$ 的第 $i$ 个层上同调，则称切—障碍复形为\emph{完美的}。

假定 $\mathcal F$ 可由概形 $Z$ 表示。按式~\eqref{eq:representable-scheme-functor} 的记号，即
$\mathcal F=h_Z$，其中 $1_Z\in h_Z(Z)=\operatorname{Hom}(Z,Z)$ 对应于泛对象
$\alpha\in\mathcal F(Z)$。为了易于说明，再假定 $\mathcal F$ 有完美切—障碍复形，且定义中的
$\mathcal E^i$ 都由整个 $Z$ 上的局部自由层诱导，而不仅仅局部存在。
\SourcePage{200}
仍用 $\mathcal E^i$ 表示这些层。下面的例题~\ref{ex:zero-scheme-obstruction-theory} 给出这种完美障碍理论的简单例子。
对 $i=1,2$，令 $E^i$ 为满足 $\mathcal O(E^i)=\mathcal E^i$ 的向量丛。
Li 与 Tian 用障碍理论构造一个远非唯一的 \emph{Kuranishi 映射}
\[
F:\widehat E^1\longrightarrow E^2,
\]
其中 $\widehat E^1$ 为 $E^1$ 沿零截面的形式完备化。这一 Kuranishi 映射 $F$ 是通常映射~\eqref{eq:kuranishi-obstruction-map} 的形式相对版本。
$F$ 的特征性质是：在每个 $z\in Z$ 上的纤维中，$F$ 的零点轨迹与 $Z$ 在 $z$ 处的完备化形式同构。
更精确地说，$F$ 本身的零点轨迹形式同构于 $Z\times Z$ 沿对角线的完备化 $\widehat Z$。
特别地，Kuranishi 映射给出 $\widehat Z$ 到 $\widehat E^1$ 的嵌入。

现在构造\emph{虚法锥}便很简单。考虑 $\widehat Z$ 在 $\widehat E^1$ 中的法锥，它是
$\widehat Z$ 上的锥。\emph{虚法锥} $C^{\mathcal E^\bullet}$ 是该锥在
$Z\subset\widehat Z$ 上的限制。注意，通过 $Z\subset\widehat Z$，$Z$ 作为 $E^1$ 的零截面嵌入。
虚法锥的维数等于 $E^1$ 的秩。再用 Kuranishi 映射，得到
$C^{\mathcal E^\bullet}$ 到 $E^2$ 的嵌入。\emph{虚拟基本类}定义为
$s^*[C^{\mathcal E^\bullet}]$，其中 $s$ 为 $E^2$ 的零截面。立即可见，它是 $Z$ 上维数
$\operatorname{rk}(E^1)-\operatorname{rk}(E^2)$ 的圈类。
这一差值可识别为 $Z$ 的预期维数；更准确地说，这也成为预期维数的一种定义。
由完美切—障碍复形的定义，显然这一差值完全由切—障碍复形决定，而与 $\mathcal E^i$ 的选择无关。
不那么显然、但由 \SourceCite{LTi2} 证明的是：虚拟基本类本身也与 $\mathcal E^i$ 和 Kuranishi 映射的选择无关。

\SourceExampleNumber{7.1.4}{1}
\begin{example}\label{ex:zero-scheme-obstruction-theory}
为了将这些内容具体化，再回到光滑概形 $Y$、向量丛 $E$ 与截面 $s\in H^0(E)$ 的情形。
令 $\mathcal F$ 为由 $Z=Z(s)$ 表示的函子。换言之，
\[
\mathcal F(S)=\{\text{所有满足 }f^*(s)=0\text{ 的态射 }f:S\to Y\}.
\]
给定 $\mathcal F(S)$ 的元素 $f:S\to Y$，我们断言：由
\[
f^*\Theta_Y\longrightarrow f^*\mathcal O(E)
\]
的核与余核定义的复形 $T^\bullet\mathcal F(f)$，满足完美切—障碍复形所需的条件。
上述复形中的映射，是截面 $s:Y\to E$ 的微分与自然投影 $T_E\to E$ 的复合；后者将 $E$ 的切丛沿零截面投影到零截面的法丛，而这个法丛同构于 $E$ 本身。
它显然由全局定义在 $Z$ 上的复形诱导，即
\begin{equation}\label{eq:zero-scheme-tangent-obstruction-complex}
\Theta_Y|_Z\longrightarrow\mathcal O(E)|_Z,
\end{equation}
这个复形充当上面的 $\mathcal E^\bullet$。按照 \SourceCite{LTi2} 的构造，虚法锥
$C^{\mathcal E^\bullet}$ 与法锥 $C_ZY$ 相同，其到 $E$ 的嵌入也正是刚才描述的嵌入。
于是，虚拟基本类就是 $s^*[C_ZY]$。这表明，虚拟基本类推广了\hyperref[subsec:normal-cone]{第~\ref*{subsec:normal-cone} 小节}的法锥构造。
\end{example}

到目前为止，只在函子 $\mathcal F$ 可由概形表示、且复形 $\mathcal E^\bullet$ 全局存在的特殊情形中构造了虚拟基本类。
\SourcePage{201}
\SourceCite{LTi2} 表明，这些假设可以放宽。特别地，对于
$\overline{\mathcal M}_{g,n}(X,\beta)$ 这种具有粗模空间的不可表示函子，也能应用这一构造定义虚拟基本类。

为了看清对 $\overline{\mathcal M}_{g,n}(X,\beta)$ 如何实现，考虑仿射概形 $S$ 上亏格 $g$、
带 $n$ 个标记的稳定映射族 $f:\mathcal C\to X$，见定义~\ref{def:stable-map-family}。
以 $D\subset\mathcal C$ 表示各截面 $s_i:S\to\mathcal C$ 的像的并。
Li 与 Tian 考察的切—障碍复形，将 $f$ 对应于函子复形
\begin{equation}\label{eq:stable-map-relative-obstruction-complex}
T^i\overline{\mathcal M}_{g,n}(X,\beta)(f)(\mathcal N)
=\operatorname{Ext}_{\mathcal C/S}^i
\left(f^*\Omega_X^1\longrightarrow\Omega_{\mathcal C/S}(D),
\mathcal O_{\mathcal C}\otimes_{\mathcal O_S}\mathcal N\right),
\end{equation}
其中 $i=1,2$。Li 与 Tian 证明，这是完美复形。因此，将
$\overline{\mathcal M}_{g,n}(X,\beta)$ 的\emph{虚拟基本类}定义为上述完美切—障碍复形的虚拟基本类。
这一类记作 $[\overline{\mathcal M}_{g,n}(X,\beta)]^{\mathrm{virt}}$，其维数为
$\operatorname{rk}(E^1)-\operatorname{rk}(E^2)$。由式~\eqref{eq:stable-map-ext-sequence} 紧前面的表达式，可见
$[\overline{\mathcal M}_{g,n}(X,\beta)]^{\mathrm{virt}}$ 具有预期维数~\eqref{eq:stable-map-expected-dimension}。

下面简要介绍 Behrend 与 Fantechi \SourceCite{BF} 的相关虚拟基本类构造。不会深入描述这一构造，因为这需要大幅转入叠的理论。叠的语言仍参见 \SourceCite{DM} 或 \SourceCite{BEFFGK}，代数叠的相交理论见 \SourceCite{Vistoli}。

Deligne--Mumford 叠是对角线非分歧的代数叠。\SourceCite{BF} 的构造始于以下观察：
任何 Deligne--Mumford 叠 $\mathcal M$ 都有一个\emph{内蕴法锥} $C_{\mathcal M}$。
它是 $\mathcal M$ 上纯维数为 $0$ 的叠。可如下局部描述内蕴法锥。
由于 $\mathcal M$ 有由概形给出的平展开覆盖，可以在一个这样的概形 $U$ 上工作。
选取嵌入 $U\hookrightarrow W$，其中 $W$ 光滑。令 $\mathcal I$ 为 $U$ 在 $W$ 中的理想层，
照常用式~\eqref{eq:normal-cone-definition} 的右端定义法锥 $C_UW$。微分映射
$\mathcal I\to\Omega_W^1$ 将 $f$ 映为 $df$，诱导映射
\[
\bigoplus_k\mathcal I^k/\mathcal I^{k+1}
\longrightarrow\bigoplus_k\operatorname{Sym}^k
(\Omega_W^1/\mathcal I\Omega_W^1).
\]
取 $\operatorname{Spec}$，得到映射
\[
T_W|_U\longrightarrow C_UW,
\]
因为 $T_W=\operatorname{Spec}(\bigoplus_{k=0}^\infty\operatorname{Sym}^k\Omega_W^1)$。
它给出一个作用，从而可取叠商；内蕴法锥 $C_U$ 定义为 $C_UW$ 关于 $T_W|_U$ 的叠商。
$C_U$ 与 $W$ 的选择无关，所以选取平展覆盖 $\mathcal M=\bigcup_iU_i$，其中各 $U_i$ 为概形时，
局部内蕴法锥 $C_{U_i}$ 能粘合成 $C_{\mathcal M}$。\index{内蕴法锥|hyperpage}

在上述构造中，把法锥 $C_UW$ 换成法层，也得到\emph{内蕴法层} $\mathcal N_{\mathcal M}$ 的构造。
事实上，$C_{\mathcal M}$ 嵌入 $\mathcal N_{\mathcal M}$。这一嵌入局部由满射
\[
\bigoplus_k\operatorname{Sym}^k(\mathcal I/\mathcal I^2)
\longrightarrow\bigoplus_k\mathcal I^k/\mathcal I^{k+1}
\]
诱导。

现在，$\mathcal M$ 有余切复形 $L^\bullet_{\mathcal M}$ \SourceCite{Illusie}。
粗略地说，$h^0(L^\bullet_{\mathcal M})$ 控制形变，$h^{-1}(L^\bullet_{\mathcal M})$ 控制障碍。
在 \SourceCite{BF} 中，$\mathcal M$ 的\emph{障碍理论}粗略指 $\mathcal M$ 上的层复形
$\mathcal E^\bullet$ 及态射 $\mathcal E^\bullet\to L^\bullet_{\mathcal M}$，
它在 $h^0$ 上为同构，在 $h^{-1}$ 上为满射。更精确的处理应在导出范畴中进行。
注意，这与 \SourceCite{LTi2} 中的切—障碍复形类似：
\SourcePage{202}
切空间 $T^1\mathcal F$ 固定，但障碍 $T^2\mathcal F$ 有一定的灵活性。

若 $\mathcal E^0\to\mathcal E^1$ 为 $\mathcal M$ 上 Abel 层的两项复形，可取
$\mathcal E^0$ 在 $\mathcal E^1$ 上的作用的叠商。上面构造内蕴法锥时，已经见过这种商的例子。
现在，若 $\mathcal E^\bullet$ 为任意长度的复形，按照 \SourceCite{BF}，
可将 $h^1/h^0(\mathcal E^\bullet)$ 定义为：$\mathcal E^1\to\mathcal E^2$ 的核，
关于 $\mathcal E^{-1}\to\mathcal E^0$ 的余核的叠商。

\begin{proposition}\label{prop:intrinsic-normal-sheaf}
\[
h^1/h^0\bigl((L^\bullet_{\mathcal M})^\vee\bigr)
\simeq\mathcal N_{\mathcal M}.
\]
\end{proposition}
\begin{proof}
这由 \SourceCite[第 66--69 页]{BF} 得到。符号 $(L^\bullet_{\mathcal M})^\vee$ 表示导出范畴中的“对偶”。
\end{proof}

现在可以构造虚拟基本类。若 $h^1/h^0((\mathcal E^\bullet)^\vee)$ 在 $\mathcal M$ 上光滑，
则障碍理论 $\mathcal E^\bullet$ 是\emph{完美的}。由命题~\ref{prop:intrinsic-normal-sheaf} 及完美障碍理论的定义，
$\mathcal N_{\mathcal M}$ 嵌入 $h^1/h^0((\mathcal E^\bullet)^\vee)$，因而 $C_{\mathcal M}$ 也嵌入其中。
令 $C$ 为 $(E^{-1})^*$ 与 $C_{\mathcal M}$ 在 $h^1/h^0((\mathcal E^\bullet)^\vee)$ 上的纤维积，其中
$\mathcal O(E^{-1})=\mathcal E^{-1}$。这是向量丛 $(E^{-1})^*$ 中的普通锥。
\emph{虚拟基本类}定义为 $C$ 与 $(E^{-1})^*$ 的零截面的交。

下面用 \SourceCite{BF} 的方法继续前面的例子。

\begingroup
\SourceExampleNumber{7.1.4}{1}
\renewcommand{\theHexam}{7.1.4.revisited.\arabic{exam}}
\begin{example}[再访]\label{ex:zero-scheme-bf-obstruction}
回到向量丛 $E$ 的截面 $s$ 的零点轨迹 $Z\subset Y$。令 $\mathcal E^\bullet$ 为
式~\eqref{eq:zero-scheme-tangent-obstruction-complex} 的对偶复形
$\mathcal O(E^*)|_Z\to\Omega_Y^1|_Z$，两项的次数分别为 $-1$ 和 $0$。
由于 $Y$ 光滑，有 $L_Y^\bullet\simeq\Omega_Y^1$。
截面 $s$ 给出的方程诱导 $\mathcal O(E^*)|_Z\to\mathcal I/\mathcal I^2$，
连同次数 $0$ 上 $\Omega_Y^1|_Z$ 的自然映射，得到 $\mathcal E^\bullet\to L_Z^\bullet$。
这使 $\mathcal E^\bullet$ 成为完美障碍理论。
于是，$h^1/h^0((\mathcal E^\bullet)^\vee)$ 是 $E|_Z$ 关于 $\Theta_Y|_Z$ 的叠商，
可与式~\eqref{eq:zero-scheme-tangent-obstruction-complex} 比较。
接着把虚拟基本类算为丛 $E|_Z$ 中锥 $C$ 与零截面的交。
由内蕴法锥的定义，$C=C_ZY$，到 $E|_Z$ 的嵌入也与本小节开头所给的嵌入相同。
因此，虚拟基本类确实等于 $s^*[C_ZY]$。
\end{example}
\endgroup

现在将这一构造用于 $\overline{\mathcal M}_{g,n}(X,\beta)$。
沿用 \SourceCite{Behrend}，定义虚拟基本类的一个相对版本，它更容易应用。
令 $\mathfrak M_{g,n}$ 为满足定义~\ref{def:stable-map} 中条件 (i)、(ii) 的带 $n$ 个标记的亏格 $g$ 曲线的叠，
这些曲线称为\emph{预稳定曲线}。于是有叠映射
\begin{equation}\label{eq:stable-map-prestable-forgetting}
F:\overline{\mathcal M}_{g,n}(X,\beta)\longrightarrow\mathfrak M_{g,n},
\end{equation}
它忘掉映射，但不作稳定化。上述构造的相对版本表明，$F$ 有内蕴法锥 $C_F$ 与法层 $\mathcal N_F$。
再考虑叠映射
\[
\pi_{n+1}:\overline{\mathcal M}_{g,n+1}(X,\beta)
\longrightarrow\overline{\mathcal M}_{g,n}(X,\beta),
\]
它忘掉最后一个点 $p_{n+1}$，并收缩因此出现的任何不稳定分支。
如\hyperref[subsec:gravitational-correlator-definition]{第~\ref*{subsec:gravitational-correlator-definition} 小节}所说明、并由 \SourceCite{BM} 证明，泛\IndexLocation{indexloc:1282:202}稳定映射在叠的意义下由 $\pi_{n+1}$ 与\IndexLocation{indexloc:0276:202}赋值映射
\[
e_{n+1}:\overline{\mathcal M}_{g,n+1}(X,\beta)\longrightarrow X
\]
组成。

\SourcePage{203}
所需完美障碍理论为
$\mathcal E^\bullet=(R\pi_{n+1*}e_{n+1}^*\Theta_X)^\vee$。
于是，$C_F$ 嵌入
\[
h^1/h^0((\mathcal E^\bullet)^\vee)=h^1/h^0(R\pi_{n+1*}e_{n+1}^*\Theta_X).
\]
同样，可将它拉回为普通向量丛中的普通锥，与零截面相交，得到\emph{相对虚拟基本类}
\[
[\overline{\mathcal M}_{g,n}(X,\beta)/\mathfrak M_{g,n},
R\pi_{n+1*}e_{n+1}^*\Theta_X]
\]
位于双变 Chow 群
$A^*(\overline{\mathcal M}_{g,n}(X,\beta)\to\mathfrak M_{g,n})$ 中。
双变 Chow 群见 \SourceCite{Fulton2} 与 \SourceCite{Vistoli}。
最后，将 $\overline{\mathcal M}_{g,n}(X,\beta)$ 的\emph{虚拟基本类}定义为
\[
[\overline{\mathcal M}_{g,n}(X,\beta)]^{\mathrm{virt}}
=[\mathfrak M_{g,n}]\cap
[\overline{\mathcal M}_{g,n}(X,\beta)/\mathfrak M_{g,n},
R\pi_{n+1*}e_{n+1}^*\Theta_X],
\]
它位于 $A_*(\overline{\mathcal M}_{g,n}(X,\beta))$ 中。

注意，上述各种虚拟基本类定义都位于 Chow 群或双变 Chow 群中。
本章余下部分经常将虚拟基本类与其关联的同调类识别，这一过程会舍弃部分信息。
这样做是为了与\hyperref[sec:gw-symplectic-definition]{第~\ref*{sec:gw-symplectic-definition} 节}讨论的辛几何版本的虚拟基本类比较，
也为了在\hyperref[chap:localization]{第~\ref*{chap:localization} 章}及后续章节中避免使用等变 Chow 群。

\subsection{计算虚拟基本类}\label{subsec:computing-virtual-class}
\index{虚拟基本类!计算|hyperpage}\index{Euler类@Euler 类!虚拟基本类|hyperpage}
这里不穷尽虚拟基本类的计算与运算，而是介绍四种理解得较为清楚的情形。
首先是一种特别容易求出 $[\overline M_{0,n}(X,\beta)]^{\mathrm{virt}}$ 的情况。

\begin{definition}\label{def:convex-variety}
\SourceCite{KoM1} 若光滑射影簇 $X$ 对所有亏格 $0$ 的稳定映射 $f:C\to X$，都满足
$H^1(C,f^*T_X)=0$，则称 $X$ 为\emph{凸的}。
\end{definition}
\index{凸簇|hyperpage}
当 $X$ 凸时，可证明 $\overline M_{0,n}(X,\beta)$ 是具有预期维数的轨形 \SourceCite{FP}。
特别地，有基本类 $[\overline M_{0,n}(X,\beta)]\in H_*(\overline M_{0,n}(X,\beta),\mathbb Q)$。
不过，凸性还意味着到 $X$ 的亏格 $0$ 稳定映射的模函子无障碍。
使用上述任一种方法，都化为平凡丛中的平凡相交，故本情形中
$[\overline M_{0,n}(X,\beta)]^{\mathrm{virt}}=[\overline M_{0,n}(X,\beta)]$。
射影空间以及更一般的齐性空间，都是凸的。

\hyperref[subsec:moduli-to-gw]{第~\ref*{subsec:moduli-to-gw} 小节}讨论了 $\overline M_{g,n}(X,\beta)$ 是具有预期维数的轨形的情形。
对这样的 $X$，仍有
\[
[\overline M_{g,n}(X,\beta)]^{\mathrm{virt}}=[\overline M_{g,n}(X,\beta)].
\]
原因是：用 \SourceCite{LTi2} 的语言，切—\IndexLocation{indexloc:1328:203}障碍复形由关联模叠
$\overline{\mathcal M}_{g,n}(X,\beta)$ 的切层诱导；用 \SourceCite{BF} 的语言，
完美障碍理论由叠的余切层诱导。与上一段类似，得到所需等式。

接着计算五次三维簇的一个虚拟基本类。

\SourceExampleNumber{7.1.5}{1}
\begin{example}\label{ex:quintic-virtual-class}
设 $V\subset\mathbb P^4$ 为光滑五次三维簇，$d>0$。本例目标是构造虚拟基本类
$[\overline M_{0,0}(V,d\ell)]^{\mathrm{virt}}$。
先证明自然映射 $\overline M_{0,0}(V,d\ell)\to\overline M_{0,0}(\mathbb P^4,d)$ 为嵌入。
为此，令 $s$ 为 $\mathcal O_{\mathbb P^4}(5)$ 上在 $V$ 上消失的截面，考虑函子
\begin{equation}\label{eq:quintic-stable-map-zero-functor}
F(S)=\{f:\mathcal C\to\mathbb P^4:
f\text{ 为 }S\text{ 上亏格 }0\text{、次数 }d\text{ 的稳定映射，且 }f^*(s)=0\}/\sim,
\end{equation}
其中 $\sim$ 表示 $S$ 上稳定映射的同构。这是关联于
$\overline M_{0,0}(V,d\ell)$ 的函子，因此得到嵌入
$\overline M_{0,0}(V,d\ell)\hookrightarrow\overline M_{0,0}(\mathbb P^4,d)$。

\SourcePage{204}
可如下描述这一嵌入。令 $\mathcal V_d$ 为叠
$\overline{\mathcal M}_{0,0}(\mathbb P^4,d)$ 上的秩 $5d+1$ 向量丛，
其在稳定映射 $f:C\to\mathbb P^4$ 处的纤维为 $H^0(C,f^*\mathcal O_{\mathbb P^4}(5))$。
定义 $V$ 的截面 $s$ 确定 $\overline{\mathcal M}_{0,0}(\mathbb P^4,d)$ 上
$\mathcal V_d$ 的截面 $\widetilde s$；由式~\eqref{eq:quintic-stable-map-zero-functor}，
$\overline{\mathcal M}_{0,0}(V,d\ell)$ 是 $\widetilde s$ 的零点轨迹。
这里使用叠，是因为 $\mathcal V_d$ 未必是轨形 $\overline M_{0,0}(\mathbb P^4,d)$ 上的向量丛。
这样表示 $\overline{\mathcal M}_{0,0}(V,d\ell)$，说明我们几乎处于\hyperref[subsec:normal-cone]{第~\ref*{subsec:normal-cone} 小节}的情形：
有向量丛 $E=\mathcal V_d$，底空间 $Y=\overline{\mathcal M}_{0,0}(\mathbb P^4,d)$ 光滑，
这里是叠意义下的光滑，而非作为簇光滑，且 $Z(\widetilde s)=\overline{\mathcal M}_{0,0}(V,d\ell)$
是截面 $\widetilde s$ 的零点轨迹。于是，法锥构造给出 $\widetilde s^*[C_ZY]$。
它是零圈，因为 $\mathcal V_d$ 的秩为 $5d+1$，由式~\eqref{eq:stable-map-expected-dimension}，
$\overline M_{0,0}(\mathbb P^4,d)$ 的维数也是 $5d+1$。

我们断言，$\widetilde s^*[C_ZY]$ 就是虚拟基本类
$[\overline M_{0,0}(V,d\ell)]^{\mathrm{virt}}$。采用 Li 与 Tian 的方法证明。
由例题~\ref{ex:zero-scheme-obstruction-theory}，$\widetilde s^*[C_ZY]$ 为下列映射的核与余核
给出的完美切—障碍复形的虚拟基本类，参见式~\eqref{eq:zero-scheme-tangent-obstruction-complex}：
\begin{equation}\label{eq:quintic-zero-obstruction-complex}
\Theta_{\overline{\mathcal M}_{0,0}(\mathbb P^4,d)}
\big|_{\overline{\mathcal M}_{0,0}(V,d\ell)}
\longrightarrow
\mathcal V_d\big|_{\overline{\mathcal M}_{0,0}(V,d\ell)}.
\end{equation}
因此，只需证明式~\eqref{eq:quintic-zero-obstruction-complex} 的核与余核，与
式~\eqref{eq:stable-map-relative-obstruction-complex} 为
$\overline{\mathcal M}_{0,0}(V,d\ell)$ 定义的切—障碍复形一致。

为此，令 $f:\mathcal C\to V$ 为概形 $S$ 上的稳定映射，$\pi:\mathcal C\to S$ 为结构映射。
考虑复形的短正合列
\[
\begin{array}{ccccccccc}
0&\longrightarrow&f^*\mathcal O_{\mathbb P^4}(-5)&\longrightarrow&
f^*\Omega_{\mathbb P^4}^1&\longrightarrow&f^*\Omega_V^1&\longrightarrow&0\\
&&\downarrow&&\downarrow&&\downarrow&&\\
0&\longrightarrow&0&\longrightarrow&
\Omega_{\mathcal C/S}^1&=&\Omega_{\mathcal C/S}^1&\longrightarrow&0.
\end{array}
\]
再考察相对 $\operatorname{Ext}$ 的长正合列：
\begin{equation}\label{eq:quintic-relative-ext-sequence}
\begin{aligned}
0&\longrightarrow
\operatorname{Ext}_{\mathcal C/S}^1(f^*\Omega_V^1\to\Omega_{\mathcal C/S}^1,\mathcal O_{\mathcal C})\\
&\longrightarrow
\operatorname{Ext}_{\mathcal C/S}^1(f^*\Omega_{\mathbb P^4}^1\to\Omega_{\mathcal C/S}^1,\mathcal O_{\mathcal C})\\
&\longrightarrow
\operatorname{Ext}_{\mathcal C/S}^0(f^*\mathcal O_{\mathbb P^4}(-5),\mathcal O_{\mathcal C})\\
&\longrightarrow
\operatorname{Ext}_{\mathcal C/S}^2(f^*\Omega_V^1\to\Omega_{\mathcal C/S}^1,\mathcal O_{\mathcal C})
\longrightarrow0.
\end{aligned}
\end{equation}
末尾的 $\to0$ 来自 $\mathbb P^4$ 的凸性。

考察式~\eqref{eq:quintic-relative-ext-sequence} 的非零项。第一项是
$\overline{\mathcal M}_{0,0}(V,d\ell)$ 的切空间拉回到 $S$，
第二项是 $\overline{\mathcal M}_{0,0}(\mathbb P^4,d)$ 的切空间拉回到 $S$，
第三项是 $\mathcal V_d$ 拉回到 $S$，第四项为障碍空间 $T^2$。
由式~\eqref{eq:quintic-relative-ext-sequence} 的正合性，可将
式~\eqref{eq:quintic-zero-obstruction-complex} 的核识别为
$\operatorname{Ext}_{\mathcal C/S}^1(f^*\Omega_V^1\to\Omega_{\mathcal C/S}^1,\mathcal O_{\mathcal C})$，
将余核识别为
$\operatorname{Ext}_{\mathcal C/S}^2(f^*\Omega_V^1\to\Omega_{\mathcal C/S}^1,\mathcal O_{\mathcal C})$。
这正是式~\eqref{eq:stable-map-relative-obstruction-complex} 的切—障碍复形，断言得证。

细心的读者会发现，这一断言稍微宽松地使用了记号。丛 $\mathcal V_d$ 及其 Euler 类
仅定义在叠 $\overline{\mathcal M}_{0,0}(\mathbb P^4,d)$ 上，而非空间 $\overline M_{0,0}(\mathbb P^4,d)$ 上，
但我们却在粗模空间 $\overline M_{0,0}(V,d\ell)$ 上定义虚拟基本类。
这是因为一旦在叠上得到 Euler 类，就能将其推到粗模空间，而不必要求
$\overline M_{0,0}(V,d)$ 上存在相关的丛。全书中将继续采用这种宽松记法。

等式 $[\overline M_{0,0}(V,d\ell)]^{\mathrm{virt}}=\widetilde s^*[C_ZY]$ 有一个很好的推论。
令 $i:\overline{\mathcal M}_{0,0}(V,d\ell)\hookrightarrow
\overline{\mathcal M}_{0,0}(\mathbb P^4,d)$ 为上述嵌入。引理~\ref{lem:normal-cone-euler-class} 蕴含
\begin{equation}\label{eq:quintic-virtual-class-pushforward}
i_*([\overline{\mathcal M}_{0,0}(V,d\ell)]^{\mathrm{virt}})
=i_*(\widetilde s^*[C_ZY])
=c_{5d+1}(\mathcal V_d)\cap[\overline{\mathcal M}_{0,0}(\mathbb P^4,d)].
\end{equation}
\SourcePage{205}
这一等式成立于 $\overline{\mathcal M}_{0,0}(\mathbb P^4,d)$ 上；
推到 $\overline M_{0,0}(\mathbb P^4,d)$ 后，得到代表丛 $\mathcal V_d$ 的 Euler 类的零圈。
由此将得到一个很好的 Gromov--Witten 不变量公式，例题~\ref{ex:quintic-gw-euler-integral} 将研究它。

按例题~\ref{ex:quintic-lines-normal-cone} 的记号，$G(2,5)$ 上的丛
$\operatorname{Sym}^5U^*$ 就是这里考虑的 $\overline M_{0,0}(\mathbb P^4,1)$ 上的丛
$\mathcal V_1$。因此，前一例题的嵌入
$\overline M_{0,0}(V,\ell)\hookrightarrow G(2,5)$，由上述
$\overline{\mathcal M}_{0,0}(V,d\ell)$ 的描述得到。
\end{example}

计算虚拟基本类的第三个方面涉及“超额维数”。在相交理论中，当一个交的维数大于预期维数时，
就出现超额维数。在某些良好情形下，可用所谓超额法丛计算相交类。
清楚的讨论见 \SourceCite[第 6.3 节]{Fulton2}。

“超额维数”的思想也适用于当前情形。将
$\overline M_{g,n}(X,\beta)$ 的\emph{超额维数}定义为实际维数
$\dim\overline M_{g,n}(X,\beta)$ 与预期维数~\eqref{eq:stable-map-expected-dimension} 之差。
若超额维数为正，等式
$[\overline M_{g,n}(X,\beta)]^{\mathrm{virt}}=[\overline M_{g,n}(X,\beta)]$ 显然不可能成立，
但有时仍能用 Euler 类计算虚拟基本类。
若 $F:\overline{\mathcal M}_{g,n}\to\mathfrak M_{g,n}$ 为式~\eqref{eq:stable-map-prestable-forgetting}
中的叠映射，则有以下结果。\index{超额维数|hyperpage}

\begin{proposition}\label{prop:excess-obstruction-euler-class}
若定义 $[\overline M_{g,n}(X,\beta)]^{\mathrm{virt}}$ 时使用的层
$R^1\pi_{n+1*}e_{n+1}^*\Theta_X$ 局部自由，其秩 $e$ 等于超额维数，则
\[
[\overline M_{g,n}(X,\beta)]^{\mathrm{virt}}
=c_e(R^1\pi_{n+1*}e_{n+1}^*\Theta_X)
\cap\bigl([\overline{\mathcal M}_{g,n}(X,\beta)/\mathfrak M_{g,n}]
\cap[\mathfrak M_{g,n}]\bigr),
\]
其中 $[\overline{\mathcal M}_{g,n}(X,\beta)/\mathfrak M_{g,n}]
\in A^*(\overline{\mathcal M}_{g,n}(X,\beta)\to\mathfrak M_{g,n})$
是 \SourceCite{Fulton2} 意义下的基本类。
\end{proposition}
\begin{proof}
见 \SourceCite{Behrend,Getzler1}。
\end{proof}

讨论亏格 $1$ 的 Gromov--Witten 不变量时，将在例题~\ref{ex:plane-genus-one} 使用这一命题。
注意，即使 $\overline M_{g,n}(X,\beta)$ 光滑，命题~\ref{prop:excess-obstruction-euler-class} 的假设也未必成立。
原因在于：$\overline M_{g,n}(X,\beta)$ 中映射的障碍并不位于 $H^1(C,f^*\Theta_X)$ 本身，
而是位于它的商
$\operatorname{Ext}_C^2(f^*\Omega_X^1\to\Omega_C^1(\sum_i p_i),\mathcal O_C)$；
这一空间处理起来更微妙。\hyperref[subsec:nodal-quintic-double-covers]{第~\ref*{subsec:nodal-quintic-double-covers} 小节}将给出说明。

第四个、也是最后一个观察，涉及 $\overline M_{g,n}(X,\beta)$ 与
$\overline M_{g,n-1}(X,\beta)$ 的虚拟基本类之间的相容性。考虑交换图
\begin{figure}[H]
\centering
\includegraphics[width=0.568\textwidth]{7-1-forgetting-diagram.png}
\end{figure}
其中 $\phi_i$ 忘掉第一个标记点并作稳定化，$\pi_i$ 忘掉最后一个标记点，
$e_i:\overline{\mathcal M}_{g,i}(X,\beta)\to X$ 在最后一个点处赋值。
由基变换，
\[
R^1\pi_{n+1*}e_{n+1}^*\Theta_X
\simeq\phi_n^*R^1\pi_{n*}e_n^*\Theta_X.
\]
\SourcePage{206}
由此得到
\begin{equation}\label{eq:virtual-class-forgetting-compatibility}
[\overline M_{g,n}(X,\beta)]^{\mathrm{virt}}
=\phi_n^*[\overline M_{g,n-1}(X,\beta)]^{\mathrm{virt}}.
\end{equation}
换言之，第 $n$ 个标记点的形变理论无障碍，见 \SourceCite[公理 IV]{Behrend}。
相容性~\eqref{eq:virtual-class-forgetting-compatibility} 将在\hyperref[chap:quantum-differential-equations]{第~\ref*{chap:quantum-differential-equations} 章}中有用。
应指出，忘却映射 $\phi_n:\overline M_{g,n}(X,\beta)\to\overline M_{g,n-1}(X,\beta)$ 并非总存在。
问题是，当 $\beta=0$ 时，忘掉一个标记点并作稳定化，可能使整条曲线消失。
例如，$\phi_3:\overline M_{0,3}(X,0)\to\overline M_{0,2}(X,0)$ 不存在，
因为 $\overline M_{0,3}(X,0)=\overline M_{0,3}\times X$，而 $\overline M_{0,2}(X,0)=\varnothing$。
幸运的是，不难证明：若 $n+2g\geq4$，或者 $\beta\ne0$ 且 $n\geq1$，则 $\phi_n$ 存在。
对 $\pi_n$ 也一样。

\subsection{定义 Gromov--Witten 类与不变量}\label{subsec:gw-algebraic-classes}
\index{Gromov--Witten不变量@Gromov--Witten 不变量!代数几何中的定义|hyperpage}
有了虚拟基本类 $\xi=[\overline M_{g,n}(X,\beta)]^{\mathrm{virt}}$，终于能严格定义
Gromov--Witten 类与不变量。回忆，由式~\eqref{eq:stable-map-evaluation-forgetting}，
有映射 $\pi:\overline M_{g,n}(X,\beta)\to X^n\times\overline M_{g,n}$；
以 $p_i$，$i=1,2$，表示到 $X^n\times\overline M_{g,n}$ 第 $i$ 个因子的投影。

\begin{definition}\label{def:gw-algebraic-invariants}
设 $\beta\in H_2(X,\mathbb Z)$ 为同调类，$\alpha_1,\ldots,\alpha_n\in H^*(X,\mathbb Q)$ 为上同调类。
\begin{enumerate}[label=(\roman*)]
\item 若 $n+2g\geq3$，则 Gromov--Witten 类 $I_{g,n,\beta}(\alpha_1,\ldots,\alpha_n)$ 为
$H^*(\overline M_{g,n},\mathbb Q)$ 中由下式定义的上同调类：
\[
I_{g,n,\beta}(\alpha_1,\ldots,\alpha_n)
=PD^{-1}p_{2*}\bigl(p_1^*(\alpha_1\otimes\cdots\otimes\alpha_n)\cap\pi_*(\xi)\bigr),
\]
其中 $\xi=[\overline M_{g,n}(X,\beta)]^{\mathrm{virt}}$ 为
$\overline M_{g,n}(X,\beta)$ 的虚拟基本类，$PD$ 为 Poincaré 对偶。
\item 对 $n,g\geq0$，Gromov--Witten 不变量
$\langle I_{g,n,\beta}\rangle(\alpha_1,\ldots,\alpha_n)$ 为由
\[
\langle I_{g,n,\beta}\rangle(\alpha_1,\ldots,\alpha_n)
=\int_\xi e_1^*(\alpha_1)\cup\cdots\cup e_n^*(\alpha_n)
\]
定义的有理数，其中 $e_i$ 如式~\eqref{eq:gw-evaluation-integral} 所述。
\end{enumerate}
\end{definition}
注意，这些定义与式~\eqref{eq:gw-fundamental-class-definition} 和~\eqref{eq:gw-evaluation-integral} 相同，
只是 $\xi$ 改为虚拟基本类。此外，若 $n+2g\geq3$，可证明
\[
\langle I_{g,n,\beta}\rangle(\alpha_1,\ldots,\alpha_n)
=\int_{\overline M_{g,n}}I_{g,n,\beta}(\alpha_1,\ldots,\alpha_n).
\]
因此，$n+2g\geq3$ 时，Gromov--Witten 不变量由相应的 Gromov--Witten 类确定。

\hyperref[sec:gw-properties]{第~\ref*{sec:gw-properties} 节}将研究 Gromov--Witten 类与不变量的性质，\hyperref[sec:computing-gw-i]{第~\ref*{sec:computing-gw-i} 节}将计算一些例子。
还应说明：虽然这里用 $\mathbb Q$ 上的同调与上同调处理 Gromov--Witten 类，
许多文献使用的却是 Chow 群。\SourceCite{Behrend} 使用\hyperref[subsec:virtual-fundamental-class]{第~\ref*{subsec:virtual-fundamental-class} 小节}先前讨论的相对虚拟基本类版本，
给出了类似的 Gromov--Witten 不变量定义。

本节以 \SourceCite{Kontsevich2} 给出的五次三维簇某些 Gromov--Witten 不变量的公式结束。
这一结果将在\hyperref[chap:mirror-theorem]{第~\ref*{chap:mirror-theorem} 章}发挥重要作用。

\SourcePage{207}
\SourceExampleNumber{7.1.6}{1}
\begin{example}\label{ex:quintic-gw-euler-integral}
设 $V\subset\mathbb P^4$ 为光滑五次三维簇，仍以 $\ell\in H_2(V,\mathbb Z)$ 表示直线的类。
例题~\ref{ex:quintic-virtual-class} 已说明，虚拟基本类
$[\overline M_{0,0}(V,d\ell)]^{\mathrm{virt}}$ 为零圈。
因此，由定义~\ref{def:gw-algebraic-invariants}，
\[
\langle I_{0,0,d\ell}\rangle
=\int_{[\overline M_{0,0}(V,d\ell)]^{\mathrm{virt}}}1
=\text{零圈 }[\overline M_{0,0}(V,d\ell)]^{\mathrm{virt}}\text{ 的次数}.
\]
可通过嵌入 $i:\overline M_{0,0}(V,d\ell)\hookrightarrow\overline M_{0,0}(\mathbb P^4,d)$，
在 $\overline M_{0,0}(\mathbb P^4,d)$ 上计算这一次数。
由式~\eqref{eq:quintic-virtual-class-pushforward}，
\[
i_*[\overline M_{0,0}(V,d\ell)]^{\mathrm{virt}}
=c_{5d+1}(\mathcal V_d)\cap[\overline M_{0,0}(\mathbb P^4,d)],
\]
其中 $\mathcal V_d$ 为 $\overline{\mathcal M}_{0,0}(\mathbb P^4,d)$ 上的向量丛，
在亏格 $0$ 的稳定映射 $f:C\to\mathbb P^4$ 处的纤维为 $H^0(C,f^*\mathcal O_{\mathbb P^4}(5))$。
由于 $c_{5d+1}(\mathcal V_d)$ 为最高次数的上同调类，
零圈 $c_{5d+1}(\mathcal V_d)\cap[\overline M_{0,0}(\mathbb P^4,d)]$ 的次数，
由在 $\overline M_{0,0}(\mathbb P^4,d)$ 的基本类上积分给出。因此
\[
\langle I_{0,0,d\ell}\rangle
=\int_{\overline M_{0,0}(\mathbb P^4,d)}c_{5d+1}(\mathcal V_d).
\]
这一公式首次见于 \SourceCite[第 2.2 节]{Kontsevich2}。
这里稍宽松地把 $\mathcal V_d$ 写成定义在 $\overline M_{0,0}(\mathbb P^4,d)$ 上的丛；
如前所述，其 Euler 类在 $\overline M_{0,0}(\mathbb P^4,d)$ 上给出良定义的零圈。
\hyperref[subsec:kontsevich-localization]{第~\ref*{subsec:kontsevich-localization} 小节}将说明如何用等变上同调与局部化计算这一积分。
\end{example}

\section{辛几何中的定义}\label{sec:gw-symplectic-definition}
\index{Gromov--Witten不变量@Gromov--Witten 不变量!辛几何中的定义|hyperpage}
本节简要概述 Gromov--Witten 不变量的辛几何表述所用的思想与技术。
这些定义以辛几何中 $J$ 全纯映射的理论为基础，带来了对量子上同调的最早严格处理。
部分基本思想见 \SourceCite{Witten3}；对半正辛流形的亏格 $0$ 辛 Gromov--Witten 不变量的首次显式处理，
则归功于 \SourceCite{Ruan2} 与 \SourceCite{RT1}。\SourceCite{MS} 可作为这一方法的入门书。
半正情形的更高亏格不变量，先在 \SourceCite{RT1} 中为固定复结构定义，
随后在 \SourceCite{RT2} 中推广到变化的复结构。
现在已有若干一般辛流形的 Gromov--Witten 不变量定义，见 \SourceCite{LTi3,Siebert1,Ruan3,FO}。
本节概述省略的许多细节，均可查阅这些文献。

\subsection{辛流形与 $J$ 全纯映射}\label{subsec:j-holomorphic-maps}
\index{几乎复结构|hyperpage}\index{几乎Kähler流形@几乎 Kähler 流形|hyperpage}
设 $M$ 为实维数 $2n$ 的紧辛流形。即 $M$ 上有闭 $2$ 形式 $\omega$，使 $\omega^n$ 在每一点非零。
每个 Kähler 流形都有由 Kähler 形式给出的自然辛结构；辛几何中的一个有力思想是，
任何辛流形 $(M,\omega)$ 都是\emph{几乎 Kähler} 的。
特别地，Kähler 流形的复结构与 Kähler 度量，可以换成几乎复结构及与 $\omega$ 关联的 Riemann 度量。

辛流形上的几乎复结构可满足若干条件。下面的定义正式给出其中两个条件。

\SourcePage{208}
\begin{definition}\label{def:tamed-compatible-complex-structure}
设 $(M,\omega)$ 为辛流形，$J$ 为 $M$ 上的几乎复结构。若
\[
\omega(X,JX)>0,\qquad
p\in M,\quad0\ne X\in T_pM,
\]
则称 $J$ \emph{受 $\omega$ 驯服}。
此外，若 $J$ 受 $\omega$ 驯服，且对 $M$ 上所有向量场 $X,Y$ 都有
\[
\omega(JX,JY)=\omega(X,Y),
\]
则称 $J$ \emph{与 $\omega$ 相容}。
\end{definition}
\index{驯服的几乎复结构|hyperpage}\index{相容的几乎复结构|hyperpage}
相容性等价于 $g(X,Y)=\omega(X,JY)$ 是 $M$ 上的 Riemann 度量。
满足上述任一条件的 $J$ 总存在，而且彼此之间都可通过形变连接。
例如，这意味着：在辛流形上，可用任意受 $\omega$ 驯服的 $J$，
赋予切丛 $TM$ 复结构，从而定义它的 Chern 类 $c_i(TM)$。
若 $J$ 受 $\omega$ 驯服，则
\[
\mu_J(X,Y)=\frac12\bigl(\omega(X,JY)+\omega(Y,JX)\bigr)
\]
定义一个度量；当 $J$ 与 $\omega$ 相容时，它就是上述 $g$。

使用\IndexLocation{indexloc:0048:208}受驯服的几乎复结构，其好处是更容易得到“一般”情形。
另一方面，若施加更强的相容性条件，往往能作出更精确的陈述。

这里的主要处理沿用 \SourceCite{Ruan2} 开创的方法，也讨论 \SourceCite{RT1,RT2} 中对它的推广；
这些工作为\emph{半正}辛流形发展了 Gromov--Witten 不变量理论。
半正是指：对每个可由映射 $f:S^2\to M$ 代表的同调类 $\beta\in H_2(M,\mathbb Z)$，
不会同时出现
\[
\int_\beta\omega>0,\qquad
3-\tfrac12\dim_{\mathbb R}M\leq\int_\beta c_1(TM)<0.
\]
例如，每个 Calabi--Yau 流形在其自然辛结构下都是半正的，$\mathbb P^r$ 也是如此。
还要注意，每个复维数不超过 $3$ 的光滑射影簇都是半正的。
提醒读者，文献中“半正”有时采用稍有不同的含义。\index{半正辛流形|hyperpage}

接下来讨论辛流形 $M$ 上的曲线。在代数几何中，用方程定义曲线及其他子簇。
在辛几何中这一做法不适用，因此改令 $C$ 为 Riemann 曲面，考虑
$C^\infty$ 映射 $f:C\to M$，它在以下意义下“全纯”。

\begin{definition}\label{def:j-holomorphic-map}
设 $(M,\omega)$ 为辛流形，$J$ 为受 $\omega$ 驯服的几乎复结构。
$J$ 全纯映射是 Riemann 曲面 $C$ 到 $M$ 的 $C^\infty$ 映射 $f:C\to M$，
其微分 $df:TC\to TM$ 相对于 $C$ 的复结构与 $M$ 的几乎复结构，是复向量丛映射。
\end{definition}
\index{J全纯映射@$J$ 全纯映射|hyperpage}
注意，映射 $f:C\to M$ 有一个可作用于它的 $\overline\partial$ 算子，定义为
\begin{equation}\label{eq:symplectic-cauchy-riemann-operator}
\overline\partial_J=d+J\circ d\circ j_C,
\end{equation}
其中 $j_C$ 为 $C$ 的可积几乎复结构。
$J$ 全纯性等价于 Cauchy--Riemann 方程 $\overline\partial_Jf=0$；
对 $\overline\partial_J$ 的定义施加 $J$ 并使用 $J^2=-1$，即可验证。
特别地，若 $J$ 可积，本情形中的 $J$ 全纯映射就是参数化的复曲线。

\SourcePage{209}
辛 Gromov--Witten 不变量的基本思想，类似于式~\eqref{eq:gw-pointed-maps} 的做法：
希望计数曲线
\begin{equation}\label{eq:symplectic-gw-pointed-maps}
f:C\to M,\qquad f_*[C]=\beta,\quad f(p_i)\in Z_i,\quad i=1,\ldots,n,
\end{equation}
其中，像先前一样，$p_1,\ldots,p_n$ 为 $C$ 的固定点，
$\beta\in H_2(M,\mathbb Z)$ 为同调类，$Z_1,\ldots,Z_n$ 为 $M$ 上处于一般位置的圈。
在 \SourceCite{Ruan2,RT2} 的方法中，通过要求 $f:C\to M$ 对固定的一般 $J$ 为 $J$ 全纯映射，
使这一表述精确化；实际上，还需要扰动 Cauchy--Riemann 方程。
一般性的作用，是确保有关模空间具有预期维数。所得不变量依赖 $\omega$，不依赖所选 $J$。
这与代数情形一致，因为\hyperref[sec:gw-properties]{第~\ref*{sec:gw-properties} 节}将看到，\hyperref[sec:gw-algebraic-definition]{第~\ref*{sec:gw-algebraic-definition} 节}的 Gromov--Witten 不变量在复结构形变下不变。
从镜像对称看，这是合理的：这些不变量出现在 A 模型关联函数中，而后者仅依赖 Kähler 模。

辛几何方法的一个优点是：选用一般几乎复结构时，代数情形中的许多困难，
例如模空间奇异或维数不符，便消失了。
事实上，与代数情形相比，式~\eqref{eq:symplectic-gw-pointed-maps} 已更接近 Gromov--Witten 不变量的真正定义。
不过，在较新的辛 Gromov--Witten 不变量方法中，可去掉半正假设，同时构造虚拟基本类的对应物。
这时，也不再需要实际模空间达到预期维数。

为了描述这些方法，必须了解有关模空间。下面转入这一主题。

\subsection{模空间}\label{subsec:symplectic-map-moduli}
\index{J全纯映射@$J$ 全纯映射!模空间|hyperpage}\index{扰动的J全纯映射@扰动的 $J$ 全纯映射|hyperpage}
先沿用 \SourceCite{Ruan2} 的处理，介绍性论述见 \SourceCite{MS}。
这里的思想，是研究类为 $\beta\in H_2(M,\mathbb Z)$ 的 $J$ 全纯映射 $f:C\to M$。
固定 $C$ 的复结构与 $M$ 的几乎复结构 $J$。出于技术原因，需要假定 $f$ 是\emph{简单的}，
即非常值且不能分解为 $C\to C'\to M$，其中 $C\to C'$ 的次数大于 $1$。
若 $f$ 不是简单的，称它为\emph{\IndexLocation{indexloc:0653:209}多重覆盖}。令
\[
\mathcal M(C,J,\beta)=
\{f:C\to M:f\text{ 为简单的 }J\text{ 全纯映射，且 }f_*[C]=\beta\}.
\]
第一个结果是：当 $J$ 为受 $\omega$ 驯服的一般几乎复结构时，
$\mathcal M(C,J,\beta)$ 自然为实流形，其维数为
\begin{equation}\label{eq:fixed-curve-map-index}
(1-g)\dim_{\mathbb R}M+2\int_\beta c_1(TM).
\end{equation}
为了证明，将 $\mathcal M(C,J,\beta)$ 表示为两个无穷维流形之间某个映射的纤维。
对一般 $J$，这一映射具有 Fredholm 线性化，隐函数定理便表明纤维
$\mathcal M(C,J,\beta)$ 为有限维流形。
纤维维数由 Fredholm 指数给出，而 Hirzebruch--\allowbreak Riemann--\allowbreak Roch 定理给出式~\eqref{eq:fixed-curve-map-index}。
顺便指出，类似分析也用于构造 Donaldson 理论与 Seiberg--Witten 理论中的模空间。

流形 $\mathcal M(C,J,\beta)$ 对 $J$ 的依赖也很好。
更精确地说，随着 $J$ 变化，只要 $J$ 足够一般，各 $\mathcal M(C,J,\beta)$ 就彼此协边。
不过，由于 $\mathcal M(C,J,\beta)$ 通常非紧，这一结果没有想象中那么强。
\hyperref[subsec:instanton-numbers]{第~\ref*{subsec:instanton-numbers} 小节}猜想~\ref{conj:symplectic-instantons} 之后的讨论，将着重指出非紧性带来的一个困难。\index{协边|hyperpage}

\SourcePage{210}
接着介绍 Ruan 与 Tian \SourceCite{RT1} 对这些空间采用的另一种方法。
这里扰动 $J$ 全纯映射的概念；部分原因是为了处理多重覆盖。
而且，正如下文将指出的，扰动能得到紧化模空间。

回忆，$J$ 全纯意味着 Cauchy--Riemann 方程 $\overline\partial_Jf=0$。
在积空间 $C\times M$ 上，令
\[
\nu:p_1^*TC\longrightarrow p_2^*TM
\]
为共轭线性映射，其中 $p_i$，$i=1,2$，为到 $C\times M$ 第 $i$ 个因子的投影。
考虑\emph{非齐次 Cauchy--Riemann 方程}
\[
(\overline\partial_Jf)(p)=\nu(p,f(p)),\qquad p\in C.
\]
若 $f:C\to M$ 满足此方程，称它为\emph{扰动的 $J$ 全纯映射}。
假定 $J,\nu$ 都一般，考虑集合
\[
\mathcal M(C,J,\nu,\beta)
=\{f:C\to M:\overline\partial_Jf=\nu,\ f_*[C]=\beta\}.
\]
注意，不再要求 $f$ 简单。这是采用扰动的 $J$ 全纯映射的优点之一。
与齐次情形 $\nu=0$ 的关系，见 \SourceCite[第 2.10 节]{RT1}。
可证明，$\mathcal M(C,J,\nu,\beta)$ 自然为流形，维数为式~\eqref{eq:fixed-curve-map-index}。

到目前为止，曲线 $C$ 的复结构固定。\SourceCite{RT1} 中这一情形产生的不变量，
可用来定义包含所有亏格的拓扑量子场论（\hyperref[sec:topological-qft]{附录 B.6}）。
若允许复结构变化，用物理语言说就是拓扑理论\emph{与引力耦合}，
便得到先前所研究的模空间 $\overline M_{g,n}(X,\beta)$ 的辛几何对应物。
由于需要空间 $\overline M_{g,n}$，必须假定 $2g+n\geq3$。
不过，$\overline M_{g,n}$ 的奇点要求谨慎处理。
在 \SourceCite{RT2} 中，Ruan 与 Tian 定义空间 $\mathcal M^\mu_{g,n}(J,\nu,\beta)_I$，
粗略地说，它由满足 $f_*[C]=\beta$ 的所有扰动的 $J$ 全纯映射
\[
f:(C,p_1,\ldots,p_n)\longrightarrow M
\]
组成。这里 $C$ 光滑，亏格为 $g$；上标 $\mu$ 提醒我们，实际上是在
$\overline M_{g,n}$ 的有限覆盖 $\overline M^\mu_{g,n}$ 上工作。
通常通过固定一个级结构来实现这一覆盖，它是获得泛曲线所必需的。
下标 $I$ 则表示限制到 $\overline M^\mu_{g,n}$ 中由自同构群平凡的带 $n$ 个标记曲线组成的子集。

与先前一样，$\mathcal M^\mu_{g,n}(J,\nu,\beta)_I$ 自然为流形，维数为
\begin{equation}\label{eq:varying-curve-map-index}
2(1-g)\left(\tfrac12\dim_{\mathbb R}M-3\right)
+2\int_\beta c_1(TM)+2n.
\end{equation}
当 $M$ 为射影簇 $X$ 时，有 $\int_\beta c_1(TX)=-\int_\beta\omega_X$。
因此，这一维数恰为 $\overline M_{g,n}(X,\beta)$ 的预期实维数，
可与式~\eqref{eq:stable-map-expected-dimension} 比较。

这样，从扰动的 $J$ 全纯映射构造了两类空间：
$\mathcal M(C,J,\nu,\beta)$ 与 $\mathcal M^\mu_{g,n}(J,\nu,\beta)_I$。
各空间都有称为 \emph{Gromov--Uhlenbeck 紧化}的自然紧化。
由于 $(M,\omega)$ 半正，可证明这些紧化通过增加有限多个实余维至少 $2$ 的层得到。
下面将看到，这两个空间给出两种不同但密切相关的 Gromov--Witten 不变量。\index{Gromov--Uhlenbeck紧化@Gromov--Uhlenbeck 紧化|hyperpage}

\SourcePage{211}
现在介绍 \SourceCite{LTi3} 提出的第三种模空间方法。
这一方法使用与 $\omega$ 相容的几乎复结构 $J$，并引入\emph{稳定 $C^\ell$ 映射}。
它们是稳定映射 $f:C\to M$ 的 $C^\ell$ 版本。
曲线 $C$ 与稳定映射定义中的完全相同，标记点仍为光滑点。
区别在于，$f$ 不再是代数态射，而满足以下条件：
对 $C$ 的每个分支 $C_i$，$f$ 与正规化映射 $\widetilde C_i\to C_i$ 的复合是 $C^\ell$ 的。
此外，不仅对 $f$ 为常值的分支施加稳定性条件，而且对满足 $f_*[C_i]=0$ 的分支都施加这一条件。
$C$ 的双全纯同构给出 $C^\ell$ 映射之间自然的等价概念。

在这一情形中，还有相应的 $J$ 全纯稳定映射概念：
要求各映射 $\widetilde C_i\to C_i\to M$ 为 $J$ 全纯映射。

若固定亏格 $g$、标记点数目 $n$ 与同调类 $\beta=f_*[C]$，所得稳定 $C^\ell$ 映射的等价类集合记为
$\overline{\mathcal F}^{\ell}_\beta(M,g,n)$。
这一空间有自然拓扑，是要研究的模空间。\SourceCite{Ruan3,Siebert1} 的方法中也使用类似空间。

\subsection{辛 Gromov--Witten 不变量}\label{subsec:symplectic-gw-invariants}
\index{伪流形|hyperpage}\index{虚邻域|hyperpage}\index{有向Euler类@有向 Euler 类|hyperpage}
空间 $\mathcal M(C,J,\nu,\beta)$ 与 $\mathcal M^\mu_{g,n}(J,\nu,\beta)_I$ 有良好的紧化，
意味着它们几乎具有基本类。问题在于，边界 $\partial$ 是否有邻域 $U$，使 $\partial$ 为 $U$ 的形变收缩，
见 \SourceCite[注 4.3]{RT1}。若有正确的基本类，便可用类似式~\eqref{eq:gw-fundamental-class-definition}
的公式定义 Gromov--Witten 不变量。

由于没有一个良好的基本类，有两种处理方法：
\begin{itemize}
\item 用横截性论证，将定义化为有向交中的点数。这是 \SourceCite{RT1,RT2} 的方法。
\item 如\hyperref[sec:gw-algebraic-definition]{第~\ref*{sec:gw-algebraic-definition} 节}的虚拟基本类讨论一样，在更大的空间中工作。
在 \SourceCite{Ruan3} 中，Ruan 使用\emph{虚邻域}，通过将模空间嵌入一个 V 流形得到它。
\SourceCite{Siebert1} 有相关方法；\SourceCite{LTi3} 中也以类似方式使用上面定义的
空间 $\overline{\mathcal F}^{\ell}_\beta(M,g,n)$。
\end{itemize}
第二种方法更有力，因为它适用于所有辛流形，而非仅限半正辛流形。
不过，这里的讨论主要采用第一种方法，同时简要提及 \SourceCite{LTi3} 的一些进一步思想，
作为更一般技术的代表。

先从 $\mathcal M(C,J,\nu,\beta)$ 给出的不变量开始。目标是定义\emph{混合 Gromov--Witten 不变量}
\begin{equation}\label{eq:mixed-gw-invariant}
\Phi_{g,\beta}(\alpha_1,\ldots,\alpha_k\mid\beta_1,\ldots,\beta_l),
\end{equation}
其中 $\alpha_i,\beta_j\in H^*(M,\mathbb Z)$ 为满足
\begin{equation}\label{eq:mixed-gw-degree-condition}
\sum_{i=1}^k\deg\alpha_i+\sum_{j=1}^l\deg\beta_j
=(1-g)\dim_{\mathbb R}M+2\int_\beta c_1(TM)+2l
\end{equation}
的上同调类。在 \SourceCite{RT1} 中，Ruan 与 Tian 使用同调类，因此这里的 $\alpha_i,\beta_j$
与他们所用的类互为对偶。粗略思想如下：固定 $C$ 上的点 $p_1,\ldots,p_k$。
混合不变量~\eqref{eq:mixed-gw-invariant} 应是 $\mathcal M(C,J,\nu,\beta)$ 中满足
$f(p_i)$ 位于 $\alpha_i$ 的对偶圈、$f(C)$ 遇到 $\beta_j$ 的对偶圈的映射 $f:C\to M$ 的数目。
\index{混合Gromov--Witten不变量@混合 Gromov--Witten 不变量|hyperpage}

\SourcePage{212}
为了精确定义，用 $M$ 的\emph{伪流形}表示 $\alpha_i,\beta_j$ 的对偶圈。
这里，伪流形指连续映射 $F:Y\to M$，其中 $Y$ 分层，每个较低层在相邻较高层中至少有余维 $2$，
而 $F$ 在每一层上光滑。若 $(F_i,Y_i)$ 与 $(G_j,Z_j)$ 分别表示 $\alpha_i,\beta_j$ 的对偶圈，则
\[
F=\prod_{i=1}^kF_i\times\prod_{j=1}^lG_j:
\prod_{i=1}^kY_i\times\prod_{j=1}^lZ_j\longrightarrow M^{k+l}
\]
为伪流形，它表示与上同调类
$\alpha_1\otimes\cdots\otimes\alpha_k\otimes\beta_1\otimes\cdots\otimes\beta_l
\in H^*(M^{k+l},\mathbb Z)$ 对偶的圈。

另一方面，有固定点 $p_1,\ldots,p_k\in C$。它们给出\IndexLocation{indexloc:0276:212}赋值映射
\[
e:\mathcal M(C,J,\nu,\beta)\times C^l\longrightarrow M^{k+l},
\]
将 $(f,q_1,\ldots,q_l)$ 映为
$(f(p_1),\ldots,f(p_k),f(q_1),\ldots,f(q_l))\in M^{k+l}$。

于是，有两个到 $M^{k+l}$ 的映射。对一般 $J,\nu$，可证明两者表示横截相交的圈；
特别地，交点出现在光滑部分。次数条件~\eqref{eq:mixed-gw-degree-condition} 还表明，
$\operatorname{Im}(F)$ 与 $\operatorname{Im}(e)$ 表示互补维数的圈。
故按原像计数的相交集只有有限多个点。
由于所有对象都有自然定向，每个相交点 $P$ 都有由定向确定的符号 $\epsilon(P)=\pm1$。

\begin{definition}\label{def:mixed-gw-invariant}
设 $(M,\omega)$ 为半正辛流形，$\beta\in H_2(M,\mathbb Z)$ 为同调类。
给定满足式~\eqref{eq:mixed-gw-degree-condition} 的
$\alpha_i,\beta_j\in H^*(M,\mathbb Z)$，混合 Gromov--Witten 不变量定义为
\[
\Phi_{g,\beta}(\alpha_1,\ldots,\alpha_k\mid\beta_1,\ldots,\beta_l)
=\sum_{P\in F\times_{M^{k+l}}e}\epsilon(P).
\]
\end{definition}

Ruan 与 Tian \SourceCite{RT1} 证明，
$\Phi_{g,\beta}(\alpha_1,\ldots,\alpha_k\mid\beta_1,\ldots,\beta_l)$
与所选 $J,\nu$、点 $p_1,\ldots,p_k\in C$、$C$ 的复结构，
以及表示 $\alpha_i,\beta_j$ 的伪流形无关。他们还讨论了\IndexLocation{indexloc:0546:212}混合不变量的许多性质与应用。

可将这一构造与\hyperref[sec:gw-algebraic-definition]{第~\ref*{sec:gw-algebraic-definition} 节}的代数不变量联系起来。
设 $(M,\omega)$ 是带 Kähler 形式的光滑射影簇。令
\[
\pi:\overline M_{g,k+l}(M,\beta)\longrightarrow\overline M_{g,k}
\]
将稳定映射忘掉映射本身及最后 $l$ 个标记点，并收缩因此出现的不稳定分支。
这是式~\eqref{eq:stable-map-evaluation-forgetting} 中 $\pi_2$ 的推广；$\pi_2$ 对应 $l=0$。
令 $[pt]$ 表示 $\overline M_{g,k}$ 上一点的上同调类，
它代表带固定标记点 $(p_1,\ldots,p_k)$ 的固定曲线 $C$。
则不变量 $\Phi_{g,\beta}(\alpha_1,\ldots,\alpha_k\mid\beta_1,\ldots,\beta_l)$ 在思想上就是
\[
\int_{[\overline M_{g,k+l}(M,\beta)]^{\mathrm{virt}}}
\left(\prod_{i=1}^ke_i^*(\alpha_i)\right)\cup
\left(\prod_{j=1}^le_{k+j}^*(\beta_j)\right)\cup\pi^*[pt].
\]

不过，我们更关心\hyperref[sec:gw-algebraic-definition]{第~\ref*{sec:gw-algebraic-definition} 节}所讨论类型的 Gromov--Witten 不变量。
因此，转而研究由 $\mathcal M^\mu_{g,n}(J,\nu,\beta)_I$ 构造的不变量。
它们是\emph{辛 Gromov--Witten 不变量}
\begin{equation}\label{eq:symplectic-gw-invariant}
\Psi_{g,n,\beta}(\gamma;\alpha_1,\ldots,\alpha_n),
\end{equation}
\SourcePage{213}
为满足
\[
\deg\gamma+\sum_{i=1}^n\deg\alpha_i
=2(1-g)\left(\tfrac12\dim_{\mathbb R}M-3\right)
+2\int_\beta c_1(TM)+2n
\]
的上同调类 $\gamma\in H^*(\overline M_{g,n},\mathbb Q)$、
$\alpha_i\in H^*(M,\mathbb Z)$ 定义。
朴素思想是：若 $Z$ 为 $\overline M_{g,n}$ 中与 $\gamma$ 对偶的圈，
$Z_i$ 为 $M$ 中与 $\alpha_i$ 对偶的圈，则
$\Psi_{g,n,\beta}(\gamma;\alpha_1,\ldots,\alpha_n)$ 应是满足
$(C,p_1,\ldots,p_n)\in Z\subset\overline M_{g,n}$、
$f_*[C]=\beta$、$f(p_i)\in Z_i$ 的扰动 $J$ 全纯映射
$f:(C,p_1,\ldots,p_n)\to M$ 的数目。

式~\eqref{eq:symplectic-gw-invariant} 的精确定义类似混合情形，但细节更复杂：
并非 $H_*(\overline M_{g,n},\mathbb Z)$ 中所有类都由所需的伪流形表示，
而使用 $\mathcal M^\mu_{g,n}(J,\nu,\beta)_I$ 意味着是在有限覆盖
$\overline M^\mu_{g,n}\to\overline M_{g,n}$ 上工作。
这样先定义依赖有限覆盖选择的不变量
$\Psi^\mu_{g,n,\beta}(\gamma;\alpha_1,\ldots,\alpha_n)$。
再除以覆盖次数，得到与 $\mu$ 无关的
$\Psi_{g,n,\beta}(\gamma;\alpha_1,\ldots,\alpha_n)$。
这里不展开细节，完整定义见 \SourceCite{RT2}。

注意，固定 $\alpha_1,\ldots,\alpha_n$ 后，对所有 $\gamma$ 知道
$\Psi_{g,n,\beta}(\gamma;\alpha_1,\ldots,\alpha_n)$，
就唯一确定类
$\Psi_{g,n,\beta}(\alpha_1,\ldots,\alpha_n)\in H^*(\overline M_{g,n},\mathbb Q)$，满足
\[
\int_{\overline M_{g,n}}\gamma\cup
\Psi_{g,n,\beta}(\alpha_1,\ldots,\alpha_n)
=\Psi_{g,n,\beta}(\gamma;\alpha_1,\ldots,\alpha_n).
\]
称 $\Psi_{g,n,\beta}(\alpha_1,\ldots,\alpha_n)$ 为\emph{辛 Gromov--Witten 类}。
与代数情形相同，辛类可看成映射
\begin{equation}\label{eq:symplectic-gw-class-map}
\Psi_{g,n,\beta}:H^*(M)^{\otimes n}\longrightarrow H^*(\overline M_{g,n},\mathbb Q).
\end{equation}
后面将看到，代数 Gromov--Witten 类与辛 Gromov--Witten 类具有相似性质。

混合 Gromov--Witten 不变量与辛 Gromov--Witten 不变量之间有一些很好的关系。
例如，混合不变量固定了复结构，因而下式并不令人意外：
\[
\Phi_{g,\beta}(\alpha_1,\ldots,\alpha_n\mid)
=\Psi_{g,n,\beta}([pt];\alpha_1,\ldots,\alpha_n).
\]
为了陈述一般关系，令 $K_{k,l}$ 为 $\overline M_{g,k+l}$ 中以下集合的闭包所给出的圈：
\[
\{(C,p_1,\ldots,p_{k+l})\in\overline M_{g,k+l}:
(C,p_1,\ldots,p_k)\text{ 为 }\overline M_{g,k}\text{ 中的固定点}\}.
\]
则可证明
\[
\Phi_{g,\beta}(\alpha_1,\ldots,\alpha_k\mid\beta_1,\ldots,\beta_l)
=\Psi_{g,k+l,\beta}([K_{k,l}];
\alpha_1,\ldots,\alpha_k,\beta_1,\ldots,\beta_l),
\]
这意味着所有混合不变量都是辛不变量的特殊情形。
一个尤其相关的情况是 $g=0$。由于 $\overline M_{0,3}$ 只有一点，
故
\begin{equation}\label{eq:genus-zero-mixed-symplectic}
\Phi_{0,\beta}(\alpha_1,\alpha_2,\alpha_3\mid\alpha_4,\ldots,\alpha_n)
=\Psi_{0,n,\beta}(1;\alpha_1,\ldots,\alpha_n).
\end{equation}
对此公式的很好说明见 \SourceCite{Voisin3}。

本节最后简要提及 \SourceCite{LTi3} 方法中的一些思想。
令 $\overline{\mathcal F}^{\ell}_\beta(M,g,n)$ 为\hyperref[subsec:symplectic-map-moduli]{第~\ref*{subsec:symplectic-map-moduli} 小节}末尾描述的模空间。
\SourcePage{214}
在其上如下定义广义丛 $E$。给定 $C$ 光滑的稳定 $C^\ell$ 映射 $f:C\to M$，
令 $\Lambda_f^{0,1}$ 为
\[
\operatorname{Hom}(TC,f^*TM)
\]
中满足 $\nu\circ j_C=-J\circ\nu$ 的连续截面 $\nu$ 的空间。
这些截面应看作 $C$ 上取值于 $f^*(TM)$ 的 $(0,1)$ 形式。
注意，Cauchy--Riemann 算子 $\overline\partial_Jf$ 属于 $\Lambda_f^{0,1}$，
这可用式~\eqref{eq:symplectic-cauchy-riemann-operator}、$J^2=-1$ 与 $j_C^2=-1$ 立即算出。
当 $C$ 不光滑时，$\Lambda_f^{0,1}$ 的定义需稍作修改，细节见 \SourceCite{LTi3}；
修改与定义 $\overline{\mathcal F}^{\ell}_\beta(M,g,n)$ 的等价关系相容。
结果是，各 $\Lambda_f^{0,1}$ 成为 $\overline{\mathcal F}^{\ell}_\beta(M,g,n)$ 上广义丛 $E$ 的纤维。

由构造，$\overline\partial_J$ 诱导 $E$ 的全局截面。
于是，我们得到无穷维空间上一个无穷维丛的截面。
思想是寻找一种对象，使它发挥有限维情形中丛的 Euler 类的作用。
\SourceCite{LTi3} 证明，当 $\ell\geq2$ 时，这一截面赋予 $E$ \emph{广义 Fredholm 轨形丛}结构，
这一概念由 \SourceCite{LTi3} 引入。主要结果是：这些数据确定一个\emph{有向 Euler 类}，
它是丛的底空间上的同调类。
这个类支于 $E$ 的截面 $\overline\partial_J$ 的一个精心选取的扰动的零点轨迹。
所研究情形中，该同调类的维数恰好是预期维数~\eqref{eq:varying-curve-map-index}。
这一有向 Euler 类充当虚拟基本类，容易按熟悉的方法定义辛 Gromov--Witten 不变量。

\SourceCite{LTi3} 的 Gromov--Witten 不变量推广了 \SourceCite{RT2} 的不变量，
后者只在 $M$ 半正时定义。注意，不必假定 $2g+n\geq3$。
\SourceCite{Siebert1} 也有相关构造。
人们相信各种辛 Gromov--Witten 不变量彼此一致，但截至本书写作时，尚无已写出的证明。
下一节将讨论：各种辛不变量与\hyperref[sec:gw-algebraic-definition]{第~\ref*{sec:gw-algebraic-definition} 节}的代数 Gromov--Witten 不变量，在多大程度上相同。

\section{Gromov--Witten 类的性质}\label{sec:gw-properties}
接下来的任务，是研究 Gromov--Witten 类，以及前两节各种定义之间的关系。
先描述所有定义共有的性质。

\subsection{Gromov--Witten 类的公理}\label{subsec:gw-axioms}
\index{Gromov--Witten类@Gromov--Witten 类!公理|hyperpage}\index{分裂公理|hyperpage}\index{约化公理|hyperpage}\index{除子公理|hyperpage}\index{基本类公理|hyperpage}\index{点映射公理|hyperpage}\index{形变公理|hyperpage}\index{复合律|hyperpage}
在 \SourceCite{KoM1} 中，Kontsevich 与 Manin 为 Gromov--Witten 类提出了一组公理。
现在已知，\hyperref[sec:gw-algebraic-definition]{第~\ref*{sec:gw-algebraic-definition} 节}的代数类 $I_{g,n,\beta}$ 与\hyperref[sec:gw-symplectic-definition]{第~\ref*{sec:gw-symplectic-definition} 节}的辛类 $\Psi_{g,n,\beta}$ 都满足这些公理。
代数情形的证明见 \SourceCite{BM,LTi2,BF,Behrend}；
辛情形的证明见 \SourceCite{RT1,RT2,LTi3}。
还将加入一条形变公理；它在代数情形中由 \SourceCite{LTi2} 证明，
在辛情形中则由辛不变性得到。

为简便起见，为 $I_{g,n,\beta}$ 陈述公理。
回忆\hyperref[sec:gw-algebraic-definition]{第~\ref*{sec:gw-algebraic-definition} 节}：若 $X$ 为光滑射影簇，则对 $g\geq0$、$n\geq0$、$n+2g\geq3$，
Gromov--Witten 类为映射
\[
I_{g,n,\beta}:H^*(X,\mathbb Q)^{\otimes n}
\longrightarrow H^*(\overline M_{g,n},\mathbb Q).
\]
\SourcePage{215}
此外，对 $g,n\geq0$，Gromov--Witten 不变量为映射
\[
\langle I_{g,n,\beta}\rangle:H^*(X,\mathbb Q)^{\otimes n}\longrightarrow\mathbb Q.
\]
当 $n+2g\geq3$ 时，它们满足关系
\[
\langle I_{g,n,\beta}\rangle(\alpha_1,\ldots,\alpha_n)
=\int_{\overline M_{g,n}}I_{g,n,\beta}(\alpha_1,\ldots,\alpha_n).
\]

理解下面的公理时，应回忆式~\eqref{eq:gw-pointed-maps} 对 Gromov--Witten 类的朴素解释：
若 $Z_i$ 为 $X$ 中与 $\alpha_i$ 对偶的圈，
则 $I_{g,n,\beta}(\alpha_1,\ldots,\alpha_n)$ 应是由那些
$(C,p_1,\ldots,p_n)\in\overline M_{g,n}$ 给出的上同调类，其中可以找到 $f$ 满足
\begin{equation}\label{eq:gw-axiom-naive-stable-map}
f:(C,p_1,\ldots,p_n)\longrightarrow X\text{ 为稳定映射},\quad
f(p_i)\in Z_i,\quad f_*[C]=\beta.
\end{equation}
带着这一想法，进入各条公理。

\begin{axiom*}{线性}
$I_{g,n,\beta}$ 对每个变元都是线性的。
\end{axiom*}
朴素地说，这是因为圈的和就是它们的并，计入各自的重数。

\begin{axiom*}{有效性}
对光滑射影簇 $X$，若 $\beta$ 不是有效类，则 $I_{g,n,\beta}=0$。
\end{axiom*}
这是合理的，因为只要 $f:C\to X$ 全纯，$f_*[C]$ 就有效。
对辛流形 $(M,\omega)$，这一公理换成以下观察：只要 $\int_\beta\omega<0$，不变量就消失。

\begin{axiom*}{次数}
若 $\alpha_1,\ldots,\alpha_n$ 为齐次上同调类，则
$I_{g,n,\beta}(\alpha_1,\ldots,\alpha_n)\in H^*(\overline M_{g,n},\mathbb Q)$ 的次数为
\[
2(g-1)\dim X+2\int_\beta\omega_X+\sum_{i=1}^n\deg\alpha_i.
\]
\end{axiom*}
当模空间 $\overline M_{g,n}(X,\beta)$ 具有预期维数~\eqref{eq:stable-map-expected-dimension}，
并有良好的基本类时，这由式~\eqref{eq:gw-gysin-definition} 得到。
一般地，由于\IndexLocation{indexloc:1415:215}虚拟基本类具有预期次数，次数公理由定义~\ref{def:gw-algebraic-invariants} 得到。

次数公理意味着，$I_{g,n,\beta}(\alpha_1,\ldots,\alpha_n)$ 为最高次数的类，当且仅当
\begin{equation}\label{eq:gw-top-degree-condition}
\sum_{i=1}^n\deg\alpha_i
=2(1-g)\dim X-2\int_\beta\omega_X+2(3g-3+n).
\end{equation}
因此，考虑 Gromov--Witten 不变量 $\langle I_{g,n,\beta}\rangle(\alpha_1,\ldots,\alpha_n)$ 时，
总可假定式~\eqref{eq:gw-top-degree-condition} 成立，否则不变量为零。
由定义~\ref{def:gw-algebraic-invariants}，这对所有 $n,g\geq0$ 都有效。
文献中，有时用“代数次数” $\tfrac12\deg\alpha_i$ 代替 $\deg\alpha_i$；
如此，余维 $i$ 的子簇 $Z$ 的对偶类 $[Z]$ 具有次数 $i$。
处理偶次数上同调时，这一约定最为有用。

\begin{axiom*}{等变性}
对称群 $S_n$ 自然作用于 $H^*(X,\mathbb Q)^{\otimes n}$ 与 $H^*(\overline M_{g,n},\mathbb Q)$。
在前者上采用分次置换作用，在后者上通过置换 $(C,p_1,\ldots,p_n)$ 的标记点 $p_i$ 作用。
映射
\[
I_{g,n,\beta}:H^*(X,\mathbb Q)^{\otimes n}
\longrightarrow H^*(\overline M_{g,n},\mathbb Q)
\]
是 $S_n$ 等变的。
\end{axiom*}
这一性质背后的直观理由，应由式~\eqref{eq:gw-axiom-naive-stable-map} 看出。
对 Gromov--Witten 不变量，等变性意味着
\[
\begin{aligned}
&\langle I_{g,n,\beta}\rangle(\alpha_1,\ldots,\alpha_i,\alpha_{i+1},\ldots,\alpha_n)\\
&\quad=(-1)^{\deg\alpha_i\deg\alpha_{i+1}}
\langle I_{g,n,\beta}\rangle(\alpha_1,\ldots,\alpha_{i+1},\alpha_i,\ldots,\alpha_n).
\end{aligned}
\]

\SourcePage{216}
\begin{axiom*}{基本类}
若 $n+2g\geq4$，忘掉最后一点得到自然映射
$\pi_n:\overline M_{g,n}\to\overline M_{g,n-1}$。
若 $[X]\in H^0(X,\mathbb Q)$ 为 $X$ 的基本类的对偶，即单位类，则
\[
I_{g,n,\beta}(\alpha_1,\ldots,\alpha_{n-1},[X])
=\pi_n^*I_{g,n-1,\beta}(\alpha_1,\ldots,\alpha_{n-1}).
\]
\end{axiom*}
用式~\eqref{eq:gw-axiom-naive-stable-map} 理解，这是合理的：
$f(p_n)\in X$ 对 $p_n$ 不施加条件。
对 Gromov--Witten 不变量，基本类公理蕴含
\begin{equation}\label{eq:gw-fundamental-class-vanishing}
\langle I_{g,n,\beta}\rangle(\alpha_1,\ldots,\alpha_{n-1},[X])=0.
\end{equation}
因为这一不变量只有在
$I_{g,n,\beta}(\alpha_1,\ldots,\alpha_{n-1},[X])$ 为最高次数类时才可能非零。
但 $\overline M_{g,n-1}$ 的维数更小，故相应次数的
$I_{g,n-1,\beta}(\alpha_1,\ldots,\alpha_{n-1})$ 为零。
注意，只要稳定映射的忘却映射 $\pi_n$ 有定义，式~\eqref{eq:gw-fundamental-class-vanishing} 就成立。
\hyperref[subsec:computing-virtual-class]{第~\ref*{subsec:computing-virtual-class} 小节}末尾的讨论表明，$n+2g\geq4$，或者 $\beta\ne0$ 且 $n\geq1$，都保证这一点。

\begin{axiom*}{除子}
若 $n+2g\geq4$，取与上一公理相同的
$\pi_n:\overline M_{g,n}\to\overline M_{g,n-1}$。若 $\alpha_n\in H^2(X,\mathbb Q)$，则
\[
\pi_{n!}I_{g,n,\beta}(\alpha_1,\ldots,\alpha_{n-1},\alpha_n)
=\left(\int_\beta\alpha_n\right)
I_{g,n-1,\beta}(\alpha_1,\ldots,\alpha_{n-1}).
\]
\end{axiom*}
用式~\eqref{eq:gw-axiom-naive-stable-map} 看其理由：
令稳定映射 $f:(C,p_1,\ldots,p_n)\to X$ 满足 $f_*[C]=\beta$，
且 $f(p_i)\in Z_i$，$i=1,\ldots,n-1$。
则 $f(p_n)$ 必须位于 $f(C)\cap Z_n$，其相交数为 $\beta\cdot Z_n$，
即 $p_n$ 按重数计有 $\int_\beta\alpha_n$ 种选择。由此得到上述公式。

对 Gromov--Witten 不变量，若 $\alpha_n\in H^2(X,\mathbb Q)$ 且
式~\eqref{eq:gw-top-degree-condition} 成立，则
\[
\langle I_{g,n,\beta}\rangle(\alpha_1,\ldots,\alpha_{n-1},\alpha_n)
=\left(\int_\beta\alpha_n\right)
\langle I_{g,n-1,\beta}\rangle(\alpha_1,\ldots,\alpha_{n-1}).
\]
如基本类公理，只要稳定映射的 $\pi_n$ 有定义，这一等式就成立；
由\hyperref[subsec:computing-virtual-class]{第~\ref*{subsec:computing-virtual-class} 小节}，$n+2g\geq4$，或者 $\beta\ne0$ 且 $n\geq1$，都足够。

\begin{axiom*}{点映射}
当 $\beta=0$ 时，若 $g=0$ 且各 $\alpha_i$ 为齐次上同调类，则
\[
I_{0,n,0}(\alpha_1,\ldots,\alpha_n)=
\begin{cases}
\displaystyle\left(\int_X\alpha_1\cup\cdots\cup\alpha_n\right)
[\overline M_{0,n}],&\displaystyle\sum_{i=1}^n\deg\alpha_i=2\dim X,\\
0,&\text{其他情形},
\end{cases}
\]
其中 $[\overline M_{0,n}]$ 表示其基本类的对偶，即单位类。
\end{axiom*}
为了看清这一点，注意满足 $f_*[C]=0$ 的映射必为常值。
因此，式~\eqref{eq:gw-axiom-naive-stable-map} 化为 $f(C)\in Z_1\cap\cdots\cap Z_n$。
当 $\sum_{i=1}^n\deg\alpha_i=2\dim X$ 时，由次数公理，该类的次数为 $0$，
从而得到上述公式；本情形中的其他类都为零。

这对 Gromov--Witten 不变量 $\langle I_{0,n,0}\rangle$ 有一个很好的推论。
将条件 $\sum_{i=1}^n\deg\alpha_i=2\dim X$ 与式~\eqref{eq:gw-top-degree-condition} 比较，
可见 $n=3$ 很特殊。更精确地，
\begin{equation}\label{eq:gw-degree-zero-three-point}
\langle I_{0,n,0}\rangle(\alpha_1,\ldots,\alpha_n)
=\begin{cases}
\displaystyle\int_X\alpha_1\cup\alpha_2\cup\alpha_3,&n=3,\\
0,&\text{其他情形}.
\end{cases}
\end{equation}
遗憾的是，这种推理仅适用于 $g=0$。
问题在于，$\overline M_{g,n}(X,0)=\overline M_{g,n}\times X$，
只有 $g=0$ 时才具有预期维数。
对所有 $g$ 都有 $I_{g,n,0}$ 的公式，但它们可能相当复杂。
后面的式~\eqref{eq:genus-one-constant} 将给出 $g=1$ 的情形；一般情形的困难见 \SourceCite[第 2.2.5 节]{KoM1}。

\SourcePage{217}
\begin{axiom*}{分裂}
设 $g=g_1+g_2$、$n=n_1+n_2$，且 $n_i+2g_i\geq2$。
对亏格分别为 $g_1,g_2$ 的稳定曲线
$(C_1,p_1,\ldots,p_{n_1+1})$ 与 $(C_2,q_1,\ldots,q_{n_2+1})$，
将 $p_{n_1+1}$ 与 $q_1$ 识别，把 $C_1\cup C_2$ 粘合为曲线 $C$，
得到亏格 $g$ 的稳定曲线 $(C,p_1,\ldots,p_{n_1},q_2,\ldots,q_{n_2+1})$。
因此，有映射
\begin{equation}\label{eq:gw-splitting-gluing-map}
\varphi:\overline M_{g_1,n_1+1}\times\overline M_{g_2,n_2+1}
\longrightarrow\overline M_{g,n}.
\end{equation}
分裂公理断言：
\[
\begin{aligned}
\varphi^*I_{g,n,\beta}(\alpha_1,\ldots,\alpha_n)
={}&\sum_{\beta_1+\beta_2=\beta}\sum_{i,j}g^{ij}\\
&\quad\cdot I_{g_1,n_1+1,\beta_1}(\alpha_1,\ldots,\alpha_{n_1},T_i)\\
&\quad\otimes I_{g_2,n_2+1,\beta_2}(T_j,\alpha_{n_1+1},\ldots,\alpha_n).
\end{aligned}
\]
其中 $T_i$ 为 $H^*(X,\mathbb Q)$ 的齐次基，$(g^{ij})$ 是矩阵 $(g_{ij})$ 的逆，
$g_{ij}=\int_XT_i\cup T_j$。等价地，对角线在 $H^*(X\times X,\mathbb Q)$ 中的上同调类为
$\sum_{i,j}g^{ij}T_i\otimes T_j$。由有效性公理，上述求和只有有限项。
\end{axiom*}
这一公理陈述起来稍复杂。为看清公式的来源，首先注意：
式~\eqref{eq:gw-axiom-naive-stable-map} 中映射在 $\varphi$ 下的逆像，由映射
\[
f:(C_1\cup C_2,p_1,\ldots,p_{n_1},p,q,p_{n_1+1},\ldots,p_n)\longrightarrow X
\]
组成，满足 $f_*[C_1]+f_*[C_2]=\beta$、$f(p_i)\in Z_i$ 及 $f(p)=f(q)$。
第一个条件对应于上述公式中 $\beta_1+\beta_2=\beta$ 的分解；
最后一个条件表示 $(f(p),f(q))$ 应位于 $X\times X$ 的对角线上。
由于对角线的类为 $\sum_{i,j}g^{ij}T_i\otimes T_j$，得到上述公式。

\begin{axiom*}{约化}
将最后两个标记点粘合，得到自然映射
$\psi:\overline M_{g-1,n+2}\to\overline M_{g,n}$。
用上一公理的对角线类及 $(g^{ij})$，有
\[
\psi^*I_{g,n,\beta}(\alpha_1,\ldots,\alpha_n)
=\sum_{i,j}g^{ij}I_{g-1,n+2,\beta}(\alpha_1,\ldots,\alpha_n,T_i,T_j).
\]
\end{axiom*}
这是合理的，因为式~\eqref{eq:gw-axiom-naive-stable-map} 中映射在 $\psi$ 下的逆像，
由亏格 $g-1$ 的映射 $f:(C,p_1,\ldots,p_{n+2})\to X$ 组成，满足
$f_*[C]=\beta$、$f(p_i)\in Z_i$，$i=1,\ldots,n$，及 $f(p_{n+1})=f(p_{n+2})$。
最后一个条件要求 $(f(p_{n+1}),f(p_{n+2}))$ 位于 $X\times X$ 的对角线上。
如上，对角线类为 $\sum_{i,j}g^{ij}T_i\otimes T_j$，公式随之成立。

\begin{axiom*}{形变}
设 $f:\mathcal X\to T$ 为光滑固有映射，底空间 $T$ 连通；对 $t\in T$，置 $X_t=f^{-1}(t)$。
于是，对每个 $t\in T$ 与 $\beta_t\in H_2(X_t,\mathbb Z)$，有映射
\[
I^{X_t}_{g,n,\beta_t}:H^*(X_t,\mathbb Q)^{\otimes n}
\longrightarrow H^*(\overline M_{g,n},\mathbb Q).
\]
若 $\beta_t$ 是同调局部系统 $H_2(X_t,\mathbb Z)$ 的局部常值截面，
$\alpha_1,\ldots,\alpha_n$ 是 $H^*(X_t,\mathbb Q)$ 的局部常值截面，则
$I^{X_t}_{g,n,\beta_t}(\alpha_1,\ldots,\alpha_n)$ 不随 $t$ 改变。
\end{axiom*}
其直观理由是：Gromov--Witten 类应在复结构形变下不变。
采用 $I_{g,n,\beta}$ 的辛几何定义时，这一点很清楚。

这些公理中，分裂公理与\IndexLocation{indexloc:0506:217}约化公理也许最重要。
\hyperref[chap:quantum-cohomology]{第~\ref*{chap:quantum-cohomology} 章}将看到，它们与量子上同调的结合律证明密切相关。
在 Ruan 与 Tian 的论文中，这两条公理合称\emph{复合律}。

除上述公理外，\SourceCite{KoM1} 还包含一条\emph{\IndexLocation{indexloc:0504:217}动机公理}：
Gromov--Witten 类应由
\SourcePage{218}
$X^n\times\overline M_{g,n}$ 的 Chow 环中适当次数的类诱导。
由于代数 Gromov--Witten 类用虚拟基本类
$[\overline{\mathcal M}_{g,n}(X,\beta)]^{\mathrm{virt}}$ 定义，这条公理随之成立。

应强调，上述各公理的直观解释不是证明。
不过，在辛情形中，某些“论证”确实与实际证明相当接近，
但复合律除外，见 \SourceCite{RT1,RT2}。
另一方面，代数情形的证明 \SourceCite{LTi2,Behrend} 更为精细。
如 \SourceCite{BM} 所解释，上述公理归结为虚拟基本类
$[\overline{\mathcal M}_{g,n}(X,\beta)]^{\mathrm{virt}}$ 之间五种关键的相容性。

\subsection{各种定义的等价性}\label{subsec:gw-equivalence}
\index{Gromov--Witten不变量@Gromov--Witten 不变量!各种定义的等价性|hyperpage}
\hyperref[sec:gw-algebraic-definition]{第~\ref*{sec:gw-algebraic-definition} 节}和\hyperref[sec:gw-symplectic-definition]{第~\ref*{sec:gw-symplectic-definition} 节}介绍了 Gromov--Witten 类与不变量的各种定义。
代数情形中，有两个定义，因为 \SourceCite{LTi2} 与 \SourceCite{BF,Behrend}
给出的虚拟基本类构造略有不同。如\hyperref[subsec:virtual-fundamental-class]{第~\ref*{subsec:virtual-fundamental-class} 小节}所述，二者应等价，但证明尚未写出。

辛情形中，\SourceCite{LTi3} 的不变量推广了 \SourceCite{RT2} 的不变量，
此外还有 \SourceCite{Siebert1} 给出的辛不变量。
同样，预期这些定义给出相同不变量，但尚未证明。

一个很好的进展是，代数不变量与辛不变量的一致性证明已有实质推进。
特别地，\SourceCite{LTi4} 证明了 \SourceCite{LTi2} 定义的代数不变量与
\SourceCite{LTi3} 定义的辛不变量一致；
\SourceCite{Siebert2} 证明了 \SourceCite{Behrend} 的代数不变量与
\SourceCite{Siebert1} 的辛不变量一致。
因此，似乎很可能所有定义在共同适用范围内都一致。

下面是一种已知这种一致性成立的情形。
用 $\langle I_{g,n,\beta}\rangle^{\mathrm{LRT}}$ 表示：
由 \SourceCite{LTi2} 的代数构造，或由 \SourceCite{LTi3} 的辛构造定义的 Gromov--Witten 不变量，
后者推广了 \SourceCite{RT2}。
又用 $\langle I_{g,n,\beta}\rangle^{\mathrm{BFS}}$ 表示：
由 \SourceCite{Behrend} 的代数构造，或由 \SourceCite{Siebert1} 的辛构造定义的 Gromov--Witten 不变量；
前者使用 \SourceCite{BF} 中的相对虚拟基本类。于是有以下结果。

\begin{theorem}\label{thm:homogeneous-gw-equivalence}
设 $X$ 为光滑射影簇，某个线性代数群在其上可迁地作用。
取 $\alpha_1,\ldots,\alpha_n\in H^*(X,\mathbb Q)$，则
\[
\langle I_{0,n,\beta}\rangle^{\mathrm{LRT}}(\alpha_1,\ldots,\alpha_n)
=\langle I_{0,n,\beta}\rangle^{\mathrm{BFS}}(\alpha_1,\ldots,\alpha_n).
\]
此外，当次数条件~\eqref{eq:gw-top-degree-condition} 成立、各类以一般位置的圈表示时，
这些不变量按重数计数满足式~\eqref{eq:gw-pointed-maps} 的带标记映射 $f:\mathbb P^1\to X$。
\end{theorem}
\begin{proof}
为检验等式，可以采用 $\langle I_{0,n,\beta}\rangle^{\mathrm{LRT}}$
与 $\langle I_{0,n,\beta}\rangle^{\mathrm{BFS}}$ 各自的代数定义。
\SourceCite{FP} 的引理 13 表明，计数定义给出代数 Gromov--Witten 不变量。
特别地，虚拟基本类等于基本类；采用虚拟基本类的任一种代数定义都如此。
群作用保证能找到与 $\alpha_i$ 对偶、处于一般位置的圈 $Z_i$。

在辛几何方面，Li 与 Tian 在 \SourceCite{LTi1} 中利用可迁群作用证明：
采用 $X$ 给定的复结构 $J_0$ 并取 $\nu=0$，得到相同数目。
对 \SourceCite{Siebert1} 的辛定义，也有类似结论。
\end{proof}

定理~\ref{thm:homogeneous-gw-equivalence} 断言这些不变量相等，其证明还表明：
它们确实计数相应类型的有理曲线。下一节将看到，一般 Gromov--Witten 不变量的\IndexLocation{indexloc:0534:218}枚举意义，
通常更难理解。

\SourcePage{219}
最后还应提到，除了迄今讨论的代数与辛 Gromov--Witten 不变量外，
还有一些相关不变量也值得研究。在第 \hyperref[chap:localization]{\ref*{chap:localization}}、\hyperref[chap:quantum-differential-equations]{\ref*{chap:quantum-differential-equations}} 章中，我们将遇到其中两种：
\emph{等变 Gromov--Witten 不变量}与\emph{引力关联数}。
它们在 Givental 关于镜像对称的工作中，以及相关思想和后续发展中，都起着关键作用。
\index{Gromov--Witten不变量@Gromov--Witten 不变量!等变不变量|hyperpage}
\index{引力关联数|hyperpage}

\hypertarget{source:219-batch5}{}
\section{计算 Gromov--Witten 不变量，I}\label{sec:computing-gw-i}
既然已经了解 Gromov--Witten 类与不变量的定义和性质，现在就来计算一些例子。
本节将讨论各种不同空间 $X$ 的 Gromov--Witten 不变量
$\langle I_{g,n,\beta}\rangle(\alpha_1,\ldots,\alpha_n)$。
大多数例子的亏格为 $g=0$ 或 $1$。鉴于 Gromov--Witten 不变量的直观解释~\eqref{eq:gw-pointed-maps}，
我们还将探讨它们的枚举意义。我们会发现，这些数与枚举几何之间的关系有时颇为复杂。

\hyperref[chap:quantum-cohomology]{第~\ref*{chap:quantum-cohomology} 章}将看到，量子上同调是由不变量 $\langle I_{0,n,\beta}\rangle$ 构造出来的。
该章还将介绍 \emph{Gromov--Witten 势}，它有一些非常有趣的性质，
尤其能帮助解释本节的若干例子。因此，本节实际上是讨论 Gromov--Witten 不变量的计算及其枚举意义的两节中的第一节。

\subsection{树级 Gromov--Witten 类}\label{subsec:tree-level-gw}
\index{树级 Gromov--Witten 不变量|hyperpage}\index{重构定理|hyperpage}
\index{Gromov--Witten类@Gromov--Witten 类!树级|hyperpage}
\index{重构定理|hyperpage}
当 $g=0$ 时，各类 $I_{0,n,\beta}$ 构成一个\emph{树级 Gromov--Witten 类系统}。
“树级”一词是说，亏格为零的稳定曲线是由若干 $\mathbb P^1$ 构成的树
（对应于量子场论中的树形 Feynman 图，这一术语就源自那里）。
由于始终要求 $n+2g\geq3$，树级类在 $n\geq3$ 时定义。
模空间 $\overline M_{0,n}$ 是维数为 $n-3$ 的轨形，其上同调已得到充分了解。
因此，对于 $I_{0,n,\beta}$，我们能说出许多事情。特别地，有下面这个 Kontsevich 与 Manin 的\emph{重构定理}。

\begin{theorem}\label{thm:gw-reconstruction}
设 $X$ 为光滑射影簇，且 $H^*(X,\mathbb Q)$ 由 $H^2(X,\mathbb Q)$ 生成。
再假设我们已经知道所有满足 $-\int_\beta\omega_X\leq\dim X+1$、$\deg\alpha_3=2$ 的
$\beta\in H_2(X,\mathbb Z)$ 所对应的 Gromov--Witten 不变量
$\langle I_{0,3,\beta}\rangle(\alpha_1,\alpha_2,\alpha_3)$。
那么，对所有 $\beta\in H_2(X,\mathbb Z)$，都能确定全部树级 Gromov--Witten 类
$I_{0,n,\beta}(\alpha_1,\ldots,\alpha_n)$。
\end{theorem}
\begin{proof}
先证明所有树级类都能由 Gromov--Witten 不变量重构。对 $n$ 作归纳；
由于 $\overline M_{0,3}$ 是一点，可假设 $n>3$。

为描述 $\overline M_{0,n}$ 的上同调，注意式~\eqref{eq:gw-splitting-gluing-map} 中映射 $\varphi$ 的像是
$\overline M_{0,n}$ 上的一个除子。对点 $p_1,\ldots,p_n$ 作置换，还能得到映射 $\varphi$ 的其他版本，
从而得到 $\overline M_{0,n}$ 上的其他除子。由 \SourceCite{Keel}，上同调环
$H^*(\overline M_{0,n},\mathbb Q)$ 由这些除子的类生成。
利用这一点，\SourceCite{KoM1} 证明：$H^*(\overline M_{0,n},\mathbb Q)$ 中不处于最高次数的上同调类，
由它与这些除子的交唯一确定。

现取一个类 $I_{0,n,\beta}(\alpha_1,\ldots,\alpha_n)$。若它处于最高次数，
它就是 $\langle I_{0,n,\beta}\rangle(\alpha_1,\ldots,\alpha_n)$ 乘以一点的类。
若不处于最高次数，则由它与上述除子的交确定。
而与 $\varphi$ 的像所给的除子相交时，得到 $\varphi^*I_{0,n,\beta}(\alpha_1,\ldots,\alpha_n)$。
由\IndexLocation{indexloc:0507:219}分裂公理，它由 $n$ 较小的 Gromov--Witten 类构成。
\SourcePage{220}
此外，等变性公理表明，与其他除子的交也具有同样的性质。
所以 $I_{0,n,\beta}(\alpha_1,\ldots,\alpha_n)$ 由 $n$ 较小的 Gromov--Witten 类确定。
由归纳假设，这些类由 Gromov--Witten 不变量确定，第一步证毕。

下一步是研究不变量 $\langle I_{0,n,\beta}\rangle$。
其思路是：上述除子之间的某些线性关系，给出不同 $n$、$\beta$ 所对应的
$\langle I_{0,n,\beta}\rangle$ 之间的二次关系。这些关系使我们能将任意 Gromov--Witten 不变量
用定理陈述中列出的不变量表示出来。尤其，定理中的不等式很容易解释。
若各 $\alpha_i$ 为齐次类，则 $\langle I_{0,3,\beta}\rangle(\alpha_1,\alpha_2,\alpha_3)=0$，除非
\[
\sum_i\deg\alpha_i=2\dim[\overline M_{0,3}(X,\beta)]^{\mathrm{virt}}
=2\left(-\int_\beta\omega_X+\dim X\right).
\]
结合显然成立的不等式 $\sum_i\deg\alpha_i\leq2\dim X+2$
（仍假设 $\deg\alpha_3=2$），就得到所需不等式。
论证的完整细节见 \SourceCite{KoM1}。下面说明当 $X=\mathbb P^2$ 时，这一方法如何运作。
\end{proof}

\subsection{射影平面的亏格零不变量}\label{subsec:plane-genus-zero}
\index{有理曲线!射影平面的计数|hyperpage}
\index{Gromov--Witten不变量@Gromov--Witten 不变量!射影平面|hyperpage}
我们将计算平面 $\mathbb P^2$ 的树级 Gromov--Witten 不变量。
这里 $\beta=d[\ell]$，其中 $\ell\subset\mathbb P^2$ 是一条直线，$d\geq0$。
为简化记号，用 $\langle I_{0,n,d}\rangle$ 代替 $\langle I_{0,n,\beta}\rangle$。
先计算 $\langle I_{0,3,1}\rangle([pt],[pt],[\ell])$。这很容易！
给定映射 $f:(\mathbb P^1,p_1,p_2,p_3)\to\mathbb P^2$，
使 $f(p_1)$ 与 $f(p_2)$ 是预先指定的点，且 $f(p_3)\in\ell$，便得到下图：
\begin{figure}[H]\centering
\includegraphics[width=0.46\textwidth]{7-4-plane-line.png}
\end{figure}
由于次数为 $1$，$f$ 的像就是由这两个点确定的直线。
又因为 $f(p_1)$、$f(p_2)$ 已固定，而 $f(p_3)$ 由交点 $\ell\cap f(\mathbb P^1)$ 确定，
我们得到一个\emph{唯一的映射}。随即容易看出，即使允许任意源曲线 $C$，
它也是唯一满足这些标记点条件的亏格零稳定映射 $f:(C,p_1,p_2,p_3)\to\mathbb P^2$。
此外，直接可见，这个稳定映射在 Gromov--Witten 不变量的定义~\eqref{eq:gw-moduli-integral} 中出现的重数为 $1$。
因此
\[
\langle I_{0,3,1}\rangle([pt],[pt],[\ell])=1.
\]
而且，依照定理~\ref{thm:gw-reconstruction}，这是唯一需要另行计算的不变量；
因为很容易检验，$\langle I_{0,3,1}\rangle([pt],[pt],[\ell])$ 是满足该定理条件的唯一非平凡不变量。

那么，其他不变量是什么？由\IndexLocation{indexloc:0505:220}点映射公理，可假设 $d\geq1$。
为计算 $\langle I_{0,n,d}\rangle(\alpha_1,\ldots,\alpha_n)$，可假设每个 $\alpha_i$ 都是
$[\mathbb P^2]$、$[\ell]$ 或 $[pt]$。又由式~\eqref{eq:gw-top-degree-condition}，有
\begin{equation}\label{eq:plane-degree-condition}
\sum_{i=1}^n\deg\alpha_i=6d-2+2n.
\end{equation}
当 $n=0,1$ 时，此方程无解；而当 $n=2,3$ 时，可以检验：不计置换，唯一的不变量是
\[
\langle I_{0,2,1}\rangle([pt],[pt])
=\langle I_{0,3,1}\rangle([pt],[pt],[\ell])=1.
\]
\SourcePage{221}
第一个等式由除子公理得到；更简单地说，$\langle I_{0,2,1}\rangle([pt],[pt])=1$，
因为在 $\mathbb P^2$ 中，经过两个点的直线唯一。

若 $n\geq4$，则由基本类公理，只要某个 $\alpha_i=[\mathbb P^2]$，就有
$\langle I_{0,n,d}\rangle(\alpha_1,\ldots,\alpha_n)=0$。
类似地，若 $n\geq4$ 且 $\alpha_n=[\ell]$，则由除子公理，
\[
\langle I_{0,n,d}\rangle(\alpha_1,\ldots,\alpha_{n-1},[\ell])
=d\langle I_{0,n-1,d}\rangle(\alpha_1,\ldots,\alpha_{n-1}),
\]
从而可以归纳计算。尚待计算的只有不变量 $\langle I_{0,n,d}\rangle([pt],\ldots,[pt])$。
由于 $\deg[pt]=4$，式~\eqref{eq:plane-degree-condition} 给出 $n=3d-1$。
因此，我们要计算的是
\begin{equation}\label{eq:plane-nd}
N_d=\langle I_{0,3d-1,d}\rangle(\underbrace{[pt],\ldots,[pt]}_{3d-1\text{ 次}}),\qquad d\geq1.
\end{equation}
由于 $\mathbb P^2$ 上有可迁的线性群作用，定理~\ref{thm:homogeneous-gw-equivalence} 表明，
$N_d$ 是 $\mathbb P^2$ 中经过 $3d-1$ 个一般位置点的 $d$ 次有理曲线的数目。

这些数 $N_d$ 也可定义为 $d$ 次有理平面曲线的 Severi 簇的次数。
然而，传统的代数几何方法此前只能给出最初几个数。
这里的新进展是：将 $N_d$ 解释为 Gromov--Witten 不变量，可以得到递推公式
\begin{equation}\label{eq:plane-recursion}
N_d=\sum_{\substack{d=d_1+d_2\\d_1,d_2>0}}N_{d_1}N_{d_2}
\left(d_1^2d_2^2\binom{3d-4}{3d_1-2}-d_1^3d_2\binom{3d-4}{3d_1-1}\right).
\end{equation}
由于 $N_1=\langle I_{0,2,1}\rangle([pt],[pt])=1$，此关系给出
\[
N_2=1,\quad N_3=12,\quad N_4=620,\quad N_5=87304,\quad\ldots.
\]
这说明稳定映射往往比嵌入曲线更容易处理，甚至在回答关于嵌入曲线的经典问题时也是如此！

式~\eqref{eq:plane-recursion} 的证明采用定理~\ref{thm:gw-reconstruction} 证明中概述的策略。
从下面的 Gromov--Witten 类出发：
\begin{equation}\label{eq:plane-xi}
\xi=I_{0,3d,d}(\underbrace{[pt],\ldots,[pt]}_{3d-2\text{ 次}},[\ell],[\ell])
\in H^*(\overline M_{0,3d},\mathbb Q).
\end{equation}
由次数公理，该类的次数为 $2\dim\overline M_{0,3d}-2$，所以它与一个除子的交是良定义的有理数。

幸运的是，$\overline M_{0,3d}$ 有许多有趣的除子。
用集合 $\mathcal I=\{1,\ldots,3d-4,p,q,r,s\}$ 为 $\overline M_{0,3d}$ 中稳定曲线的 $3d$ 个标记点编号。
用式~\eqref{eq:plane-xi} 的记号，这意味着 $\alpha_1=\cdots=\alpha_{3d-4}=\alpha_p=\alpha_q=[pt]$，
而 $\alpha_r=\alpha_s=[\ell]$。现在将 $\mathcal I$ 分成不相交的集合 $A$ 与 $B$。
与分裂公理中一样，这给出自然映射
\[
\varphi_{A,B}:\overline M_{0,|A|+1}\times\overline M_{0,|B|+1}\longrightarrow\overline M_{0,3d}.
\]
$\varphi_{A,B}$ 的像是 $\overline M_{0,3d}$ 上的一个除子。此外，还有下面这个基本的除子线性等价关系：
\begin{equation}\label{eq:plane-boundary-relation}
\sum_{\substack{r,s\in A\\p,q\in B}}\operatorname{Im}(\varphi_{A,B})
\sim\sum_{\substack{p,r\in A\\q,s\in B}}\operatorname{Im}(\varphi_{A,B}).
\end{equation}
\SourcePage{222}
这个关系来自自然映射 $\phi:\overline M_{0,3d}\to\overline M_{0,4}$，
它遗忘最初的 $3d-4$ 个点，并收缩不稳定分支。在此映射下，式~\eqref{eq:plane-boundary-relation}
的左边是 $\phi$ 在 $\overline M_{0,4}$ 中下面这条稳定曲线上的纤维：
\begin{figure}[H]\centering
\includegraphics[width=0.31250\textwidth]{7-4-four-point-boundary.png}
\end{figure}
类似地，式~\eqref{eq:plane-boundary-relation} 的右边是 $\phi$ 在另一条稳定曲线上的纤维，
其中 $p,r$ 位于同一分支，$q,s$ 位于另一分支。交比给出
$M_{0,4}\simeq\mathbb P^1\setminus\{0,1,\infty\}$，上图中的稳定曲线是
$\overline M_{0,4}\simeq\mathbb P^1$ 的三个边界点之一。
由于这些边界点显然线性等价，$\phi$ 在这些点上的纤维也线性等价，式~\eqref{eq:plane-boundary-relation} 因而成立。

证明的下一步，是将式~\eqref{eq:plane-boundary-relation} 的两边分别与式~\eqref{eq:plane-xi}
定义的 Gromov--Witten 类 $\xi$ 相交。先看左边。$\xi$ 与 $\operatorname{Im}(\varphi_{A,B})$
的交对应于拉回 $\varphi_{A,B}^*\xi$，可用分裂公理计算。
在\hyperref[sec:gw-properties]{第~\ref*{sec:gw-properties} 节}中，我们只对指标集 $\mathcal I$
的一个特定分划陈述了分裂公理；但由等变性公理，它适用于所有分划。
因此，由于 $\mathbb P^2\times\mathbb P^2$ 中的对角线类是
\[
\Delta=[pt]\otimes[\mathbb P^2]+[\ell]\otimes[\ell]+[\mathbb P^2]\otimes[pt],
\]
分裂公理与等变性公理表明，交数 $\xi\cdot\operatorname{Im}(\varphi_{A,B})$ 由下式给出：
\begin{equation}\label{eq:plane-splitting-intersection}
\begin{aligned}
\sum_{d=d_1+d_2}\Big(&\langle I_{0,|A|+1,d_1}\rangle(\alpha_i,\ i\in A,[pt])
\cdot\langle I_{0,|B|+1,d_2}\rangle(\alpha_i,\ i\in B,[\mathbb P^2])\\
&+\langle I_{0,|A|+1,d_1}\rangle(\alpha_i,\ i\in A,[\ell])
\cdot\langle I_{0,|B|+1,d_2}\rangle(\alpha_i,\ i\in B,[\ell])\\
&+\langle I_{0,|A|+1,d_1}\rangle(\alpha_i,\ i\in A,[\mathbb P^2])
\cdot\langle I_{0,|B|+1,d_2}\rangle(\alpha_i,\ i\in B,[pt])\Big),
\end{aligned}
\end{equation}
因为 $\varphi_{A,B}^*\xi$ 是最高次数的类。

虽然公式看起来复杂，但由 Gromov--Witten 不变量的次数条件，多数项为零。
例如，考虑不变量 $\langle I_{0,|A|+1,d_1}\rangle(\alpha_i,\ i\in A,\gamma)$，
其中 $\gamma$ 是 $[pt]$、$[\ell]$、$[\mathbb P^2]$ 之一。
由对 $A$ 的条件，$\alpha_r=\alpha_s=[\ell]$，其余 $\alpha_i=[pt]$。
所以式~\eqref{eq:gw-top-degree-condition} 表明，此不变量为零，除非
\[
|A|=3d_1+2-\tfrac12\deg\gamma.
\]
由于 $|A|$ 固定，这个方程确定了 $d_1$ 与 $\deg\gamma$。
因此，式~\eqref{eq:plane-splitting-intersection} 至多只有一项非零。

现假设 $|A|,|B|\geq3$。由基本类公理，
$\langle I_{0,|A|+1,d_1}\rangle(\alpha_i,\ i\in A,[\mathbb P^2])
=\langle I_{0,|B|+1,d_2}\rangle(\alpha_i,\ i\in B,[\mathbb P^2])=0$。
所以非零项必须满足 $\gamma=[\ell]$，进而 $|A|=3d_1+1$、$|B|=3d_2-1$。
于是式~\eqref{eq:plane-splitting-intersection} 化为
\[
\langle I_{0,3d_1+2,d_1}\rangle(\underbrace{[pt],\ldots,[pt]}_{3d_1-1\text{ 次}},[\ell],[\ell],[\ell])
\cdot\langle I_{0,3d_2,d_2}\rangle(\underbrace{[pt],\ldots,[pt]}_{3d_2-1\text{ 次}},[\ell]).
\]
\SourcePage{223}
由除子公理与 $N_d$ 的定义，上式化为 $d_1^3d_2N_{d_1}N_{d_2}$。因此
\begin{equation}\label{eq:plane-boundary-positive}
\xi\cdot\operatorname{Im}(\varphi_{A,B})=d_1^3d_2N_{d_1}N_{d_2},\qquad |A|=3d_1+1.
\end{equation}
还要注意，此时 $d_1>0$、$d_2>1$。

另一方面，若 $|B|=2$，请读者证明，式~\eqref{eq:plane-splitting-intersection} 的非零项为
\[
\langle I_{0,3d-1,d-1}\rangle(\underbrace{[pt],\ldots,[pt]}_{3d-4\text{ 次}},[\ell],[\ell],[\ell])
\cdot\langle I_{0,3,1}\rangle([pt],[pt],[\ell]).
\]
利用除子公理，它化为 $(d-1)^3N_{d-1}N_1$，
即 $d_1=d-1$、$d_2=1$ 时的式~\eqref{eq:plane-boundary-positive}。
最后，若 $|A|=2$，则非零项为
\[
\langle I_{0,3,0}\rangle([\ell],[\ell],[\mathbb P^2])
\cdot\langle I_{0,3d-1,d}\rangle(\underbrace{[pt],\ldots,[pt]}_{3d-2\text{ 次}},[pt]).
\]
由\IndexLocation{indexloc:0505:223}点映射公理，乘积中的第一个因子为 $1$，第二个因子为 $N_d$。于是
\begin{equation}\label{eq:plane-boundary-zero}
\xi\cdot\operatorname{Im}(\varphi_{A,B})=N_d,\qquad |A|=2,\ d_1=0,\ d_2=d.
\end{equation}
现在容易看出，类 $\xi$ 与式~\eqref{eq:plane-boundary-relation} 左边的交为
\[
N_d+\sum_{\substack{d=d_1+d_2\\d_1,d_2>0}}N_{d_1}N_{d_2}d_1^3d_2\binom{3d-4}{3d_1-1}.
\]
为说明原因，考虑指标集 $\mathcal I=\{1,\ldots,3d-4,p,q,r,s\}$ 的分划 $A,B$，
其中 $r,s\in A$、$p,q\in B$。当 $A,B$ 变化时，式~\eqref{eq:plane-boundary-zero}
恰好出现一次，即 $A=\{r,s\}$ 时。
而当 $d=d_1+d_2$、$d_1,d_2>0$ 时，式~\eqref{eq:plane-boundary-positive}
恰好出现 $\binom{3d-4}{3d_1-1}$ 次；因为 $|A|=3d_1+1$ 意味着，
除 $r,s$ 之外，$A$ 还有 $3d_1-1$ 个元素，选自 $\mathcal I$ 最初的 $3d-4$ 个元素。

为完成式~\eqref{eq:plane-recursion} 的证明，可用类似论证说明，
$\xi$ 与式~\eqref{eq:plane-boundary-relation} 右边的交为
\[
\sum_{\substack{d=d_1+d_2\\d_1,d_2>0}}N_{d_1}N_{d_2}d_1^2d_2^2\binom{3d-4}{3d_1-2}.
\]
细节留给读者，这是运用这些公理的一道很好的练习。
由线性等价关系~\eqref{eq:plane-boundary-relation}，$\xi$ 与两边的交相同，
递推关系~\eqref{eq:plane-recursion} 立即成立。

在上述由 Kontsevich 给出的论证中，请注意：将式~\eqref{eq:plane-boundary-relation}
与 $\xi$ 相交，一边产生了 $N_d$，另一边却没有。
这对式~\eqref{eq:plane-recursion} 的证明至关重要，而这一现象有更广泛的适用性。
事实上，\IndexLocation{indexloc:1208:223}重构定理的证明也以类似方式进行：将式~\eqref{eq:plane-boundary-relation}
与各种 Gromov--Witten 类相交，说明如何用定理~\ref{thm:gw-reconstruction}
中列出的特殊不变量表示一般的树级不变量。

还可用其他方法证明 $N_d$ 的递推关系。例如，\hyperref[chap:quantum-cohomology]{第~\ref*{chap:quantum-cohomology} 章}将看到，量子上同调这一套工具
如何使上述论证自动化并加以简化。尤其，最初的严格证明 \SourceCite{RT1}
使用了量子上同调的辛几何表述。
\SourcePage{224}
也能给出更几何化的论证，省去 Gromov--Witten 类的形式体系，而遵循上面所做的基本思路，见 \SourceCite{FP}。
另外，还可以用其他代数几何方法给出完全自足的论证，见 \SourceCite{Ran1,CH1,CH2,Vakil3}。
关于数 $N_d$ 的推广，见 \SourceCite{DI,EK1}。

最后还应提到，这些方法对其他具有可迁线性群作用的簇，也能给出很好的结果。
例如，\SourceCite{FP} 计算了 $\mathbb P^3$ 与光滑二次三维簇的树级 Gromov--Witten 不变量，
$\mathbb P^n$ 的情形则在 \SourceCite{Vakil3,GPa} 中讨论。其他例子见 \SourceCite{DI}。

\subsection{Calabi--Yau 三维簇的亏格零不变量}\label{subsec:cy-genus-zero}
\index{Clemens猜想@Clemens 猜想|hyperpage}
我们将讨论 Calabi--Yau 三维簇的树级 Gromov--Witten 不变量。
这里 $V$ 表示一个光滑 Calabi--Yau 三维簇。对非零 $\beta\in H_2(V,\mathbb Z)$，置 $N_\beta=\langle I_{0,0,\beta}\rangle$。
\begin{proposition}\label{prop:cy-gw-determined}
$V$ 的亏格零 Gromov--Witten 不变量完全由各 $N_\beta$ 及 $V$ 上的交数确定。
\end{proposition}
\begin{proof}
只需考虑各 $\alpha_i$ 都是齐次类的不变量 $\langle I_{0,n,\beta}\rangle(\alpha_1,\ldots,\alpha_n)$。
由\IndexLocation{indexloc:0498:224}次数公理，可假设 $\sum_i\deg\alpha_i=2n$。
由于每个 $\alpha_i$ 或为基本类，或满足 $\deg\alpha_i\geq2$，需考虑两种情形：
\begin{itemize}
\item 某个 $\alpha_i$ 是基本类；
\item 所有 $i$ 都满足 $\deg\alpha_i=2$。
\end{itemize}
在第一种情形，由基本类公理、\IndexLocation{indexloc:0501:224}等变性公理及点映射公理，Gromov--Witten 不变量是一个交数，或为 $0$。
在第二种情形，当 $\beta\ne0$ 时，除子公理给出
\[
\langle I_{0,n,\beta}\rangle(\alpha_1,\ldots,\alpha_n)=N_\beta\int_\beta\alpha_1\cdots\int_\beta\alpha_n.
\]
而当 $\beta=0$ 时，式~\eqref{eq:gw-degree-zero-three-point} 表明，$n\ne3$ 时 $\langle I_{0,n,0}\rangle=0$，且
\[
\langle I_{0,3,0}\rangle(\alpha_1,\alpha_2,\alpha_3)=\int_V\alpha_1\cup\alpha_2\cup\alpha_3.
\]
命题证毕。
\end{proof}
\SourceExampleNumber{7.4.3}{1}
\begin{example}\label{ex:quintic-gw-types}
设 $V\subset\mathbb P^4$ 是一个一般的五次三维簇，取直线 $\ell\subset V$
（这样的直线共有 $2875$ 条）。用 $H\in H^2(V,\mathbb Z)$ 表示超平面类。
那么，不计置换，唯一可能非零的 Gromov--Witten 不变量是
\[
\langle I_{0,3,0}\rangle(\alpha_1,\alpha_2,1),\qquad\langle I_{0,n,d}\rangle(H,\ldots,H).
\]
其中，为简化记号，用 $\langle I_{0,n,d}\rangle$ 代替 $\langle I_{0,n,d[\ell]}\rangle$，而 $n,d\geq0$ 任意。
由基本类公理，$\langle I_{0,3,0}\rangle(\alpha_1,\alpha_2,1)=\int_V\alpha_1\cup\alpha_2$。
置 $N_d=\langle I_{0,0,d}\rangle$，则当 $d>0$ 时，除子公理给出
$\langle I_{0,n,d}\rangle(H,\ldots,H)=d^nN_d$。
最后，由点映射公理，$\langle I_{0,3,0}\rangle(H,H,H)=\int_VH^3=5$。
\end{example}
对于五次三维簇 $V$，上述例子将求 Gromov--Witten 不变量的问题归结为确定
$N_d=\langle I_{0,0,d}\rangle$。由例题~\ref{ex:quintic-gw-euler-integral}，
\[
N_d=\langle I_{0,0,d}\rangle=\int_{\overline M_{0,0}(\mathbb P^4,d)}c_{5d+1}(\mathcal V_d).
\]
\SourcePage{225}
当然，我们希望用 $V$ 上的有理曲线来计算 $N_d$。
由于 Clemens 猜想以及多重覆盖的存在，这仍是一个未解决的问题。
我们将分别讨论这两点，并说明它们与计算 $N_d$ 的问题有何关系。

如\hyperref[chap:quintic]{第~\ref*{chap:quintic} 章}所讨论，镜像对称的一个引人注目的方面，是它对五次三维簇 $V$ 上的有理曲线作出的预测。
困难在于，我们一般并不知道 $V$ 上给定次数的有理曲线是否只有有限多条。
这就引出了如下的 Clemens 猜想。
\index{Clemens猜想@Clemens 猜想|hyperpage}
\begin{conjecture}\label{conj:clemens}
设 $V\subset\mathbb P^4$ 是一个一般的五次三维簇。则对每个次数 $d\geq1$：
\begin{enumerate}[label=(\roman*)]
\item $V$ 中 $d$ 次不可约有理曲线 $C\subset V$ 只有有限多条。
\item 遍历所有次数时，这些曲线彼此不相交。
\item 若 $f:\mathbb P^1\to C$ 是不可约有理曲线 $C\subset V$ 的正规化，
则法丛 $N_f$ 同构于 $\mathcal O_{\mathbb P^1}(-1)\oplus\mathcal O_{\mathbb P^1}(-1)$。
\end{enumerate}
\end{conjecture}
注意，这里并未断言有理曲线都是光滑的。
那样的断言是错误的，因为一般 $V$ 中存在具有六个结点的有理平面五次曲线 $C\subset V$，见 \SourceCite{Vainsencher}。
目前已知，$V$ 上次数不超过 $9$ 的有理曲线只有有限多条。
而且，正如猜想预测的，它们都彼此不相交，且除上述平面五次曲线外，都非奇异。
此外，光滑曲线的法丛都是 $\mathcal O_{\mathbb P^1}(-1)\oplus\mathcal O_{\mathbb P^1}(-1)$。
这些结果均在 \SourceCite{Katz2,JK1} 中证明。
\hyperref[subsec:nodal-quintic-double-covers]{第~\ref*{subsec:nodal-quintic-double-covers} 小节}将说明，$V$ 中的奇异五次曲线也具有预测的法丛，因此猜想~\ref{conj:clemens} 对 $d\leq9$ 成立。

但即使 Clemens 猜想对所有 $d$ 都成立，仍有\emph{多重覆盖}的问题。
设 $C\subset V$ 是一条 $d/k$ 次不可约有理曲线，其正规化映射为 $f:\mathbb P^1\to V$。
将 $f$ 与任意 $k$ 次映射 $g:\mathbb P^1\to\mathbb P^1$ 复合，便得到稳定映射
$f\circ g:\mathbb P^1\to V$，满足 $(f\circ g)_*[\mathbb P^1]=k(d/k)=d$。
由于 $g$ 在一个族中变化，$k>1$ 时就得到正维的稳定映射族。
这说明 $\overline M_{0,0}(V,d)$ 为正维，尽管由于 $V$ 是 Calabi--Yau 三维簇，预期维数~\eqref{eq:stable-map-expected-dimension} 为 $0$。

我们需要确定上述多重覆盖族对虚拟基本类 $[\overline M_{0,0}(V,d)]^{\mathrm{virt}}$，从而对 $N_d$ 的贡献。
\hyperref[subsec:instanton-numbers]{第~\ref*{subsec:instanton-numbers} 小节}将看到：
若 $C\subset V$ 是 $d/k$ 次光滑有理曲线，法丛为
$N\simeq\mathcal O_{\mathbb P^1}(-1)\oplus\mathcal O_{\mathbb P^1}(-1)$，
则其 $k$ 次多重覆盖对 $N_d$ 的贡献为 $k^{-3}$。

暂且作一个朴素假设：对所有 $d$，$V$ 上的 $d$ 次有理曲线只有有限多条，
且全都光滑、彼此不相交、法丛为 $\mathcal O_{\mathbb P^1}(-1)\oplus\mathcal O_{\mathbb P^1}(-1)$。
若 $n_d$ 是这样的 $d$ 次曲线的总数，则上述讨论提示
\begin{equation}\label{eq:quintic-multiple-cover}
N_d=\sum_{k\mid d}n_{d/k}k^{-3},
\end{equation}
因为每条 $d/k$ 次有理曲线对 $N_d$ 都贡献 $k^{-3}$。
利用\IndexLocation{indexloc:0499:225}除子公理，式~\eqref{eq:quintic-multiple-cover} 给出
\begin{equation}\label{eq:quintic-three-divisors}
\langle I_{0,3,d}\rangle(H,H,H)=d^3\langle I_{0,0,d}\rangle=d^3N_d
=\sum_{k\mid d}n_{d/k}\left(\frac d k\right)^3=\sum_{k\mid d}n_kk^3,
\end{equation}
\SourcePage{226}
其中 $d\geq1$。这使我们能与\hyperref[chap:quintic]{第~\ref*{chap:quintic} 章}的公式建立如下联系。
回忆式~\eqref{eq:quintic-a}，A 模型关联函数为
\begin{equation}\label{eq:quintic-a-instanton}
\langle H,H,H\rangle=5+\sum_{d=1}^{\infty}n_dd^3\frac{q^d}{1-q^d}.
\end{equation}
由于此式右边可以写成
\[
5+\sum_{d=1}^{\infty}\left(\sum_{k\mid d}n_kk^3\right)q^d,
\]
上面关于 $\langle I_{0,3,d}\rangle(H,H,H)$ 的公式表明
\begin{equation}\label{eq:quintic-a-gw}
\langle H,H,H\rangle=\sum_{d=0}^{\infty}\langle I_{0,3,d}\rangle(H,H,H)q^d,
\end{equation}
这里用到了例题~\ref{ex:quintic-gw-types} 的 $\langle I_{0,3,0}\rangle(H,H,H)=5$。

这是一个重要进展：虽然上述多数等式依赖各种假设，但给出 $\langle H,H,H\rangle$
的式~\eqref{eq:quintic-a-gw} 有严格的定义。这意味着，我们终于扎实地理解了五次三维簇的 A 模型关联函数。
此外，\hyperref[chap:mirror-theorem]{第~\ref*{chap:mirror-theorem} 章}证明式~\eqref{eq:quintic-mirror-identity} 中 $\langle H,H,H\rangle$ 的精彩公式时，将使用这个严格定义的 $\langle H,H,H\rangle$。

从上述讨论可以得出两个结论。第一，Gromov--Witten 不变量 $N_d=\langle I_{0,0,d}\rangle$
正是表述五次三维簇的镜像对称所需要的量。
第二，我们仍有\emph{大量}工作要做，因为还需要在不依赖上述所有假设的情况下，将 $N_d$ 与 $V$ 上的有理曲线联系起来。
因此，应研究式~\eqref{eq:quintic-multiple-cover} 真正意味着什么。我们很快会看到，答案比预想的复杂。

\subsection{Calabi--Yau 三维簇的瞬子数}\label{subsec:instanton-numbers}
\index{瞬子数|hyperpage}\index{多重覆盖公式|hyperpage}
\index{瞬子数|hyperpage}\index{多重覆盖|hyperpage}
为更好地理解 $N_d$ 与五次三维簇上有理曲线的关系，我们采取如下策略：
先定义\emph{瞬子数} $n_d$，使式~\eqref{eq:quintic-multiple-cover} 自动成立；
再处理更精细的问题，将 $n_d$ 与 $V$ 上 $d$ 次有理曲线的实际数目联系起来。

考虑如下的一般情形。设 $V$ 是光滑 Calabi--Yau 三维簇，$\beta\in H_2(V,\mathbb Z)$ 是非零有效类。
那么，容易看出存在唯一的有理数 $n_\beta$，称为\emph{瞬子数}，使
\begin{equation}\label{eq:instanton-definition}
N_\beta=\langle I_{0,0,\beta}\rangle=\sum_{\beta=k\gamma}n_\gamma k^{-3}.
\end{equation}
求和遍历满足 $\beta=k\gamma$ 的所有正整数 $k$ 与非零有效类 $\gamma\in H_2(V,\mathbb Z)$。
$n_\beta$ 的存在唯一性如下证明：若 $\beta$ 在 $H_2(V,\mathbb Z)$ 中不可整除，直接置
$n_\beta=\langle I_{0,0,\beta}\rangle$；否则，利用式~\eqref{eq:instanton-definition} 归纳定义 $n_\beta$。
这证明了存在性，唯一性也随即得到。

有了瞬子数 $n_\beta$ 的这一定义，首先要说明式~\eqref{eq:instanton-definition} 中因子 $k^{-3}$ 的由来。
纯粹从形式上看，可以说它给出了漂亮的公式。例如，设 $\beta\in H_2(V,\mathbb Z)$
是非零有效类，$\alpha_1,\alpha_2,\alpha_3\in H^2(V,\mathbb C)$。
\SourcePage{227}
由式~\eqref{eq:instanton-definition} 与除子公理，
\begin{equation}\label{eq:instanton-three-divisors}
\begin{aligned}
\langle I_{0,3,\beta}\rangle(\alpha_1,\alpha_2,\alpha_3)
&=N_\beta\int_\beta\alpha_1\int_\beta\alpha_2\int_\beta\alpha_3\\
&=\left(\sum_{\beta=k\gamma}n_\gamma k^{-3}\right)
\int_\beta\alpha_1\int_\beta\alpha_2\int_\beta\alpha_3\\
&=\sum_{\beta=k\gamma}n_\gamma\int_\gamma\alpha_1\int_\gamma\alpha_2\int_\gamma\alpha_3.
\end{aligned}
\end{equation}
最后一个等式来自 $\beta=k\gamma$ 时的 $k^{-1}\int_\beta u=\int_\gamma u$。
注意，这推广了式~\eqref{eq:quintic-three-divisors}。

\hyperref[chap:quantum-cohomology]{第~\ref*{chap:quantum-cohomology} 章}将严格表述 A 模型关联函数这一数学概念。对于 Calabi--Yau 三维簇 $V$，关联函数为
\[
\langle\alpha_1,\alpha_2,\alpha_3\rangle
=\sum_\beta\langle I_{0,3,\beta}\rangle(\alpha_1,\alpha_2,\alpha_3)q^\beta,
\]
其中 $\alpha_i\in H^2(V,\mathbb C)$，而 $q^\beta=e^{2\pi i\int_\beta\omega}$，
$\omega$ 是复化 Kähler 类。利用点映射公理与式~\eqref{eq:instanton-three-divisors}，得到
\begin{equation}\label{eq:cy-a-correlation-instanton}
\langle\alpha_1,\alpha_2,\alpha_3\rangle
=\int_V\alpha_1\cup\alpha_2\cup\alpha_3
+\sum_{\substack{\beta\ne0\\\beta\text{ 有效}}}n_\beta\frac{q^\beta}{1-q^\beta}
\int_\beta\alpha_1\int_\beta\alpha_2\int_\beta\alpha_3.
\end{equation}
推导方式与式~\eqref{eq:quintic-a-gw} 类似。此公式在镜像对称中起重要作用，
并与量子上同调密切相关。\hyperref[subsec:a-model-correlation]{第~\ref*{subsec:a-model-correlation} 小节}将更详细地讨论它。

式~\eqref{eq:instanton-definition} 中的因子 $k^{-3}$ 还有更深刻的理由。
\hyperref[subsec:cy-genus-zero]{第~\ref*{subsec:cy-genus-zero} 小节}已略作提及，现在作更充分的说明。
设 $C\subset V$ 为光滑有理曲线，法丛为
$N\simeq\mathcal O_{\mathbb P^1}(-1)\oplus\mathcal O_{\mathbb P^1}(-1)$，
并置 $\gamma=[C]\in H_2(V,\mathbb Z)$。
我们要计算 $C$ 的 $k$ 次多重覆盖对 Gromov--Witten 不变量
$N_\beta=\langle I_{0,0,\beta}\rangle$ 的贡献，其中 $\beta=k\gamma$。

固定同构 $f:\mathbb P^1\simeq C$。则 $k$ 次覆盖 $g:\mathbb P^1\to\mathbb P^1$
给出稳定映射 $f\circ g\in\overline M_{0,0}(V,\beta)$。
由于 $g$ 自然属于 $\overline M_{0,0}(\mathbb P^1,k)$，得到嵌入
$\overline M_{0,0}(\mathbb P^1,k)\hookrightarrow\overline M_{0,0}(V,\beta)$。
记这个嵌入的像为 $M_{k,C}(V)$。我们断言，$M_{k,C}(V)$ 实际上是 $\overline M_{0,0}(V,\beta)$ 的一个连通分支。

为说明这一点，考虑在连通参数空间 $S$ 上、同调类为 $\beta$ 的任意亏格零稳定映射族
$f:\mathcal C\to V$，使 $0\in S$ 上的稳定映射属于 $M_{k,C}(V)$。
记结构态射为 $\pi:\mathcal C\to S$。关于 $C$ 的法丛 $N$ 的假设保证存在解析收缩映射
$\rho:V\to V'$，将 $C$ 收缩为一点 $p$，且在 $C$ 与 $p$ 之外是同构。
考虑复合映射 $\widetilde f=\rho\circ f$，注意 $\widetilde f(\pi^{-1}(0))$ 是点 $p$。
由 \SourceCite[引理 1.5]{CKM}，容易得到，对所有 $s\in S$，$\widetilde f(\pi^{-1}(s))$ 都是一点。
由于 $f_*[\pi^{-1}(s)]=\beta$，该点也必须是 $p$。断言得证。

虚拟基本类 $[\overline M_{0,0}(V,\beta)]^{\mathrm{virt}}$ 是零维的；
按定义，$\langle I_{0,0,\beta}\rangle$ 是它的次数。
曲线 $C$ 对 $\langle I_{0,0,\beta}\rangle$ 的贡献如下。
\begin{theorem}\label{thm:multiple-cover}
在上述关于 $C$ 的假设下，设 $\beta=k\gamma=k[C]$，
并令 $M_{k,C}(V)$ 为 $\overline M_{0,0}(V,\beta)$ 中包含 $C$ 的所有 $k$ 次多重覆盖的分支。
\SourcePage{228}
则虚拟基本类 $[\overline M_{0,0}(V,\beta)]^{\mathrm{virt}}$ 中支撑在 $M_{k,C}(V)$ 上的部分，
对 $N_\beta=\langle I_{0,0,\beta}\rangle$ 的贡献为 $k^{-3}$。
\end{theorem}
\begin{proof}
从历史上看，这一结果最初表述为以下断言：$C$ 的 $k$ 次多重覆盖对
$\langle I_{0,3,\beta}\rangle(\alpha_1,\alpha_2,\alpha_3)$ 的贡献是
\begin{equation}\label{eq:multiple-cover-three-point}
\int_\gamma\alpha_1\int_\gamma\alpha_2\int_\gamma\alpha_3,
\end{equation}
其中 $\alpha_i\in H^2(V,\mathbb C)$。
这一断言与定理~\ref{thm:multiple-cover} 的等价性可由式~\eqref{eq:instanton-three-divisors} 的推导得到。
\SourceCite{CdGP} 写出三点函数~\eqref{eq:quintic-a-instanton} 的方式，已经隐含地提出了这一猜想；
Aspinwall 与 Morrison 在 \SourceCite{AM1} 中给出了证明概要。
现在必须将他们的论证视为不完整的，因为它没有与后来任何一种 Gromov--Witten 不变量的定义作直接比较。

下面概述 Voisin 的证明 \SourceCite{Voisin2}。
与\hyperref[sec:gw-symplectic-definition]{第~\ref*{sec:gw-symplectic-definition} 节}一样，
将 $V$ 的复结构形变为一般的几乎复结构 $J$，再加入一般的非齐次项 $\nu$。
考虑扰动的 $J$ 全纯映射 $g:\mathbb P^1\to V$，满足 $g_*[\mathbb P^1]=\beta$。
用该节的记号，它们构成空间 $\mathcal M(\mathbb P^1,J,\nu,\beta)$。
若 $J$ 足够接近原复结构、$\nu$ 足够小，而 $U$ 是 $V$ 中 $C$ 的一个足够小的邻域，
则由像位于 $U$ 中的映射构成的子集 $\mathcal M^U(\mathbb P^1,J,\nu,\beta)$
是 $\mathcal M(\mathbb P^1,J,\nu,\beta)$ 的一个分支（未必连通）。
固定三个点 $p_1,p_2,p_3\in\mathbb P^1$，得到赋值映射
\[
\pi_1:\mathcal M^U(\mathbb P^1,J,\nu,\beta)\longrightarrow V^3,
\qquad g\longmapsto(g(p_1),g(p_2),g(p_3)).
\]
此映射的像为六维，具有自然定向，并能紧化为边界维数至多为四的空间。
因此，$\pi_1$ 的像给出 $H_6(V^3,\mathbb Z)$ 中一个良定义的同调类。
\SourceCite{Voisin2} 的主要结果是：由于 $C$ 的嵌入是刚性的，该类恰好是 $C\times C\times C$ 的同调类。

假定这一结果，取一般位置的 $\Gamma_1,\Gamma_2,\Gamma_3$，分别为
$\alpha_1,\alpha_2,\alpha_3$ 的 Poincaré 对偶。
由于 $C\times C\times C$ 与 $\Gamma_1\times\Gamma_2\times\Gamma_3$ 相交的点数
由式~\eqref{eq:multiple-cover-three-point} 给出（记得 $[C]=\gamma$），
$\Gamma_1\times\Gamma_2\times\Gamma_3$ 与 $\pi_1$ 的像相交也得到同样的数。
于是，辛 Gromov--Witten 不变量的计数定义表明，
对 $\langle I_{0,3,\beta}\rangle(\alpha_1,\alpha_2,\alpha_3)$ 的贡献正是式~\eqref{eq:multiple-cover-three-point}。

定理~\ref{thm:multiple-cover} 的一些较新证明直接处理 $N_\beta$，说明 $C$ 的 $k$ 次覆盖贡献 $k^{-3}$。
\hyperref[subsec:multiple-cover-localization]{第~\ref*{subsec:multiple-cover-localization} 小节}将介绍 Kontsevich 的证明方法；我们将详细处理 $d=2$ 的情形，
一般情形则参见 \SourceCite{Manin2}。
不过，该小节还将给出 Pandharipande 的一个简化论证。
此定理的其他证明见 \SourceCite{LLY,Gathmann2}；\hyperref[sec:quintic-mirror-theorem]{第~\ref*{sec:quintic-mirror-theorem} 节}将讨论 \SourceCite{LLY} 的证明。
\end{proof}
接下来讨论两个例子，说明瞬子数 $n_\beta$ 的枚举意义有时相当微妙。
即使对于五次三维簇，事情也并非如人们所希望的那样简单。

\SourceExampleNumber{7.4.4}{1}
\begin{example}\label{ex:quintic-instanton-enumerative}
设 $V\subset\mathbb P^4$ 是一个一般的五次三维簇。
将瞬子数 $n_\beta$ 记为 $n_d$，其中 $\beta=d[\ell]$，$\ell\subset V$ 为直线。
这些数正是式~\eqref{eq:quintic-multiple-cover}、\eqref{eq:quintic-a-instanton} 中出现的数，
并由镜像对称通过\hyperref[chap:quintic]{第~\ref*{chap:quintic} 章}的公式计算。
\SourcePage{229}
当 $d\leq9$ 时，$n_d$ 与 $d$ 次有理曲线的关系尤其好：
\begin{equation}\label{eq:quintic-instanton-small-degree}
d\leq9\quad\Longrightarrow\quad n_d=V\text{ 中 }d\text{ 次有理曲线的数目}.
\end{equation}
由于 Clemens 猜想对 $d\leq9$ 成立，此式可由定理~\ref{thm:multiple-cover} 得到。

下面更详细地解释原因。由于猜想~\ref{conj:clemens} 在 $d\leq9$ 时成立，
式~\eqref{eq:quintic-instanton-small-degree} 的右边有限。
用 $n'_d$ 表示 $V$ 中 $d$ 次有理曲线的数目。
如\hyperref[subsec:cy-genus-zero]{第~\ref*{subsec:cy-genus-zero} 小节}所述，
$d\leq9$ 时，$V$ 中所有 $d$ 次有理曲线都光滑（$d=5$ 时有例外），
且法丛为 $\mathcal O_{\mathbb P^1}(-1)\oplus\mathcal O_{\mathbb P^1}(-1)$，见 \SourceCite{JK1}。

为证明式~\eqref{eq:quintic-instanton-small-degree}，先假设 $d\ne5$，
并取 $\overline M_{0,0}(V,d)$ 中的稳定映射 $f:C'\to V$。
由于次数不超过 $9$ 的所有曲线彼此不相交，$C=f(C')$ 必为 $d/k$ 次光滑曲线，
其中 $k$ 是 $f:C'\to C$ 的次数。
用定理~\ref{thm:multiple-cover} 的记号，$f$ 位于 $\overline M_{0,0}(V,d)$ 的分支 $M_{k,C}(V)$ 中。
这说明 $\overline M_{0,0}(V,d)$ 是这些分支的不交并，于是定理~\ref{thm:multiple-cover} 给出
\[
N_d=\sum_{k\mid d}n'_{d/k}k^{-3}.
\]
由 $n_d$ 的唯一性，$n'_d=n_d$，所以式~\eqref{eq:quintic-instanton-small-degree} 对 $d\ne5$ 得证。

再设 $d=5$。由 \SourceCite{Vainsencher}，$V$ 中有 $17\,601\,000$ 条具有六个结点的五次有理曲线。
现取 $\overline M_{0,0}(V,5)$ 中的稳定映射 $f:C'\to V$。
此时，$f(C')$ 或为直线，或为五次光滑曲线，或为具有六个结点的五次有理曲线。
因此，$\overline M_{0,0}(V,5)$ 分解为三种相应的分支；定理~\ref{thm:multiple-cover} 适用于前两种。
余下只需考虑 $f:C'\to C$，其中 $C$ 是 $V$ 中具有六个结点的有理曲线。
容易看出，$f$ 是正规化映射 $f:\mathbb P^1\to C$。
此外，\hyperref[subsec:nodal-quintic-double-covers]{第~\ref*{subsec:nodal-quintic-double-covers} 小节}将证明 $f$ 的法丛为
$N_f=\mathcal O_{\mathbb P^1}(-1)\oplus\mathcal O_{\mathbb P^1}(-1)$。
这表明，$\overline M_{0,0}(V,5)$ 中对应于 $C$ 的分支是一个光滑点，
且它对虚拟基本类的贡献为 $1$，与定理~\ref{thm:multiple-cover} 相符。
于是，上一段的论证给出 $n'_5=n_5$，完成式~\eqref{eq:quintic-instanton-small-degree} 的证明。

基于式~\eqref{eq:quintic-instanton-small-degree}，自然会希望：
只要证明更高次数 $d$ 的 Clemens 猜想，就能自动推出相应瞬子数是\IndexLocation{indexloc:1193:229}\IndexLocation{indexloc:1199:229}\IndexLocation{indexloc:1200:229}五次三维簇上 $d$ 次有理曲线的数目。
长期以来，许多人相信这一点；然而，R. Pandharipande 最近的一个观察表明，这种希望过于朴素。
事实上，对于 $d=10$——猜想~\ref{conj:clemens} 中第一个尚未解决的情形——有下面关于 $n_{10}$ 的猜想公式：
\begin{equation}\label{eq:quintic-degree-ten}
\begin{gathered}
\text{若 Clemens 猜想对所有 }d\leq10\text{ 成立，则}\\
n_{10}=6\times17\,601\,000+V\text{ 中 }10\text{ 次有理曲线的数目}.
\end{gathered}
\end{equation}
\hyperref[subsec:nodal-quintic-double-covers]{第~\ref*{subsec:nodal-quintic-double-covers} 小节}将仔细证明此式。
关键在于，$V$ 中那 $17\,601\,000$ 条具有六个结点的五次曲线的二重覆盖，造成了复杂情形。
更确切地说，该小节将证明，每条这样的曲线 $C$ 对 $N_{10}$ 的贡献是 $6\tfrac18$，而非 $\tfrac18$。
我们会看到，这是因为二重覆盖 $f:C'\to C$ 有两种：
一种经由正规化映射分解；另一种的 $C'$ 可约，且 $f$ 在六个结点中某一个的邻域内是局部同构。
\SourcePage{230}
假设猜想~\ref{conj:clemens} 对 $d\leq10$ 成立，且每条上述 $C$ 对 $N_{10}$ 都贡献 $6\tfrac18$，
则式~\eqref{eq:quintic-degree-ten} 可如下推导。首先，由定义，
\[
N_{10}=n_1 10^{-3}+n_2 5^{-3}+n_5 2^{-3}+n_{10}.
\]
再利用 $d\leq9$ 时的 $n_d=n'_d$，以及证明式~\eqref{eq:quintic-instanton-small-degree} 时的方法，得到
\[
N_{10}=n_1 10^{-3}+n_2 5^{-3}+(n_5-17\,601\,000)2^{-3}
+17\,601\,000\cdot6\tfrac18+n'_{10},
\]
其中 $n'_{10}$ 是 $10$ 次有理曲线的数目。
将这两个等式相减，容易得到式~\eqref{eq:quintic-degree-ten}。
\end{example}

\hyperref[chap:mirror-theorem]{第~\ref*{chap:mirror-theorem} 章}证明五次三维簇 $V$ 的镜像定理后，我们就能按需要计算任意多个瞬子数 $n_d$。
但从前一个例子可见，这离知道 $V$ 上 $d$ 次有理曲线的数目仍很远。
首先，对于 $d\geq10$，证明 Clemens 猜想是一个非平凡的问题
（关于 $10\leq d\leq24$ 的一些初步进展，见 \SourceCite{JK2}）；
然后还要面对这些曲线并非全都光滑的困难。
如上所见，这类曲线的多重覆盖，并非总按定理~\ref{thm:multiple-cover} 的方式作出贡献。
事实上，只要 $d\geq10$ 是 $5$ 的倍数，就可能出现这一问题。
所以，即使\hyperref[chap:quintic]{第~\ref*{chap:quintic} 章}的公式已经全部得到证明，五次三维簇仍提出了许多未解决的问题。

下面给出一个瞬子数更微妙的例子。这里将看到，除多重覆盖之外，虚拟基本类还会以其他方式发挥作用。
\SourceExampleNumber{7.4.4}{2}
\begin{example}\label{ex:local-plane-instantons}
设 $\widehat V\subset\mathbb P(1,1,1,6,9)$ 是由下式定义的 $18$ 次超曲面：
\[
x_0^{18}+x_1^{18}+x_2^{18}+x_3^3+x_4^2=0.
\]
可以证明，$\widehat V$ 在 $(0,0,0,-1,1)$ 有唯一奇点，解消此奇点得到 Calabi--Yau 三维簇 $V$。
例外集是 $V$ 中的一个除子 $E\simeq\mathbb P^2$。
一般情形也是如此；而 $\mathbb P^2\subset V$ 表明，$V$ 有无穷多条各个次数的有理曲线。

为计算 Gromov--Witten 不变量，需要描述 $V$ 上的一些类。
依照 \SourceCite{CFKM}，令 $L,H$ 分别为由 $1$ 次与 $3$ 次方程定义的 $V$ 上的除子类。
由于爆破，有 $3L+E\in|H|$；可以证明，$L,H$ 张成 $H^2(V,\mathbb Z)$ 并生成 Kähler 锥。
此外，$\ell=L\cap E$ 是 $E\simeq\mathbb P^2$ 中的直线，
$h=L_1\cap L_2$ 是椭圆曲线，其中 $L_i\in|L|$ 为一般成员。
于是 $\ell,h$ 生成 $H^4(V,\mathbb Z)\simeq H_2(V,\mathbb Z)$。

在这一情形中，用 $\langle I_{0,3,r,s}\rangle$ 表示 $\langle I_{0,3,[rh+s\ell]}\rangle$。
则式~\eqref{eq:instanton-three-divisors} 变成
\[
\langle I_{0,3,r,s}\rangle(\alpha_1,\alpha_2,\alpha_3)
=\sum_{k\mid\gcd(r,s)}n\!\left(\frac r k,\frac s k\right)
\int_{\frac r kh+\frac s k\ell}\alpha_1
\int_{\frac r kh+\frac s k\ell}\alpha_2
\int_{\frac r kh+\frac s k\ell}\alpha_3,
\]
其中 $\alpha_1,\alpha_2,\alpha_3\in H^2(V,\mathbb C)$，
$n(r,s)$ 是类为 $C\sim rh+s\ell$ 的有理曲线 $C\subset V$ 的瞬子数。
为看看会得到怎样的数，考虑
\[
\langle I_{0,3,0,2}\rangle(L,L,L)=n(0,1)+n(0,2)2^3.
\]
由 \SourceCite{CFKM}，镜像对称预测 $n(0,1)=3$、$n(0,2)=-6$，
从而 $\langle I_{0,3,0,2}\rangle(L,L,L)=-45$。
这些数也可由定理~\ref{thm:givental-toric-mirror} 得到。因此，我们得到一个负的 Gromov--Witten 不变量。
\SourcePage{231}
这些数可作如下解释。首先，$C\sim\ell$ 表明 $C$ 是 $\mathbb P^2$ 中的直线。
这类 $C$ 的模空间是对偶的 $\mathbb P^2$。
为计算它对 $n(0,1)$ 的贡献，回忆\hyperref[subsec:normal-cone]{第~\ref*{subsec:normal-cone} 小节}
关于虚拟基本类的讨论。在命题~\ref{prop:excess-obstruction-euler-class} 的情形下，
构造归结为求障碍丛 Euler 类（即最高 Chern 类）的次数。
又因为这里处理的是嵌入曲线，直线 $C\subset\mathbb P^2$ 的障碍空间为 $H^1(N)$，
其中 $N$ 是 $C$ 在 $V$ 中的法丛。但 $H^1(N)$ 通过配对
\[
H^1(N)\otimes H^0(N)\longrightarrow H^1(\wedge^2N)
\simeq H^1(T_C^*)=\mathbb C
\]
与 $H^0(N)$ 对偶，所以全局的障碍丛是 $T^*_{\mathrm{moduli}}=T^*_{\mathbb P^2}$。
因此瞬子数为
\[
n(0,1)=\int_{\mathbb P^2}c_2(T^*_{\mathbb P^2})
=(-1)^{\dim\mathbb P^2}\chi(\mathbb P^2)=3.
\]
类似地，$C\sim2\ell$ 表明 $C$ 是 $\mathbb P^2$ 中的二次曲线。
二次曲线的模空间是 $\mathbb P^5$。由于它光滑，与上面相同的计算给出
\[
n(0,2)=\int_{\mathbb P^5}c_5(T^*_{\mathbb P^5})
=(-1)^{\dim\mathbb P^5}\chi(\mathbb P^5)=-6.
\]
这就是 \SourceCite{CFKM} 中的论证，但它并不完整：满足 $f_*[C]=2\ell$ 的稳定映射
$f:C'\to V$ 的模空间，与这里所用的二次曲线 Hilbert 概形不同。

$n(0,1)$ 也可在辛几何框架中计算，见 \SourceCite{Ruan1}。
对 $n(0,2)$ 的朴素计算得到正确数值，提示非孤立曲线可能存在一种多重覆盖公式的推广，
使我们能用浸入曲线，而不是稳定映射，更广泛地从几何上描述 $n_\beta$。
下面的猜想~\ref{conj:symplectic-instantons} 涉及其中一些微妙之处。
关于这里研究的 Calabi--Yau 三维簇 $V$ 的进一步细节与计算，见 \SourceCite{CFKM}。
\end{example}

我们将以一种定义 $n_\beta$ 的猜想性辛几何方法结束瞬子数的讨论。
在 $V$ 上固定一个辛结构（例如 Kähler 形式定义的结构）。
回忆\hyperref[subsec:symplectic-map-moduli]{第~\ref*{subsec:symplectic-map-moduli} 小节}，
$\mathcal M(\mathbb P^1,J,\beta)$ 是简单的 $J$ 全纯映射 $f:\mathbb P^1\to V$ 的空间，
其中 $J$ 是一般的几乎复结构，受 $V$ 的辛形式驯服。
由 \SourceCite[第 2 章]{MS}，每个 $f\in\mathcal M(\mathbb P^1,J,\beta)$
在一般点为浸入，意思是 $\mathbb P^1$ 有稠密开集 $U$，在其上 $df$ 单射，且对所有 $u\in U$，
$f^{-1}(f(u))=\{u\}$。
$\operatorname{PGL}(2,\mathbb C)$ 在 $\mathcal M(\mathbb P^1,J,\beta)$ 上自然作用，
我们将此作用的一个轨道定义为 $V$ 上类为 $\beta$ 的一条 \emph{$J$ 全纯曲线}。

由式~\eqref{eq:fixed-curve-map-index}，$\mathcal M(\mathbb P^1,J,\beta)$ 的实维数为 $6$。
这正是 $\operatorname{PGL}(2,\mathbb C)$ 的维数；由于作用自由，各轨道就是
$\mathcal M(\mathbb P^1,J,\beta)$ 的连通分支。
此外，每个分支 $\mathcal C$ 有自然定向，因此可定义符号 $\epsilon(\mathcal C)=+1$ 或 $-1$，
分别对应于映射 $\phi:\operatorname{PGL}(2,\mathbb C)\simeq\mathcal C$ 保持或不保持定向；
这里 $\phi$ 由群作用于固定的 $f_0\in\mathcal C$ 给出。于是提出如下猜想。
\begin{conjecture}\label{conj:symplectic-instantons}
设 $(V,\omega)$ 是带 Kähler 形式的 Calabi--Yau 三维簇。
固定一般的、受 $\omega$ 驯服的几乎复结构 $J$。则：
\SourcePage{232}
\begin{enumerate}[label=(\roman*)]
\item 对任意同调类 $\beta\in H_2(V,\mathbb Z)$，类为 $\beta$、亏格为零的浸入 $J$ 全纯曲线只有有限多条。
它们都是嵌入的，且任意两条都不相交（即使 $\beta$ 不同）。定义整数
\[
n'_\beta=\sum_{\mathcal C\in\pi_0\mathcal M(\mathbb P^1,J,\beta)}\epsilon(\mathcal C),
\]
其中求和遍历 $\mathcal M(\mathbb P^1,J,\beta)$ 的所有连通分支 $\mathcal C$。
\item 亏格零的所有 $J$ 全纯映射 $C\to V$，都是嵌入的多重覆盖。
\item 对每个 $\beta$，$n_\beta$ 可表示为
\[
n_\beta=n'_\beta+\sum_{\substack{\beta=k\gamma\\k>1}}
\sum_{\mathcal C\in\pi_0\mathcal M(\mathbb P^1,J,\gamma)}c(k,\mathcal C),
\]
其中 $c(k,\mathcal C)$ 是由 $\mathcal C$ 的 $k$ 次多重覆盖确定的整数。
\end{enumerate}
\end{conjecture}
此猜想说，$n_\beta$ 应为类 $\beta$ 中 $J$ 全纯曲线的带符号条数，
再加上当 $\beta$ 可被 $k$ 整除时，$k$ 次多重覆盖的某些贡献。
特别地，猜想~\ref{conj:symplectic-instantons} 推出每个 $n_\beta$ 都是整数。
$\langle I_{0,0,\beta}\rangle$ 的辛定义表明 $n_\beta$ 是辛不变量；
但如下面所解释，$n'_\beta$ 可能不是。

复杂性主要仍来自多重覆盖。
如\hyperref[sec:gw-symplectic-definition]{第~\ref*{sec:gw-symplectic-definition} 节}所述，
空间 $\mathcal M(\mathbb P^1,J,\beta)$ 不紧，而 Gromov--Uhlenbeck 紧化需要加入多重覆盖。
因此，甚至不能预先断定猜想~\ref{conj:symplectic-instantons} 中等式的右边是有限和。

读者也许会问，为何不更简单地猜想总有 $n_\beta=n'_\beta$？
事实上，我们不知道任何违反这个等式的例子。另一方面，\SourceCite{Taubes} 中有类似情形：
它分析了辛四维流形中法丛平凡的椭圆曲线的多重覆盖。
在那里，二重覆盖有一个相应的离散不变量，其值会改变该覆盖对 Gromov--Witten 不变量的贡献。
此外，所有 $J$ 构成的空间有一些称为\emph{墙}的实余维一轨迹；
当 $J$ 穿过一堵墙时，这种贡献的变化为整数。
穿墙时出现额外的类为 $2\beta$ 的 $J$ 全纯曲线，使 Gromov--Witten 不变量保持不变。
而二重覆盖的多重覆盖所产生的贡献变化，又由类 $2\beta$ 中这些新曲线的多重覆盖的贡献抵消。
我们感谢 \SourceCite{IP} 的作者解释 Taubes 的工作及他们自己的论文；
该论文研究 Taubes 定义的不变量与 Gromov--Witten 不变量的联系。
这番讨论促使我们提出猜想~\ref{conj:symplectic-instantons}。

可以设想，Calabi--Yau 三维簇上的有理曲线也出现类似情形：
所有 $J$ 的空间中可能存在墙；若 $V$ 的给定复结构位于墙上，
那么“一般”可能有两种含义，取决于“一般”的 $J$ 位于墙的哪一侧。
正如上一段所说的 $g=1$ 情形，穿墙可能使某些嵌入的 $J$ 全纯曲线变成多重覆盖。
我们希望为每个 $k>1$ 和嵌入的 $J$ 全纯曲线 $\mathcal C$ 定义离散不变量 $c(k,\mathcal C)$，
使穿墙时 $n'_\beta$ 的损失由 $c(k,\mathcal C)$ 的增加抵消，
从而猜想~\ref{conj:symplectic-instantons} 中的总和保持不变。
\SourcePage{233}
在一种情形下，猜想~\ref{conj:symplectic-instantons} 的一部分可以证明。
假设 $\beta$ 是\emph{本原}的，即不是 $H_2(V,\mathbb Z)$ 中另一个类的某个大于 $1$ 的整数倍。
\begin{proposition}\label{prop:primitive-instantons}
若 $\beta\in H_2(V,\mathbb Z)$ 非零且本原，则 $n_\beta=n'_\beta$，
其中 $n'_\beta$ 如猜想~\ref{conj:symplectic-instantons} 所定义。
\end{proposition}
\begin{proof}
这由 \SourceCite{Ruan2} 得到。采用 Li--Ruan--Tian 的 Gromov--Witten 不变量定义，
则用\hyperref[subsec:gw-equivalence]{第~\ref*{subsec:gw-equivalence} 小节}的记号，
$n_\beta=\langle I_{0,0,\beta}\rangle^{\mathrm{LRT}}$。
由式~\eqref{eq:instanton-three-divisors}，可以从 $\langle I_{0,3,\beta}\rangle^{\mathrm{LRT}}$ 恢复 $n_\beta$。
由于 $V$ 半正，可用不变量 $\Psi_{0,3,\beta}$ 代替 $\langle I_{0,3,\beta}\rangle^{\mathrm{LRT}}$。
取 $\alpha_1,\alpha_2,\alpha_3\in H^2(V,\mathbb Z)$。
用\hyperref[subsec:symplectic-gw-invariants]{第~\ref*{subsec:symplectic-gw-invariants} 小节}
定义的辛 Gromov--Witten 不变量 $\Psi,\Phi$，由式~\eqref{eq:genus-zero-mixed-symplectic}，
\[
\Psi_{0,3,\beta}([pt];\alpha_1,\alpha_2,\alpha_3)
=\Phi_{0,\beta}(\alpha_1,\alpha_2,\alpha_3\mid\,).
\]
可以用 \SourceCite{Ruan2} 的方法计算 $\Phi_{0,\beta}(\alpha_1,\alpha_2,\alpha_3\mid\,)$，
因为 $\beta$ 本原意味着所有满足 $f_*[\mathbb P^1]=\beta$ 的 $J$ \IndexLocation{indexloc:0655:233}全纯映射 $f:\mathbb P^1\to V$ 都是简单的。
所以 $\mathcal M(\mathbb P^1,J,\beta)$ 给出了类 $\beta$ 中的所有 $J$ 全纯 $\mathbb P^1$。
考虑\IndexLocation{indexloc:0276:233}赋值映射
\[
\mathcal M(\mathbb P^1,J,\beta)\times_{\operatorname{PGL}(2,\mathbb C)}(\mathbb P^1)^3
\longrightarrow V^3,\quad(f,p_1,p_2,p_3)\longmapsto(f(p_1),f(p_2),f(p_3)).
\]
此映射的像的闭包为六维圈 $\Gamma$。
由 \SourceCite[第 7.2 节]{MS}，$\Phi_{0,\beta}(\alpha_1,\alpha_2,\alpha_3\mid\,)$
为交数 $\Gamma\cdot(\Gamma_1\times\Gamma_2\times\Gamma_3)$，
其中与定理~\ref{thm:multiple-cover} 的证明一样，$\Gamma_i$ 是与 $\alpha_i$ 对偶的四维圈。

已知 $\mathcal M(\mathbb P^1,J,\beta)$ 为各连通分支 $\mathcal C$ 的并。
固定 $f_0\in\mathcal C$，让 $\operatorname{PGL}(2,\mathbb C)$ 作用于 $f_0$，得到保持定向的同构
$\mathcal C\simeq\epsilon(\mathcal C)\operatorname{PGL}(2,\mathbb C)$，
其中 $\epsilon(\mathcal C)$ 决定如何给群定向。
据此，上述映射变成
\[
f_0\times f_0\times f_0:\epsilon(\mathcal C)(\mathbb P^1)^3
\simeq\mathcal C\times_{\operatorname{PGL}(2,\mathbb C)}(\mathbb P^1)^3\longrightarrow V^3.
\]
由于 $f_{0*}[\mathbb P^1]=\beta$，其像与 $\Gamma_1\times\Gamma_2\times\Gamma_3$ 的交数为
$\epsilon(\mathcal C)\int_\beta\alpha_1\int_\beta\alpha_2\int_\beta\alpha_3$。
对所有分支求和，就得到所需结果。
\end{proof}
类似地，对于五次三维簇，猜想~\ref{conj:symplectic-instantons} 对所有满足 $d\leq9$、$d\ne5$ 的次数 $d$ 成立。
原因是：取可积的 $J$，就已能达到足够一般的情形，
然后用例题~\ref{ex:quintic-instanton-enumerative} 的论证得到 $n_d=n'_d$，
这里的 $n'_d$ 如猜想~\ref{conj:symplectic-instantons} 所定义。

\subsection{亏格一不变量}\label{subsec:genus-one-invariants}
\index{Gromov--Witten不变量@Gromov--Witten 不变量!亏格一|hyperpage}\index{椭圆曲线!计数|hyperpage}\index{固定j不变量@固定 $j$ 不变量|hyperpage}\index{退化瞬子|hyperpage}
\index{Gromov--Witten不变量@Gromov--Witten 不变量!亏格一|hyperpage}
当 $g=1$ 时，得到\emph{椭圆 Gromov--Witten 不变量}
$\langle I_{1,n,\beta}\rangle(\alpha_1,\ldots,\alpha_n)$。
这些不变量最初在 \SourceCite{BCOV1} 中讨论，此后已在各种背景下得到研究。
下面的例子说明，确定 Gromov--Witten 不变量的枚举意义，往往是一项微妙的工作。
\SourceExampleNumber{7.4.5}{1}
\begin{example}\label{ex:plane-genus-one}
先研究射影平面 $\mathbb P^2$ 的椭圆 Gromov--Witten 不变量。
若 $\ell\subset\mathbb P^2$ 为直线，与往常一样，用 $\langle I_{1,n,d}\rangle$ 代替 $\langle I_{1,n,d[\ell]}\rangle$。
\SourcePage{234}
首先考虑 $\langle I_{1,1,0}\rangle$，它出人意料地并非平凡。
这里处理恒定映射（次数为零），但\hyperref[sec:gw-properties]{第~\ref*{sec:gw-properties} 节}的点映射公理不适用。
不过，如 \SourceCite{KoM1} 所解释，对每个亏格都有相应版本。
对 $g=1$ 和任意 $X$，当 $\deg\alpha=2$ 时，得到
\begin{equation}\label{eq:genus-one-constant}
\langle I_{1,1,0}\rangle(\alpha)=-\frac1{24}\int_Xc_{\dim X-1}(T_X)\cup\alpha.
\end{equation}
问题在于，$n$ 点、亏格一、次数零的映射 $f:C\to X$ 的空间是
$\overline M_{1,n}\times X$，它并不具有预期维数。
所以必须像\hyperref[subsec:computing-virtual-class]{第~\ref*{subsec:computing-virtual-class} 小节}一样计算虚拟基本类。
可用命题~\ref{prop:excess-obstruction-euler-class} 完成这一计算，例题~\ref{ex:genus-one-gravitational-euler} 将详细说明，
式~\eqref{eq:genus-one-constant} 随即得到。
将它用于 $X=\mathbb P^2$，得到
\[
\langle I_{1,1,0}\rangle([\ell])=-\frac1{24}\int_{\mathbb P^2}c_1(T_{\mathbb P^2})\cup[\ell]=-\frac18.
\]
所以，Gromov--Witten 不变量既可能为分数，也可能为负！
例题~\ref{ex:genus-one-gravitational-euler} 在将 Gromov--Witten 不变量推广为\emph{引力关联数}之后，将讨论一个有趣的相关计算。

幸运的是，其他椭圆不变量具有很好的枚举意义。对 $d\geq1$，定义
\[
N_d^1=\langle I_{1,3d,d}\rangle(\underbrace{[pt],\ldots,[pt]}_{3d\text{ 次}}).
\]
可以证明，$N_d^1$ 是 $\mathbb P^2$ 中经过 $3d$ 个一般点的 $d$ 次椭圆曲线的数目。
于是 $N_1^1=N_2^1=0$、$N_3^1=1$，因为平面中经过九个一般点的光滑椭圆曲线唯一。
当 $d>3$ 时，$N_d^1$ 所计数的曲线都奇异。
容易证明：一旦知道 $\langle I_{1,1,0}\rangle([\ell])$ 与所有 $d\geq1$ 的 $N_d^1$，
就确定了全部椭圆 Gromov--Witten 不变量。

关于 $N_d^1$ 的主要结果是各种递推关系；一个例子见 \SourceCite{Getzler1}。
\SourceCite{EHX} 用 Virasoro 约束的方法预测了一个漂亮的递推式，\hyperref[subsec:virasoro-conjecture]{第~\ref*{subsec:virasoro-conjecture} 小节}将简要讨论这一方法。
此公式由 \SourceCite{Pandharipande2} 证明：若 $N_d$ 是式~\eqref{eq:plane-nd} 给出的有理曲线数目，则
\begin{equation}\label{eq:plane-genus-one-recursion}
N_d^1=\frac1{12}\binom d3N_d+
\sum_{\substack{d=d_1+d_2\\d_1,d_2>0}}
\frac{3d_1^2d_2-2d_1d_2}{9}N_{d_1}N_{d_2}^1\binom{3d-1}{3d_1-1}.
\end{equation}
\SourceCite{Pandharipande2} 的证明，使用 Getzler 在
$H^4(\overline M_{1,4},\mathbb Q)$ 中发现的余维二分层之间一个关系的几何构造，
然后应用\hyperref[chap:quantum-cohomology]{第~\ref*{chap:quantum-cohomology} 章}将讨论的 WDVV 方程。
由此可按需要计算任意多个 $N_d^1$。一个初等证明见 \SourceCite{Ran3}。

在迄今讨论的椭圆不变量中，我们允许 $j$ 不变量变化。
但也可研究 $\mathbb P^2$ 中 $j$ 不变量\emph{固定}的椭圆曲线。
更准确地，令 $E_{d,j}$ 为 $\mathbb P^2$ 中经过 $3d-1$ 个一般点、具有给定 $j$ 不变量的
$d$ 次不可约椭圆曲线数目。
\SourcePage{235}
\SourceCite{Ionel,Pandharipande1} 证明
\begin{equation}\label{eq:plane-fixed-j}
E_{d,j}=\begin{cases}
\binom{d-1}{2}N_d,&j\ne0,1728,\\
\frac13\binom{d-1}{2}N_d,&j=0,\\
\frac12\binom{d-1}{2}N_d,&j=1728,
\end{cases}
\end{equation}
其中 $N_d$ 如上。

试用 Gromov--Witten 类进行这一计算。为简便，假设 $j\ne0,1728$。
由于希望经过 $\mathbb P^2$ 中的 $3d-1$ 个点，从 Gromov--Witten 类
\[
\xi=I_{1,3d-1,d}(\underbrace{[pt],\ldots,[pt]}_{3d-1\text{ 次}})
\]
出发。由式~\eqref{eq:gw-top-degree-condition}，这个类的次数为 $6d-4$；
因 $\overline M_{1,3d-1}$ 维数为 $3d-1$，它比最高次数低二。
固定 $j$，得到由所有具有此 $j$ 不变量的带标记曲线构成的除子
$Z_j\subset\overline M_{1,3d-1}$。
朴素地，我们可能预期 $E_{d,j}$ 是 $\xi$ 与 $Z_j$ 的交，即
\begin{equation}\label{eq:plane-fixed-j-naive}
E_{d,j}=\xi\cdot Z_j.
\end{equation}
如下计算 $\xi\cdot Z_j$。
自然映射 $\overline M_{1,3d-1}\to\overline M_{1,1}\simeq\mathbb P^1$
遗忘最后的 $3d-2$ 个点，并收缩不稳定分支。由于 $j\ne0,1728$，$Z_j$ 就是 $j$ 上的纤维。
而 $\infty$ 上的纤维由所有映到结点曲线的 $\overline M_{1,3d-1}$ 中的元素组成。
容易看出，这个纤维是映射
$\psi:\overline M_{0,3d+1}\to\overline M_{1,3d-1}$ 的像，$\psi$ 将最后两个点粘合。
所以 $Z_j\sim\operatorname{Im}(\psi)$。
这意味着
\[
\begin{aligned}
\xi\cdot Z_j=\xi\cdot\operatorname{Im}(\psi)=\frac12\Big(&
\langle I_{0,3d+1,d}\rangle([pt],\ldots,[pt],[pt],[\mathbb P^2])\\
&+\langle I_{0,3d+1,d}\rangle([pt],\ldots,[pt],[\ell],[\ell])\\
&+\langle I_{0,3d+1,d}\rangle([pt],\ldots,[pt],[\mathbb P^2],[pt])\Big).
\end{aligned}
\]
这是第一次使用约化公理。由此，基本类公理与除子公理容易给出
\[
\xi\cdot Z_j=\xi\cdot\operatorname{Im}(\psi)=\frac12d^2N_d.
\]
与式~\eqref{eq:plane-fixed-j} 比较，可见式~\eqref{eq:plane-fixed-j-naive} 是错误的。
原因在于，$E_{d,j}$ 只计数 $\mathbb P^2$ 中的不可约曲线，
而 $\xi\cdot Z_j$ 还包括来自 $\overline M_{1,3d-1}$ 边界的可约曲线。
事实上，\SourceCite{Ionel} 的论证先进行上述计算，再在辛几何背景下仔细计入边界，来求 $E_{d,j}$。
还应提到，用\hyperref[sec:gw-symplectic-definition]{第~\ref*{sec:gw-symplectic-definition} 节}
定义的混合不变量与辛不变量，有
\[
\xi\cdot Z_j=\frac12\Phi_{1,d}([pt]\mid\underbrace{[pt],\ldots,[pt]}_{3d-2\text{ 次}})
=\Psi_{1,3d-1,d}(Z_j;\underbrace{[pt],\ldots,[pt]}_{3d-1\text{ 次}}),
\]
尤其参见\hyperref[subsec:symplectic-gw-invariants]{第~\ref*{subsec:symplectic-gw-invariants} 小节}中混合不变量与辛不变量的关系。
这些不变量及相关椭圆不变量的\IndexLocation{indexloc:0534:235}枚举意义，见 \SourceCite[注记 2.19]{RT2}。
\end{example}
\SourcePage{236}
\SourceExampleNumber{7.4.5}{2}
\begin{example}\label{ex:space-genus-one}
在 $\mathbb P^3$ 中取直线 $\ell$，并定义
\[
\begin{aligned}
N_{a,b}&=\langle I_{0,a+b,(a+2b)/4}\rangle
(\underbrace{[\ell],\ldots,[\ell]}_{a\text{ 次}},\underbrace{[pt],\ldots,[pt]}_{b\text{ 次}}),\\
N_{a,b}^1&=\langle I_{1,a+b,(a+2b)/4}\rangle
(\underbrace{[\ell],\ldots,[\ell]}_{a\text{ 次}},\underbrace{[pt],\ldots,[pt]}_{b\text{ 次}}),
\end{aligned}
\]
其中假设 $(a+2b)/4$ 为整数。
由定理~\ref{thm:homogeneous-gw-equivalence}，$N_{a,b}$ 是 $\mathbb P^3$ 中
经过 $a$ 条一般直线及 $b$ 个一般点的 $d=(a+2b)/4$ 次有理曲线数目。
$N_{a,b}$ 的递推关系见 \SourceCite{FP}。

椭圆不变量 $N_{a,b}^1$ 也可用 \SourceCite{Getzler1} 附录中的递推关系计算。
但与亏格零情形不同，它们不是枚举不变量。
例如，\SourceCite{Getzler1} 附录算出的\emph{所有} $N_{a,b}^1$ 都是负数，许多还是分数。
不过，若取线性组合
\[
E_{a,b}=N_{a,b}^1+\frac{2d-1}{12}N_{a,b},
\]
其中仍有 $d=(a+2b)/4$，则 \SourceCite{GeP} 表明，$E_{a,b}$ 是
$\mathbb P^3$ 中经过 $a$ 条一般直线及 $b$ 个一般点的 $d$ 次椭圆曲线数目。
计算 $N_{a,b}^1$ 的其他方法，以及对 $\mathbb P^n$ 的推广，见 \SourceCite{Vakil3,GPa}。

还应注意，也能研究 $\mathbb P^3$ 中 $j$ 不变量固定的椭圆曲线。
与 $\mathbb P^2$ 一样，可得到这类曲线的数目与 $N_{a,b}$ 之间的关系式。
精确公式见 \SourceCite{Ionel}。
\end{example}
\SourceExampleNumber{7.4.5}{3}
\begin{example}\label{ex:quintic-genus-one}
最后，回到一般的五次三维簇 $V\subset\mathbb P^4$。
从式~\eqref{eq:cy-a-correlation-instanton} 的椭圆版本开始。
令 $n_d$ 为 $V$ 上 $d$ 次有理曲线的瞬子数，令 $e_d$ 为 $V$ 上 $d$ 次“椭圆曲线”的类似瞬子数
（下面解释引号的含义）。假设这些数有限，并将 $n_d,e_d$ 理解为计数适当的刚性嵌入曲线。
那么有公式
\begin{equation}\label{eq:quintic-genus-one-instanton}
\langle I_{1,1,d}\rangle(H)=\sum_{k\mid d}\sigma\!\left(\frac d k\right)ke_k
+\frac1{12}\sum_{k\mid d}kn_k,
\end{equation}
其中 $\sigma(d/k)=\sum_{l\mid(d/k)}l$ 是 $d/k$ 的所有正因子之和。

此公式通过物理推理得到，见 \SourceCite{BCOV1}，但迄今尚无数学证明。
不过，可作如下直观解释。先看式~\eqref{eq:quintic-genus-one-instanton} 中第一个和。
设 $f:(C,p)\to V$ 为稳定映射，其像 $E=f(C)$ 为椭圆曲线，
$C\to E$ 的次数为 $d/k$，$E\subset V$ 的次数为 $k$。
那么 $E$ 是这 $e_k$ 条椭圆曲线之一；由于 $f(p)\in E\cap H$，$f(p)$ 有 $k$ 种选择。
给定 $f(p)$，$p\in C$ 有 $d/k$ 种选择，但因 $C\to E$ 是覆盖，它们彼此同构。
因此，只需计数 $E$ 的 $d/k$ 次覆盖空间。这相当于计数 $\mathbb Z^2$ 中指数为 $d/k$ 的子格，
众所周知，其数目为 $\sigma(d/k)$。

式~\eqref{eq:quintic-genus-one-instanton} 中的第二个和更微妙。例如，考虑 $k=d$。
那 $n_d$ 条 $d$ 次“有理曲线”对 $\langle I_{1,1,d}\rangle(H)$ 如何作出贡献？
取 $d$ 次有理曲线 $C_0\subset V$。如往常一样，$f(p)\in C_0\cap H$ 有 $d$ 种可能。
给定 $C_0$ 与 $f(p)$，我们想构造次数为 $1$ 的亏格一稳定映射 $f:(C,p)\to C_0$。
这似乎不可能，
\SourcePage{237}
但稳定映射的定义允许 $C$ 可约。
物理文献把这种情形称为\emph{退化瞬子}。
这里 $C=E\cup C_0$，$E$ 是在一点与 $C_0$ 相接的椭圆曲线，
$f:C\to C_0$ 在 $C_0$ 上为恒等映射，在 $E$ 上为恒定映射。
一旦选取映到 $f(p)$ 的点 $p\in C$，就得到我们所关心的这类稳定映射的一个族。

若第二个和中的 $k\mid d$ 是 $d$ 的真因子，情形更复杂，主要原因是多重覆盖。
与亏格零情形一样，需要处理有理曲线的多重覆盖；
此外，还有由椭圆曲线到 $\mathbb P^1$ 的熟知二次映射产生的多重覆盖。
物理计算 \SourceCite{BCOV1} 提示，答案仍为 $1/12$；
当底层有理曲线的法丛为 $\mathcal O_{\mathbb P^1}(-1)\oplus\mathcal O_{\mathbb P^1}(-1)$ 时，
这一点已得到数学证明，见 \SourceCite{GPa}。
所以，若 Clemens 猜想~\ref{conj:clemens} 的假设对所有整除 $d$ 的次数成立，
并且椭圆曲线也满足类似假设，彼此不相交且与有理曲线不相交，
则式~\eqref{eq:quintic-genus-one-instanton} 对该次数 $d$ 成立。

还有一个复杂之处：即使 $e_d$ 很可能有限，它也可能计数 $V$ 上光滑椭圆曲线之外的对象。
这就是我们为“椭圆曲线”加引号的原因。
最好的例子是五次三维簇上四次“椭圆曲线”的数目 $e_4$。
\SourceCite{BCOV1} 将此数算为
\[
e_4=3721431625.
\]
然而，由 \SourceCite{ES2}，一般五次三维簇上的光滑四次椭圆曲线数目是
\[
3718024750.
\]
它与 $e_4$ 不同。如 \SourceCite{ES2} 所解释，差别来自 $e_4$ 允许\emph{奇异}曲线。
具体地，$V$ 含有若干双结点的四次平面曲线，每条都产生一个对 $e_4$ 有贡献的稳定映射。
这些奇异曲线如下计数。若 $\ell$ 是 $V$ 上的 $2875$ 条直线之一，
则包含 $\ell$ 的任意平面 $P\subset\mathbb P^4$，与 $V$ 的交为 $\ell\cup C$，其中 $C$ 是平面四次曲线。
一般情况下，$C$ 光滑；但在有限多个特殊平面 $P$ 上，它出现两个结点。
\SourceCite{Vainsencher} 证明，这样的特殊平面 $P$ 共有 $1185$ 个。
由于
\[
3721431625=3718024750+1185\times2875,
\]
这些数是相符的。这又说明，在理解 Gromov--Witten 不变量的\IndexLocation{indexloc:0534:237}枚举意义时，必须十分仔细。
\end{example}

\begingroup\linespread{1.22}\selectfont
\subsection{其他例题}\label{subsec:other-gw-examples}
除已经列出的例子外，关于 Gromov--Witten 不变量与枚举几何，还有\emph{大量}其他工作。
例子的清单已经很长，下面列出一些供进一步阅读的文献。

\looseness=-1 镜像对称的枚举方面可追溯到 1991 年：\SourceCite{CdGP} 对五次三维簇上的有理曲线作出了惊人的预测。
随后，大批论文对 Calabi--Yau 流形作出枚举预测，
包括物理论文 \SourceCite{BCOV1,BCOV2,CdFKM,CFKM,HKTY1,HKTY2,BKK,AM2,Font,KT1,KT2}，
以及数学论文 \SourceCite{Katz2,Morrison1,ES2,ES1,LTe,BvS,Voisin1,Meurer,HSS,Givental2,LLY,Givental4,BCFKvS2}。
关于镜像对称枚举方面的一篇很好的综述，见 \SourceCite{Morrison6}。
\SourcePage{238}
虽然 Calabi--Yau 流形的枚举几何仍有许多未解决问题，但不少较新的工作已经转向其他类型的代数簇。
这部分源于\hyperref[sec:gw-algebraic-definition]{第~\ref*{sec:gw-algebraic-definition} 节}
所述任意光滑射影簇的 Gromov--Witten 不变量定义，它促成了上述许多例子与论文。
此外，还有对有理曲面 \SourceCite{CM,GoP,Vakil1}、Del Pezzo 曲面 \SourceCite{Vakil2}、
$\mathbb P^n$ 的爆破 \SourceCite{Gathmann1}、凸簇的爆破 \SourceCite{Gathmann2}、
完全交 \SourceCite{Beauville1}，以及 $\mathbb P^2$ 中带结点的亏格二曲线 \SourceCite{KQR} 的应用。
与此同时，人们也受到启发，探索不用 Gromov--Witten 理论能做些什么。
由此产生了 Crauder 与 Miranda 的论文 \SourceCite{CM}，Ran 的论文 \SourceCite{Ran1,Ran3}，
以及 Caporaso 与 Harris 的论文 \SourceCite{CH1,CH2,CH3}（综述见 \SourceCite{Caporaso}）。
五次三维簇及其他相关 Calabi--Yau 三维簇上的 Clemens 猜想也得到研究，见 \SourceCite{JK1,JK2,EJS}。
综观这些论文，近年来对枚举几何的兴趣迅速增长。
这种兴趣可能起源于镜像对称，但如今在代数几何中已有自身的发展动力。

\enlargethispage{\baselineskip}
\looseness=-1 这一长章至此结束。我们已经看到，Gromov--Witten 不变量与辛几何建立了非平凡的联系，
并在当代代数几何中发挥重要作用。本章的一个启示是，曲线计数比预想的更加深刻、丰富。
然而，故事尚未讲完：我们还需要了解量子上同调。
\hyperref[chap:quantum-cohomology]{第~\ref*{chap:quantum-cohomology} 章}将看到，它会进一步解释这里的例子，
并补全理解镜像对称的数学表述所需的一整套思想。
\par\endgroup
